\documentclass[11pt,twocolumn]{article}

%\usepackage{sectsty}
\usepackage{titlesec}
\usepackage{url}

\usepackage{amssymb}

\usepackage{titlesec}

\usepackage{transparent}

\usepackage{setspace}
%\doublespacing
%\onehalfspacing

\usepackage{graphicx}
%\graphicspath{ {/ext_root/openshift-oc3/mmui-caf-local/web/mmui/public/pics/} }

\newcommand{\sectionGraphics}{

\vspace{2em}\centerline{\includegraphics[scale=0.125]{logo-deco.png}}\vspace{-1em}}

\usepackage[flushmargin]{footmisc}

\setlength{\parindent}{30pt}

%\usepackage{bigfoot}

%\DeclareNewFootnote[para]{R}[Roman]
%\DeclareNewFootnote[para]{N}

%\usepackage[letterpaper, left=1.3in,right=1in,top=1in,bottom=1in]{geometry}

\newcommand{\descItem}[1]{\item[{\color{logoRed!80!logoBlue}{\fontfamily{phv}%
\fontshape{it}\fontseries{b}\selectfont\scriptsize #1}}]}
%\newcommand{\descItem}[1]{\item[#1]}

\newcommand{\descItemLarge}[1]{\item[{\color{logoRed!80!logoBlue}{\fontfamily{phv}%
\fontshape{it}\fontseries{b}\selectfont\small #1}}]}


\usepackage{enumitem}

\setlist[description]{%
  topsep=30pt,               % space before start / after end of list
  itemsep=5pt,               % space between items
  font={\bfseries\sffamily}, % set the label font
%  font={\bfseries\sffamily\color{red}}, % if colour is needed
}


\usepackage{etoolbox}

\usepackage{mathptmx}
%\usepackage{euler}
%\usepackage{amsmath}

%\usepackage[LGRgreek]{mathastext}

%\sectionfont{\fontsize{12}{15}\selectfont}
%\subsectionfont{\fontsize{11}{12}\selectfont}

\titleformat*{\subsection}{\small\bfseries}

\usepackage{eufrak}
\usepackage{wasysym}
\usepackage{textcomp}
\usepackage{amssymb}

\usepackage{microtype}


\usepackage{graphicx}

\usepackage{tikz}

\usetikzlibrary{decorations.pathmorphing}
\usetikzlibrary{positioning}
\usetikzlibrary{shapes,snakes}


\let\OldSection\section

\renewcommand\section[1]{
	\vspace{9pt}
	
	\scalebox{1.3}{\colorbox{logoPurple!50}{\hspace{1em}}}
				
	%\hspace{-1.1cm}{\protect\includegraphics[width=1cm]{e-logo.png}}	
	{\protect\transparent{0.3}{\colorbox{logoPeach!10!blue}{%
			\begin{minipage}{\linewidth}
				    \vspace{.5em}
				
	\protect\transparent{1}{\OldSection{#1%
		%{\transparent{0.5}{
		%	\protect\includegraphics[width=1cm]{e-logo.png}
		%}
	    %}
     }}
	    \end{minipage}}} }
    
    \vspace{-5em}
    	{\protect\transparent{0.3}{\colorbox{logoPeach}{%
    		\begin{minipage}{\linewidth}
    			\hspace{\linewidth}
    	\end{minipage}}} } 
    
    \vspace{5em}
    
	\vspace{-6pt}
}

%\subsectionfont{\Large}
%\subsectionfont{\fontsize{12}{15}\selectfont}

\titleformat*{\subsection}{\Large\bfseries}

\let\OldSubsection\subsection
\renewcommand\subsection[1]{

	\vspace{12pt}
	
    %{\LARGE
	%\colorbox{logoCyan}{%
	%\begin{minipage}{\linewidth}
		\OldSubsection{% 	
			\hspace{-2.75em}
			%\begin{minipage}{}
			\protect\raisebox{-5pt}{%
			\colorbox{logoCyan!50}{\hspace{2.1em}}}%
			\hspace{-5pt}{\protect\transparent{0.3}{\colorbox{logoBlue!80}{\protect\transparent{1}{%
						   \protect\raisebox{1pt}{\textit{{\large #1}}} }}}}
			%\end{minipage}
		}
	%\end{minipage}
    %}
    %}
	\vspace{-6pt}
}



\newcommand*{\MyCircle}{\begin{tikzpicture}{baseline}
	\node [draw, shape=regular polygon, regular polygon sides = 5, shape border rotate=180,
	line width=.2mm, left color=blGreen, right color=darkBlGreen, draw=darkRed, aspect=1,
	draw opacity=0.6, fill opacity=0.8, xscale=0.5, yscale=0.5,
	minimum width=1mm, 
	] 
	at (0, 0) {};
	\end{tikzpicture}}


\newcommand*{\MyLabel}{%
 %\MyBall	
 %{\tikz \draw [baseline, ball color=red, draw=red] circle (3pt);} 	
 \protect\MyCircle{\hspace{4pt}}\color{flColor}\LARGE{\textbf{\protect\raisebox{-3pt}{\theenumi}}}
}

%
\newenvironment{colorDescription}{%
\vspace{1em}\begin{description}[style=nextline]}{\end{description}}

\newcommand*{\MyDiamond}{%
	\begin{tikzpicture}{baseline}
	
	\node [draw, shape=diamond, %shape border rotate=100,
	line width=.2mm, left color=flColor, right color=darkRed, draw=logoPurple, aspect=1,
	draw opacity=0.6, fill opacity=0.8, xscale=0.5, yscale=0.5,
	minimum width=1mm, 
	] 
	at (2, 0) {};
	
	\end{tikzpicture}
	
} 

\usepackage{enumitem}

\newcommand*{\MyDiamon}{*} 


\newenvironment{colorItemize}{\vspace{1em}\begin{itemize}[label={\MyDiamond}]}{\end{itemize}}

\newenvironment{colorEnumerate}{\vspace{1em}\begin{enumerate}[label=\MyLabel]}{\end{enumerate}}

\newcommand{\PeoplesVoiceCafe}{{\relscale{0.8}{\fontfamily{fvs}\fontseries{b}\selectfont%
 Peoples' Voice Cafe}}}
 
\newcommand{\q}[1]{{\fontfamily{qcr}\selectfont ``}#1{\fontfamily{qcr}\selectfont ''}} 

\usepackage{tcolorbox}

%\usepackage[framemethod=TikZ]{mdframed}

%
\tcbuselibrary{skins}
%\tcbuselibrary{listings,breakable}
\usetikzlibrary{calc}
\usetikzlibrary{shadows}
\pgfdeclarelayer{background}
\pgfdeclarelayer{foreground}
\pgfsetlayers{background,main,foreground}

\definecolor{BaseColor}{HTML}{8533FF}

\newlength{\boxShadowOfset}
\setlength{\boxShadowOfset}{2mm}

\newcommand{\tc}[1]{
\vspace*{-2.05cm}\hspace{1cm}
\begin{tcolorbox}
[colframe=cyan!10!black,boxrule=0.5pt,arc=8pt,hbox,enhanced, 
leftrule=5pt,rightrule=3pt,toprule=0pt,toprule=0pt,
%drop fuzzy shadow southeast={BaseColor!50!white},
%frame code={
%        },
     % left=6pt,right=6pt,top=6pt,bottom=6pt,
      boxsep=0pt]      
\protect{#1}
\end{tcolorbox}      
\vspace*{0.4cm}

}

\newcommand{\cursive}[2]{{\LARGE%
%\begin{minipage}{0.5\textwidth}
\protect\tc{%
%\protect\fcolorbox{cyan!10!black}{white}{
\protect\raisebox{-2pt}{
{\color{logoBlue!90!cyan}{%
{\fontfamily{bch}\fontseries{b}\selectfont%
 \larger{#1}}}}}{{\color{logoBlue!90!cyan}{\Fontlukas\bfseries\textsc{#2}}}}%
%\end{minipage}
%}
}}}
%\end{minipage}

\newcommand{\parlead}[1]{\vspace{1.15em}\noindent{%
{\color{logoRed!80!logoBlue}{\fontfamily{bch}\fontseries{b}\selectfont\large #1}}}}

\newcommand{\itemLarger}[1]{\item {\raisebox{-2pt}{\larger{#1}}}}

\newcommand{\itemLargerLine}[2]{\item {\raisebox{-2pt}{\larger{#1}}}%
\hspace{1em}{\raisebox{-2pt}{#2}}}

\newcommand{\p}{

\vspace{.45em}}

\newcommand{\txtbl}{\hspace{2pt}{{\raisebox{1pt}{\relsize{-2}\textbullet}}}\hspace{2pt}}

\newcommand{\colorhref}[1]{\colorbox{blGreen!20}{#1}}



\PassOptionsToPackage{svgnames}{xcolor}

\usepackage[object=vectorian]{pgfornament} %%  http://altermundus.com/pages/tkz/ornament/index.html
\usepackage{lipsum,tikz}

\newcommand{\sectionline}[1]{%
  \noindent
  \begin{center}
  {\color{#1}
    \resizebox{0.5\linewidth}{1ex}
    {{%
    {\begin{tikzpicture}
    \node  (C) at (0,0) {};
    \node (D) at (9,0) {};
    \path (C) to [ornament=84] (D);
    \end{tikzpicture}}}}}%
    \end{center}
  }
%% A macro with two arguments to change ornaments and colors easily
%% Syntax -- \sectionlinetwo{<color>}{<ornament>}
\newcommand{\sectionlinetwo}[2]{%
  \nointerlineskip \vspace{.5\baselineskip}\hspace{\fill}
  {\color{#1}
    \resizebox{0.5\linewidth}{2ex}
    {{%
    {\begin{tikzpicture}
    \node  (C) at (0,0) {};
    \node (D) at (9,0) {};
    \path (C) to [ornament=#2] (D);
    \end{tikzpicture}}}}}%
    \hspace{\fill}
    \par\nointerlineskip \vspace{.5\baselineskip}
  }

\newcommand{\EnglischeLinie}{
\sectionline{darkRed!60!cyan}
}

\newcommand{\lMOSAIC}{\mbox{MOSAIC}}
\newcommand{\lfMOSAIC}{\mbox{M\small{OSAIC}}}
\newcommand{\MOSAIC}{\mbox{\small{MOSAIC}}}
\newcommand{\RAK}{\mbox{\small{RAK}}}
\newcommand{\SDK}{\mbox{\small{SDK}}}
\newcommand{\lSDK}{\mbox{SDK}}
\newcommand{\GUI}{\mbox{\small{GUI}}}
\newcommand{\qvspace}{\vspace{.4em}}
\begin{document}\title{Cognitive State Semantics and the 
Interface Theory of Meaning}
\maketitle{}
\begin{abstract}\noindent Start with the question of conditions of possibility for 
language to exist: how spoken, written, or inscribed signs 
differ from other kinds of things, enough to establish language as 
something that exists in the world.  Language is a regime of 
multiple conversation partners, an ambient surrounding where they 
share actions, perceptions, and foci of attention, and a common posit of structured rules.  
If communication is successful, each sign's referent is 
isolated, collectively/phenomenologically, from a (canonically spatiotemporal) reality 
that extends beyond it.  My goal here is to sketch this process 
with the help of formal and Computational linguistic theories, 
but keeping in sight a cognitive nuance and intersubjectivity that is 
(thankfully) an essential, unavoidable part of human language.
I will incorporate technical models mostly from Link Grammar 
and Type-Theoretic Semantics, understanding the latter as possibly 
layering on the former to complete a syntax/semantic pairing.  
I propose both cases for and limits on formal theories' 
applicability for Cognitive Linguistics/Humanities/Phenomenology, through the lens of the 
Philosophy of Science \mdash{} 
arguing how Phenomenology and Cognitive Linguistics, while distinct traditions, can work powerfully in consort.
\end{abstract}

 \saying{
 	On conna\^{\OldI}t la c\'el\`ebre affirmation de Claude L\'evi-Strauss: 
 	\q{les sciences humaines seront structurales ou ne seront pas}.  Nous aimerions lui en
 	adjoindre une autre: \q{les sciences humaines seront des sciences naturelles ou ne seront pas}. 
 	Evidemment, sauf \`a en revenir \`a un r\'eductionnisme dogmatique, une telle
 	affirmation n'est soutenable que si l'on peut suffisamment g\'en\'eraliser le concept
 	classique de \q{naturalit\'e}, le g\'en\'eraliser jusqu'\`a pouvoir y faire droit, 
 	comme \`a des ph\'enom\`enes naturels, aux ph\'enom\`enes d'organisation structurale.
 }
 \sayingsrc{Jean Petitot, \textit{Syntaxe Topologique et Grammaire Cognitive}}
 
{\vspace{-.25em}}
\saying{The nature of any entity, I propose, divides into three aspects or facets, which we may call its
	form, appearance, and substrate.  In an act of consciousness, accordingly, we must distinguish
	three fundamentally different aspects: its form or intentional structure, its appearance or
	subjective \q{feel}, and its substrate or origin.  In terms of this three-facet distinction, 
	we can define the place of consciousness in the world.}
\sayingsrc{David Woodruff Smith, \i{Mind World}}
 
{\vspace{-.25em}}

\vspace{1em}
\p{Computational and scientific data representation intrinsically 
balances two competing priorities: \q{semantic} expressiveness and 
computational tractability.  On the one hand, representations should not 
obscure important details: the formal requirements on 
representational validity should not force representations 
into structures that necessitate the elimination of meaningful 
information.  On the other hand, conversely, representations 
should have enough structural consistency that they are 
amenable to analysis and transformations driven by 
formal algorithms.
}
\p{I have just left a lot of loose ends: in particular, these comments need 
to give some meaningful definition to \q{representation}.  There 
are many media wherein \q{data} can be represented: via graphics 
(e.g., charts, \q{maps} in the cartographical sense, or \q{graphs} 
in the sense of plotted visualizations of 
mathematical functions), via printed documents (like the logarithmic 
tables or astronomic records of early modern science), or via mathematical 
equations and formulae (if a mathematical theory correctly 
predicts and quantifies empirical data \mdash{} e.g., fitting trajectories 
to elliptic or parabolic shapes \mdash{} then the numeric structures of 
the theory are a proxy representation of the corresponding data).  
Moreover, aside from these relatively formal modes of representation, 
we have the capacity to indirectly describe information via 
natural language. 
}
\p{The modern age has a further notion of representation: \i{digital} 
representations which leverage the capabilities of 
computer storage and processing to create persistent 
data repositories, with the possibility of manipulating 
data via computer programming and displaying data 
via computer software.  These digital representations incorporate 
features of the older representational media I alluded to: 
like Natural Language digital data often emanates from a 
textual representation, with data structures initialized 
from special \q{languages} like \XML{} or \JSON{}.  As with 
structured \q{printed} documents (of the astronomical-table 
or even grocery-shopping-list variety), digital representations 
often build off of a rigid structural template, like a two-dimensional 
table with rows and columns, or a hierarchical document 
whose elements can be nested documents.  And as with 
mathematical representations, digital representations can be 
analyzed as instances of spaces or structures with certain 
algebraic or syntactic properties and requirements.  
It is often possible to mathematically describe the 
full space of structures which conform to a given 
representational protocol, or the full space of transformations 
that can modify a given data structure while staying 
consistent with its enforced protocols.  Also, each data 
structure potentially has some notion of aggregateness: of 
having different \q{parts}, of being able to focus on one part 
at a time, and to \q{move} focus from one part to another.  
Given these qualities, digital representations extend and integrate  
the affordances of print, graphical, and linguistic media 
\mdash{} but they do so in an electronic environment that permits digital 
archives to be accessed via computer software.  With sotfware as their  
access point, digital representations can inherit the formal properties of 
the software that uses them, so we can think of digital representations as 
formal structures themselves, amenable to something like a mathematical analysis. 
}
\p{In short, digital representations have several general criteria: 
\i{validity} (a formal model of conformant vs. invalid 
structures\footnote{An invalid graph might be a case where an 
edge has no incident nodes; an invalid \XML{} structure is one without 
a single root element or (insofar as such a structure would be 
representable) with mismatched tags, and so forth)}); 
\i{traversability} (a notion of parthood and iterating
or refocusing among parts), and, let's say, \i{atomicity} 
(a notion of unitary parts that can be represented as 
integral wholes).  These criteria ensure that digitally 
represented data structures can work within software and 
networking representations: information is 
presented to software users by displaying atomic 
units (textually or graphically) and traversing through 
data structures to fill in, via application code, a 
visual tableau presenting compound data, with different 
unitary displays visually separated and often organized into 
coherent visual patterns, like the grid-pattern of a 
spreadsheet.  Meanwhile, atomic units can be textually 
encoded, and aggregate structure likewise notated through 
syntactic conventions that preserve atoms' boundaries, yielding 
textual serializations of data structures that can be 
shared between computing environments, allowing information to 
be copied and distributed.  Finally, digital representations 
of data structures can be marshalled into different binary forms \mdash{} 
encoded in byte- and bit-patterns \mdash{} to enable both 
persistent storage in a database and \q{runtime} presence 
as binary data that may be accessed by software applications.
}
\p{I take the time to lay out these basic principles because I want 
to emphasize the different contexts where digital representations can be 
found: in particular, textual encoding and serialization; 
graphical/interactive displays for software users; application runtimes; 
and long-term database storage.  A given representation will morph and 
mutate to accommodate these different contexts.  Moreover, these 
contexts correspond to distinct technical specializations: database engineers 
have a different perspective on data structures than network engineers designing 
protocols for encoding and exchanging data between network endpoints; and 
application designers focused on optimal human-computer-interaction 
understand digital representations as visual and interactive phenomena, 
whereas application \i{developers} need to focus on how to 
properly encode data structures for binary runtimes.  These various 
perspectives all influence the theory and technology 
behind digital data representation \mdash{} a successful 
representational paradigm needs to adapt to the operatoonal 
requirements of engineers in each of these disicplines.
}
\p{Additionally, insofar as the point of digital archives is to encode 
empirical, \q{real-world} data, a proper representational protocol 
needs to promote a synergy between the information as humans understand 
it and the data structures recognized by the technology.  For instance, 
if a published data set shares scientific data, it should be 
represented in ways that preserve scientifically significant 
details \mdash{} any derivations, descriptions, or observations 
which are intrinsic to the science's methodology, 
theory, assumptions, and experimental results.  The data needs to 
be structured according to the \q{semantics of the science}, so that 
the scientific background can be reconstructed along with future use 
of the data, even after a circuitous journey through different 
contexts, like through a database and over a network to 
a scientific-software application (maye years later).   
As formal models of data semantics have become more rigorous \mdash{} e.g. in our 
century with Ontologies and the \q{Semantic Web} \mdash{} this 
\q{semantic engineering} has become a further technical perspective 
needing consideration in the design and evaluation of digital representations.
}
\p{Over the decades, many digital representation strategies and 
\q{layouts} have been envisioned, from tabular structures 
in the manner of Relational Databases, to tree-like 
documents whose morphology is inspired by markup languages, 
to structurally looser variations on row/column 
arrangements like multi-dimensional key-value spaces or 
\q{Big Column} and other \q{\NoSQL{}}-database-inspired 
articulations.  These various paradigms are subject to 
selective pressures based on how well they meet the 
different engineering needs I have identified.  But the 
multiplicity of requirements complicates the \q{competition} between 
representational strategies: a paradigm optimized 
for one context (e.g., persistent databases) is not necessarily 
optimal for another (e.g., application and \GUI{} development).  
As a result, developers and computer scientists must 
continue to explore and collaborate on new paradigms which 
work better across contexts.
}
\p{In the plurality of representational paradigms, a significant 
evolution has been the emergent popularity of structures 
based on graphs.  The predominant influence on this 
line of development has been the Semantic Web and the specific graph 
architecture codified by \q{Web Ontologies}, although 
other graph-representation models (such as 
Conceptual Graph Semantics) were also candidates for 
potential technological \q{entrenchnment} before 
Semantic Web formats like \RDF{} and \OWL{} 
became standardized.\footnote{Why \RDF{}/\OWL{} and not, say, Conceptual 
Graphs, or the hybrid Object-Database/Graph-Database models 
studied in the 1990s, became canonized in the Semantic Web, is 
an interesting question but perhaps of mostly historical 
interest insofar as Hypergaphs can unify each of these paradigms.
}  More recently, scientists have proposed generalizations 
and/or alternatives to the Semantic Web based on 
hypergraphs in lieu of ordinary (directed, labeled) graphs.
}
\p{As might be ascertained from hypergraphs being the focus of this paper, 
I believe that representational paradigms based on hypergraphs can be superior to 
other formats and need to be studied with an integrated, generalized 
set of theories and computer tools.  I have introduced the 
topic of hypergraph data structures via general issues in 
digital representation in part to rest claims of hypergraphs' 
merit on explicit criteria.  In particular, I believe 
hypergraphs can adapt to the different technological 
contexts prerequisute to a decentralized digital ecosystem 
\mdash{} databases, networking, application implementation, visual/interactive 
software design, and semantic expressiveness \mdash{} more effeectively 
than other paradigms, like regular graphs or \SQL{}-style tables.
}
\p{I suspect other scientists and engineers have 
similar intuitions, because there has been a recent uptick in 
research on hypergraphs in various disciplines, such 
as Category Theory, Information Management, Artificial Intelligence, 
and Natural Language Processing.  Compared to 
the Semantic Web, however, there is a noticeable lack 
of standard tools or formats for expressing hypergraph 
data or sharing hypergraph structures across multiple 
applications and environments.  Whereas \RDF{} and \OWL{} are 
definitively associated with the Semantic Web, so that 
supporting these standards is a basic entry point for 
Semantic Web technologies, there is no comparable 
consensus on the underlying theory of hypergraph data 
in general.    
}
\p{This situation may be explained in part by subtle differences on 
what the word \q{hypergraph} is supposed to mean in different 
settings, so the overal terrain of hypergraphs is 
divided into distinct mathematical models, often without a 
rigorous theory of their interrelationships.  
Another problem is perhaps a failure to appreciate how hypergraphs 
are structurally different from ordinary graphs, so that 
hypergraph-shaped data may be imprecisely treated as 
a mere variation upon or a special structuration within 
ordinary graphs.  For eample, Semantic Web structures such as 
\RDF{} \q{bags} and \q{collections} introduce a kind of 
hierarchical organization to Semantic Web graphs \mdash{} qualifying as 
a form of hypergraph \mdash{} but these protocols are not described as 
enabling a transition from a networking framework based on 
directed labeled graphs to one based on hypergraphs.  Instead, 
hypergraphs become \i{de facto} embedded within ordinary graphs, 
exploiting the representational flexibility which also 
makes the Semantic Web suitable for spreadsheets, \XML{}, and 
other structures that are not graphs at a \i{prima facie} level.
The problem is that while hypergraphs can be \i{encoded} in 
ordinary graphs with a suitable labeling convention, 
their distinct structural advantages are lost as 
explicit architectural features once hypergraphs are 
\q{encoded} in conventional Semantic Web formats.
}
\p{As I will outine in the first section, there are at least some half-dozen 
different structures that may certainly be described as hypergraphs, generalizing 
various underlying graph models (or indeed, even if we restrict attention to 
labeled directed graphs).  Within this space of possibilities, 
recent computer-science-related) research into hypergraphs 
seems to emphasize two different topics: first, the 
representational potential for hypergraphs as a 
general morphology for encoding structured information-resources 
(so as a general tactic for data warehousing and modeling); and, second, the 
specific possibilities of using hypergraphs to model computer code and computational 
processes.  In the second subject-area, hypergraphs are studied 
as a medium for expressing details about computer algorithms as a 
way to reason about executable computer code as a structured system, and 
perhaps even to design execution engines (virtual machines to 
run computer code that is in a suitable representation).  This 
latter research sometimes appears to present a mathematical 
picture of hypergraphs as an abstract model 
of computing procedures and evaluations \mdash{} a kind of 
graph-based interpretation of the lambda calculus \mdash{} but 
sometimes is also marshaled into implemented execution 
environments, as in the OpenCog and lmnTal frameworks.
}
\p{My goal in this paper is to consider what a unified hypergraph 
paradigm can look like \mdash{} how the different strands which in their 
own way embody hypergraph structures can be woven together 
into a multi-purpose whole.  My emphasis is on computer implementations 
rather than mathematics \mdash{} that is, I will not present axioms 
or theorems formalizing different sorts of hypergraphs, though I 
ackowledge that such descriptions are possible.  Instead, my 
aim is to describe what might constitute a general software 
library or toolkit that can adapt to different hypergraph 
models and use-cases.  This paper is accompanied by a \q{data set} 
which involves a library (written in \Cpp{}) for creating and 
manipulating hypergraphs of different varieties, and also a 
\q{virtual machine} for modeling and realizing 
computatational procedures via hypergraph structures.  The 
point of the accompanying code is to 
demonstrate that a generalized hypergraph representation 
is possible, and that in addition to use-cases for modeling 
data structures such a representation can be used at the 
core of a virtual processor.  The code includes simple tests and 
\q{scripts} that can be executed via the provided virtual-machine 
code.
}
\p{Studying hypergraphs from a computational (and 
\q{implementational}) perspective \mdash{} not just as 
mathematical objects \mdash{} introduces useful details 
that can add depth to our overall understanding of 
hypergraphs.  For example, one question is how hypergraphs 
can be initialized \mdash{} how software systems can build up 
hypergraph structures, piece by piece.  This is related to the question of 
proper serialization formats for hypergraphs.  Given some 
specific hypergraph data aggregate \Hy{}, it is important to 
have a textual encoding where \Hy{} can be mapped to a 
character-string, shared as a document, and then reconstructed 
with the same hypergraph structure.  This raises the question of 
when and whether two hypergraphs are properly isomorphic 
\mdash{} such that serializing and then deserializing a 
hypergraph yields \q{the same} hypergraph \mdash{} and also the problem 
of validating and parsing textual representations of hypergraphs.  
The theory of parsing \i{serializations} of hypergraphs \mdash{} 
textual encodings constructed according to a protocol wherein their 
grammar and morphology are suitable to 
expressing and then rebuilding hypergraph structures \mdash{} 
then becomes an extension of hypergraph theory itself.
}
\p{In addition to creating hypergraphs via textual 
representations in the manner of markup languages 
(like \XML{} or \RDF{}), we can also consider the incremental 
accumulation of hypergraph structures via minimal 
units of transformation, providing a kind of 
\q{Intermediate Representation} embodying the form of a 
hypergraph via a sequence of effectively (at least on one scale 
of analysis) atomic operations.  For any variety of 
hypergraph, then, we can consider which is a suitable 
Intermediate Representation for conformant spaces \mdash{} in particular, 
which set of primitive options can, iteratively, construct 
any representative of the space of instantations of a 
given kind of hypergraph.  Insofar as a software framework 
seeks to work with several hypergraph varieties, we can consider 
how several different such operation-sets can be combined.  
}
\p{Moreover, we can consider both serialization and Intermediate Representation 
as parallel tactics for building hypergraphs.  
An Intermediate Representation language can be used to carry out 
transformations between hypergraphs \mdash{} producing \IR{} code based on 
the first hypergraph from which the second can be populated.  In conjunction 
with serialization and parsers, a hypergraph can then be realized 
via several stages, with one form of hypergraph used as an intermediary 
because it has a convenient grammar and works with a parsing engine; 
from that preliminary graph a new graph can be created via \IR{} code.  
Grammars, serializaion and deserialization, and \IR{} thereby become 
part of the formal architecture defining a hypergraph ecosystem and 
inter-graph transformations (analogous to the trio of \XML{} parsers, 
\XML{} transforms, and the \XML{} Document Object Model).
}
\p{Continuing the \XML{} analogy, note that \XML{} is not just a serialization specification: 
the full \XML{} technology defined requirements on a series of tools 
operating on \XML{} data, not just \XML{} files: tools 
for traversing \XML{} documents (including a programming interface 
for how traversal options are to be operationalized, e.g. 
via Object-Oriented methods like \parentNode{} and \nextSibling{}); 
for parsing \XML{} files into traversable data structures (including 
requirements on the traversals which the derived structurees 
\mdash{} the so-called \q{post-processing infoset} \mdash{} must 
support); and for inter-document mapping 
(via \XSLT{} or \XML{}-FO).  The Document Object Model and 
post-processing infoset is the structural core of the \XML{} 
format, even more than any surface-level syntax (the grammar 
for tags, attributes, and so forth).  Analogous specifictions 
would have to be developed for a generalized Hypergraph representation.
}
\p{The code discussed here does not purport to provide a mature or 
decisive implementation of a general-purpose hypergraph ecosystem, 
but rather to demonstrate by example what components may comprise 
such a framework and how they may interoperate.  The code includes 
hyprgraph models but also the preliminary and 
intermediary structures that can connct hypergraphs to 
a surrounding computational infrastructure \mdash{} parsers, 
Intermediate Representation, application-integration 
logic (such as mapping hypergraph nodes to 
application-specific data types), and a \q{virtual machine} 
for realizing hypergraphs as executable code.  
}
\p{The remainder of this paper will focus on three 
subjects: delimitating the proper range of hypergraph 
models; execution models and the \q{virtual machine}; and 
a review of formal structurs (like grammars and \IR{} 
formats) which are structurally different 
than hypergraphs but can help tie hypergraph 
representations togther into a unified ecosystem.  
I will also, in conclusion, sketch arguments to the effect that 
hypergraphs offer a plausible foundation for reasoning 
about structured/computation data and date types in general.  
I have spoken only informally about hypergraphs 
themselves, taking for granted that we have a general 
picture of what hypergraphs are, but this now 
demands more rigorous treatment.  There are actually several 
different research and engineering communities 
that talk about hypergraphs, in each case describing 
some sort of generalization of regular (often 
labeled and/or directed) graphs, but these generalizations 
fall into different kinds.  Hence there are several different 
varieties of hypergraph, and I will try 
to outline a general theory of hypergraphs in these 
various forms.
}
\p{}

\input{section1.ngml}
\section{The Illogic of Syntax}
\p{As I understand it, a non-trivial truth-theoretic semantics 
requires more than a holistic association between 
sentences and propositional content: it requires that this 
association be established \i{by linguistic means} and 
\i{on linguistic grounds} (syntax, semantics, pragmatics).  
I will present several arguments against this possibility, in 
the general cases \mdash{} that is, against the possibility that 
for \i{typical} sentences we can analyze syntactic form through 
the lens of the logical structure of propositions signified 
via a sentence; or analyze natural-language semantics through a 
logically well-structured semantics of propositions.  
I will emphasize to issues: first, that the architecture of 
linguistic performances \i{does not}, in the general case, 
\i{recapitulate propositional structure}; and, second, 
that language acts work through gaps in logical specificity that 
complicate how we should theorize the triangular relation between 
surface language, propositional content, and side-effect meanings.
}
\p{}
\p{}

\section{Channels and Carriers}  
\p{\phantomsection\label{sThree}The term \q{channels of communication} is usually associated with
process calculi, wherein channels pass values between procedures,
which in turn are combined sequentially or in parallel.
The possibility of \q{bidirectional} channels between
concurrent processes gives channels more structure, and
points to the wider scope of process calculi, compared to
lambda calculi.  There is a well-known \q{embedding} of
(classical) \mOldLambda{}-Calculus into process calculi, although
there does not appear to be comparable quantities of
research into process-algebraic
interpretations of \q{expanded} \mOldLambda{}-calculi with
objects, exceptions, and so forth.  Since process algebras
and calculi prioritize the analysis of computations
that run in parallel \mdash{} and the algebraic combinations
of procedures, such as concurrent vs. sequential \mdash{} the
narrower themes of \mOldLambda{}-abstraction in one single
procedure may seem tangential to serious process analysis.
}
\p{If real, such an assumption is unfortunate,
because the underlying semantics
of communicating values between disparate computing procedures
\mdash{} even apart from issues of concurrency, or from issues
concerning whether two superficially different \mOldLambda{}-expressions
are structurally equivalent (i.e., apart from the
main analytic themes of process- and \mOldLambda{}-calculi,
respectively) \mdash{}
is important in its own right.  Any value passed between
procedures \i{may} be reinterpreted as belonging to a
different type (the number \litOne{} can be an integer in one
context that gets interpreted as the decimal \litOnePtZ{} somewhere
else), which \i{may} even cause a modification (\litOnePtO{} gets
truncated to just \litOne{}, say, if the decimal is passed
to a function expecting an integer).  Moreover, inter-procedure
communications \i{may} be subject to gatekeeping and/or overloading
which blocks the target procedure from running if the
passed values are incorrect (relative to some specification) or
which selects one from a set of possible procedures \mbox{to run based
on the values passed for each occasion}.
}
\p{These processual issues give rise to different operators than conventional
process-calculi, but they can still be introduced as algebraic
structures on groups of procedures.  Let \piOne{}, \piTwo{}, and \piThree{}
represent three procedures; we can use notation like \pisbl{} to represent
a \q{guarded} inter-procedure combination, wherein a gatekeeper is
in effect that may \i{block} \piTwo{}.  We may also use notation like
\pisfo{} to represent a \q{polymorphic} inter-procedure combination,
wherein a gatekeeper selects one of \piTwo{} or \piThree{} as the
proper procedure to target based on passed values.  Note that the former
case is a special example of the latter if we assign \piThree{} in the
first notation as a \q{null} procedure, \makebox{which does nothing,
having no actions or effects}.
}
\p{These inter-procedure relations have nontrivial semantics
even if we do not include concurrency or bidirectional
communication, so that \q{channels} are closer in spirit to
lambda abstraction than to stateful \q{message-lines} as in
the Calculus of Communicating Systems.
As the expansion of lambda calculi
toward Object-Orientation (in the \q{Sigma} or \sigmaCalculus{}) and
other variations shows, even a simpler notion of channels,
generalized from lambda abstractions (from which
bidirectional mutable channels can then be modeled) is
not uninteresting.  In the literature,
criteria like \q{guarded choice}, supplementing
underlying process calculi, appears to be understood
principally at the level of \i{procedures} \mdash{} for one example,
the case of \q{choosing} one of several procedures is
similar to an operator that unifies multiple procedures
into \i{one} procedure with an \q{either or} logic:
\q{\piOne{} \i{or} \piTwo{}} is a procedure which may on some
occasions execute as \piOne{} and elsewhere as \piTwo{}
(see for instance \cite{EricWalkingshaw}).  With
this either-or logic defined on procedures, inter-procedure combinations
(where one procedure sends values to a second, causing the second to begin)
which are guarded and polymorphic can be modeled as combinations
where the second procedure is an either-or sum of multiple
more-specific procedures (possibly including one \q{do nothing} no-op).
}
\p{This model, however, neglects the details which guide how either-or
choices are understood to be made, according to the intended
theoretical models.  An important class of guarded-choice cases is
where choices are made on the basis specifically of the \i{values}
sent between procedures \mdash{} say, dispatching to different
function-bodies according to whether or not a number is in a range
\rRan{}.  These cases can perhaps be described as \i{localized guarded choice}
because all information pertinent to the guard is derived from
the passed values (and not from any global or ambient state).
We can further identify \i{constructor-localized} guarded choice as
cases where the \q{guards} are simply the value-constructors for
values passed in to the target procedure (maybe cast).  In
a language like \Cpp{} one or more \q{hidden} functions can execute
before a called function actually begins, insofar as
\q{copy constructors} may be in effect, controlling how values are copied
from one function's context to another.  Without deliberately
modeling process calculi at a technical level, then, \Cpp{} does
provide a \q{hook} where various process-related ideas can be implemented.
}
\p{Insofar as Localized Guarded Choice bases guarded-choice semantics
solely on inter-procedure passed values, I would argue that it lies
in a theoretically intermediate position between \mOldLambda{}-style and
process calculi.  From one perspective, Localized Guarded Choice is a
variation on inter-procedure combination, wherein one already-executing
procedure initiates a second by passing one or more values between them.
This perspective is perhaps the more intuitive when seen from the
point of view of the prior procedure.
}
\p{However, Localized Guarded Choice is
also a variation on \mOldLambda{}-abstraction, because it presumes that
the abstracting of some symbol in the target procedure's implementation
is not abstracted indiscriminately, but rather can only
be furnished with values conformant to some
specification.  We can envision a \q{Local Guard Calculus}, maybe
dubbed \mGamma{}-calculus, which works on the assumption that
for each \mOldLambda{}-abstraction there is a corresponding
value constructor which executes as a procedure before a target
procedure is called (and may even prevent the target
procedure from starting, e.g. by throwing an exception).  From the
perspective of the prior, calling procedure, this hidden execution
appears to be an added procedural layer, an intermediary
function that lies between itself and the target.  From the
perspective of the target procedure, however, the intermediary
value-constructors appear more as logical guarantees keyed to
its implementational specifications \mdash{} i.e., as abstraction-refinements.
In short, depending on perspective, this hypothetical
Local Guard Calculus can be seen as either a special case of
Lambda or Process calculi, respectively.
}
\p{Since the relevant issues of gatekeeping and overloading are
\q{intermediate} in this sense, they are not central foci of
either genre of calculi, and call for a distinct terminology and
conceptualization.  I will accordingly think of function implementations
as for practical purposes \q{procedures} which communicate via
\q{channels}, but my sketch of \q{Channel Algebra} is not a
process calculus or process algebra \i{per se}, and will
skirt around foundational process-oriented topics like concurrency,
or the possibility of stateful, bidirectional channels.
Moreover, my ambition is to
model source code and behavior via semantic graphs which
can fit into the \RDF{} (and \DH{}/\CH{}) ecosystems,
a reconstruction that can
start at the simpler non-concurrent, single-threaded level.
So for the duration of this section I
will attend exclusively to the semantics of
functions passing values to other functions so as to
initiate the procedures which the target functions implement;
setting aside considerations such as whether the prior
function \q{blocks} until the second function completes
(which represents sequential-process semantics) or,
instead, continues running, concurrently.
}
\decoline{}
\p{Suppose one function calls a second.  From a high level perspective,
this has several consequences that can be semantic-graph represented
\mdash{} among others, that the calling function depends on an
implementation of the callee being available \mdash{} but at the
source code level the key consequence is that a node representing
source tokens which designate functional values enters into
different semantic relations (modeled by different kinds
of edge-annotations) than nodes marking other types
of values and literals.  Suppose we have an edge-annotation
that \nodex{} is a value passed to \nodef{}; this graph is only
semantically well-formed if \nodef{}'s representatum has
functional type (by analogy to the \mbox{well-formedness criteria, at
least in typed settings, of \lambdaxfx{})}.
}
\vspace{-1em}
\p{This motivates the following: suppose we have a Directed
Hypergraph, where the
nodes for each hyper-edge represent source-code tokens (specifically,
symbols and literals).  Via the relevant Source Code Ontology, we can
assert that certain
edge-annotations are only possible if a token (in subject or object position)
designates a value passed to a function.  From the various edge-annotation
kinds which meet this criteria, we can define a set of \q{channel kinds}.
}
\p{For every function called, there is a corresponding
function implementation which has its own graph representation.
Assume this implementation includes a \i{signature} and a \i{body}.
With sufficient compiler work, the body can be expressed as an
Abstract Syntax Tree and, with further analysis, an
Abstract Semantic Graph \mdash{} which in turn will identify semantic
connections such as the scoped relation between a symbol
appearing in the signature and the same symbol in the
body \mdash{} these symbols will all share the same
value (unless a symbol in a nested lexical scope \q{hides}
the symbol seen in the parent scope).
}
\p{This implicitly
assumes that symbols \q{hold} values; to make the
notion explicit, I will say that symbols are
\i{carriers} for values.  It may be appropriate
to consider literals as carriers (as demonstrated by
the \Cpp{} \operatorqq{}, the conversion of character-string
literals to typed values is not always trivial, nor
necessarily fixed by the language engine).  For now, though,
assume that carriers correspond to lexically-scoped symbols
(I also ignore globals).  Carriers do not necessarily hold a
value at every point in the execution of a program; they
may be \q{preinitialized}, and also \q{retired} (the
latter meaning they no longer hold a meaningful value;
consider deleted pointers or references to out-of-scope
symbols).  A carrier may pass through a
\q{career} from preinitialized to initialized, maybe then
changing to hold \q{different} values, and maybe then
retired.  I assume a carrier is associated with a single type
throughout its career, and can only hold values appropriate for its
type.  The variety of possible
careers for carriers is not directly tied to its type: a
carrier which cannot change values (be reinitialized)
is not necessarily holding a \const{}-typed value.
}
\p{Having introduced the basic idea of \q{carriers} I will now consider
carrier operations in more detail, before then expanding on the
theory of carriers by considering how carriers group into channels.
}
\spsubsection{Carriers and Careers}
\p{In material fact, computers work with \q{numbers} that
are electrical signals; but in practice, we reason about
software via computer code which abstracts and is
materially removed from actual running programs.  In these
terms, we can note that \q{computers} (in this
somewhat abstract sense) do not deal \i{directly} with
numbers (or indeed other kinds of values), but rather
deal with \q{symbols} that express, or serve as placeholders
for, or \q{carry} values.  Even \q{literal} values are composed
of symbols: the number (say) \oneTwoThree{}, which has a particular
physical incarnation in electrons, is represented in source
code by a Unicode or \ascii{} character string (which materially
bears little resemblance to a \q{digital} \oneTwoThree{}).
}
\p{This primordial fact is orthogonal to type systems.  The underlying
notion of type theory is of \i{typed values} \mdash{} the idea
that the whole universe of values that may be encountered
during a computation can be classified into types, so
a particular value always has one single declared type
(though it may be a direct candidate for classification
as other types as well).  There is a core distinction
between \i{types} and \i{sets}: two distinct types can
be instantiated by the same set of possible values, and
one type can cover a different spectrum of values on different
computers.  For example, any type which has an expandable data collection
\mdash{} say, a list of numbers \mdash{} can have values which
would take ever larger amounts of memory; e.g., as lists grow
larger (I made essentially the same
comment about pointer-lists on page
\hyperref[pointerlist]{\pageref{pointerlist}}).
Since different computers have different available
memory at different times, there is no fixed relation between
the set of \q{inhabitants} of these types which can be
represented on different computers, or on one single computer
at different times.  Type theory therefore belongs to
a logicomathematical (or philosophical) foundation distinct from
set theory.  Instead of defining types from sets of values,
types are instead built up from more primitive types via
several operators, most notably \q{products}, or \q{tuples}
wherein a fixed number of values of (in the general case)
distinct types are grouped together
(as I outlined from one
perspective within subsection \sectsym{}\hyperref[types]{1.4}).
}
\p{But type theory still consider these \q{values} as abstract, primitive
constructs.  Both types and values are in essence undefinable concepts,
for which there are no deeper concepts in terms of which they can be defined
(although this de-facto non-definition can be given a rather
elegant gloss with Category Theory).  Recognizing that the type-value
pair does not itself include the medium through which a value
is \i{represented} in computer code, we can say that an equally
primordial concept can be (say) a \i{carrier}, a symbolic expression
which \q{holds} a value, or the anticipation of one.  So we introduce
\q{carriers} as a third fundamental concept.  All carriers, in
this theory, have one single (declared) type; and all carriers
can hold one (or in some contexts many) values; so there is
an interrelationship between the notions \i{type}, \i{value},
and \i{carrier}.
}
\p{Carriers come in two broad groups: \i{literals} (like \qOneTwoThree{})
which embody a single value everywhere they occur in source
code, and \i{symbols} (like \xSymbol{})
which can have different values when they
occur in different places in source code.  There may be many
\q{instances} of a single source code symbol in one code base,
each of which, if they occur in different \q{scopes}, may
have their own value unrelated to the others.  In general one
symbol might carry different values at different time,
though some programming environments may restrict this variability.
}
\p{\i{Immutable} symbols are tightly bound to their values; they have
one single value for the full time that they have any value
at all, and therefore act as \q{transparent} proxies for their value
\mdash{} especially if languages discourage constructs like
non-initialized values and \addressOf{} operators (which
can cause symbols to hold no-longer-meaningful values).
}
\p{Because \q{unitialized} carriers and \q{dangling pointers}
are coding errors, within \q{correct} code, carriers
and values are bound tightly enough that the whole carrier/value
distinction might be considered an artifact of programming practice,
out of place in a rigorous discussion of programming languages
(as logicomathematical systems, in some sense).  But even
if the \q{career} of symbols is unremarkable, we cannot avoid
in some contexts \mdash{} within a debugger and/or an \IDE{}
(Integrated Development Environment, software for writing programs),
for example \mdash{} needing to formally distinguish the carrier
from the value which it holds, or recognize that carriers
can potentially be in a \q{state} where, at some point in which
they are relevant for code analysis or evaluation, they
do not yet (or do not any longer) hold meaningful values.
The \q{trajectory} of carrier \q{lifetime} \mdash{} from being declared,
to being initialized, to falling \q{out of scope} or otherwise \q{retired}
\mdash{} should be integrated into our formal inventory of programming
constructs, not relegated to an informal \q{metalanguage} suitable
for discussing computer code as practical documents but not as formal systems.
I'd make the analogy to how lambda calculus models \q{symbol rewriting}
as a formal part of its theory, rather than a notational intuition
which reflects mathematical conventions but is seen as part
of mathematical \q{metalanguage} rather than mathematical \q{language}.
The history of mathematical foundations reveals how convention, notation,
and metalanguage tends to become codified into theory, models, and
(mathematical) language proper.  Likewise, we can develop a formal
theory of carriers which models that carriers have different states
\mdash{} including ones where they do not hold meaningful values
\mdash{} separate and apart from whether the underlying type
system allows for mutable value-holders, or pointers, references
to lexically scoped symbols, and other code formations that
can lead to \q{undefined behavior}.
}
\p{Consider a carrier which has been declared, so it is assigned a type and belongs to
some given scope in its source code, but has not yet been initialized.
The carrier does not therefore, in the theory, hold any value at all.  This
is different from holding a default value, or a null value, or
any other value that a type system might introduce to represent
\q{missing} information.  In practice, of course, insofar as a \q{carrier} refers
to a segment of memory, the carrier will be seen by the computer as having
\q{some} (effectively random) value; some bit pattern left behind by some
other calculation, which obviously has no \q{meaning}.  Programming languages
usually talk about the (mis)use of uninitialized variables (i.e., in this
context, carriers) as leading to \q{undefined behavior} and therefore beyond any
structured reasoning about code, and leave it at that.  We can be more rigorous
however and define being uninitialized as a \i{state} of a carrier, as is
holding a meaningful value (once it \i{is} initialized) and then, some time
later, no longer holding a meaningful value, because its associated memory
has been deemed recyclable (typically, once a symbol falls out of scope).  Any
carrier therefore has at least three kinds of state, which we can
call \i{pre-initialized}, \i{initialized}, and \i{retired}.\footnote{Again, though,
we should not think of (say) \q{Preinitialized} as a \i{value} with a \i{type}.
A carrier declared as type \intthrtwo{}, say (32-bit integer) can only hold
values of this type, when it holds any value at all, which precludes
interpreting \q{Preinitialized} as some kind of \q{null} value which is somehow
arranged to be an instance of \intthrtwo{}.
}
}
\p{In this theory, carriers are the basic means by which values are represented
within computer code, including representing the \q{communication} of
values between different parts of code source, such as calling a function.
The \q{information flow} modeled by a function-call includes values
held by carriers at the function-call-site being transferred to
carriers at the function-implementation site.  This motivates
the idea of a \q{transfer} of values between carriers, a kind of primitive
operation on carriers, linking disparate pieces of code.  It also
illustrates that the symbols used to name function parameters, as
part of function signatures, should be considered \q{carriers} analogous
to lexically-scoped and declared symbols.
}
\p{Taking this further, we can define a \i{channel} as a list of carriers which, by
inter-carrier transfers, signify (or orchestrate) the passage of data into and out
of function bodies (note that this usage varies somewhat from process calculi,
where a channel would correspond roughly to what is here called a single carrier;
here channels in the general case are composed of multiple carriers).
I'll use the notation \opTransfer{} to represent inter-carrier transfer: let
\carrOne{} and \carrTwo{} be carriers, then \carrOneOpTransferTwo{} is a
transfer \q{operator} (note that \opTransfer{} is non-commutative; the
\q{transfer} happens in a fixed direction), marking the
logical moment when a value is moved from code-point to code-point.
The \opTransfer{} is intended
to model several scenarios, including \q{forced coercions} where the associated
value is modified. Meanwhile, without further details a \q{transfer} can be
generalized to \i{channels} in multiple ways.
If \carrOne{} and \carrTwo{} are carriers which belong to two channels
(\chanOne{}, \chanTwo{}), then \carrOneOpTransferTwo{} elevates to
a transfer between the channels
\mdash{} but this needs two indices to be concrete: the notation has to
specify which carrier in \chanOne{} transfers to which carrier in \chanTwo{}.
For example, consider the basic function-composition \fofg{}:
\fDotOfGX{} = \fOfGx{}.  The analogous \q{transfer} notation
would be, say, \gOpTransferOneOneF{}:
here the first carrier in the \returnch{} channel
of \funG{} transfers to the first carrier in the \lambdach{} channel of \funF{}
(the subscripts indicate the respective positions).
}
\p{Most symbols in a function body (and corresponding nodes in a
semantic graph) accordingly represent carriers, which
are either passed in to a function or lexically declared
in a function body.  Assume each function body corresponds
with one lexical scope which can have subscopes
(the nature of these scopes and how they fit in
graph representation will be addressed later in this section).
The \i{declared} carriers are initialized with
values returned from other functions (perhaps the
current function called recursively), which can
include constructors that work on literals (so, the
carrier-binding in source code can look like
a simple assignment to a literal, as in \intieqzero{}).
In sum, whether they are passed \i{to} a  function or
declared \i{in} a function, carriers are only initialized
\mdash{} and only participate in the overall semantics
of a program \mdash{} insofar as they are passed to other functions
or bound to their return values.
}
\p{Furthermore, both of these cases introduce associations between
different carriers in different areas of source code.  When a carrier
is passed \i{to} a function, there is a corresponding carrier
(declared in the callee's signature) that receives the
former's value: \q{calling a function} means
transferring values between carriers present at the site of
the function call to those present in the function's
implementation.  Sometimes this works in reverse: a function's
return may cause the value of one of its carriers to be
transferred to a carrier in the caller (whatever
carrier is bound to the caller's return value).
}
\p{Let \carOne{} and \carTwo{} be two carriers.  The \opTransfer{} operator
(representing a value passed from \carOne{} to \carTwo{}) encompasses
several specific cases.  These include:
\begin{enumerate}\eli{}  Value transfer directly between two carriers in one scope, like \aeqb{} or \aceqb{}.
\eli{}  A value transferred between one carrier in one function body when the return
value of that function is assigned to a carrier at the call site, as in \yeqfx{}
when \fSym{} returns with \retFive{}, so the value \five{} is transferred to \ySym{}.
\eli{}  \label{transf}A value transferred between a
carrier at a call-site and a carrier in the
called function's body.  Given \yeqfx{} and \fFun{} declared as, say,
\fIntI{}, then the value in carrier \xSym{} at the call-site is transferred to
the carrier \iSym{} in the function body.  In particular, every node
in the called function's code-graph whose vertex represents a source-code
token representing symbol \iSym{} then becomes a carrier whose value is
that transferred from \xSym{}.
\eli{}  A value transferred between a \returnch{} channel and either a
\lambdach{} or \sigmach{} channel, as part of a nested expression or a
\q{chain of method calls}.  So in \hfx{}, the value held by the carrier
in \fFuns{}'s \returnch{} channel is transferred to the first carrier in
\hFun{}'s \lambdach{}.  An analogous \returnch{}\opTransfer{}\sigmach{}
transfer is seen in code like \fxdoth{}: the value
in \fFuns{}'s \returnch{} channel becomes the value in \hFun{}'s \sigmach{},
i.e., its \q{\this{}} (we can use \opTransfer{}
as a notation between \i{channels} in this case because we understand
the Channel Algebra in force to restrict the size of both
\returnch{} and \sigmach{} to be at most one carrier).
\end{enumerate}
}
\vspace{-1em}
%Define two narrower operators
%modeling a transfer of values between them: \carOnetoTwos{}
%means direct assignment (between carriers in the same
%scope, as in \aeqb{});
\p{Let \carOnetoTwof{} be the special case of \opTransfer{}
corresponding to item \hyperref[transf]{(3)}: a transfer
effectuated by a function call, where
\carOne{} is at the call site and \carTwo{} is part of a
function's signature.  If \fOne{} calls \fTwo{} then
\carOne{} is in \fOne{}'s context, \carTwo{} is in \fTwo{}'s
context, and \carTwo{} is initialized with a copy of
\carOne{}'s value prior to \fTwo{} executing.
A \i{channel} then becomes a
collection of carriers which are found in the
scope of one function and can be on the
right hand side of an \carOnetoTwof{} operator.\footnote{In general, we might say that two carriers are \q{convoluted}
if there is a value passed between them or if their values
somehow become interdependent (as an example of
interdependence without direct transfer, consider tests that
that two lists are the same size, or two numbers
monotone increasing, as a runtime disambiguation of
dependent-typed polymorphic implementations).  Depending
on context, convolution can refer to a structure in
program graphs in the abstract or to an event in
program execution: modeling a program as a series
of discrete points in time \mdash{} each point inhabited by a
small change in program state \mdash{} two carriers are convoluted
at the time-point where a value-transfer occurs (or the
steps toward some kind of gatekeeping check get initiated).
}
}
\p{To flesh out Channels' \q{transfer semantics} further, I will refer back
to the model of function-implementations as represented in code
graphs.  If we assume that all code in a computer program is found
in some function-body, then we can assume that any function-call
operates in the context of some other function-body.  In particular,
any carrier-transfer caused by a function call involves a link between
nodes in two different code graphs (I set aside the case of recursive
functions \mdash{} those which call themselves \mdash{} for this discussion).
}
\p{So, to review my discussion in this section so far, I started with the
process-algebraic notion of inter-procedure combinations; but whereas process
calculi enlarge this notion by distinguishing different syntheses of procedures
(such as, concurrent versus sequential), I have focused instead on the
communication of values between procedures.  The semantics of inter-procedure
value-transfers is unexpectedly detailed, because it has to recognize
the possibility of nontrivial copy semantics, casts, overloading,
and perhaps \q{gatekeeping}.  Furthermore, in addition to these
semantic details, analysis of value-transfers is particularly significant
in the context of Source Code Ontologies and \RDF{} or
Directed Hypergraph representations
of computer code.  This is because code-graphs give us a rigorous foundation
for modeling computer programs as sets of function-implementations
which call one another.  Instead of abstractly talking about
\q{procedures} as conceptual primitives, we can see procedures as
embodied in code-graphs (and function-values as
constructed from them, which I emphasized last section).
\q{Passing values between} procedures is
then explicitly a family of relationships between nodes
(or hypernodes) in disparate code-graphs, and the various semantic
nuances associated with some such transfers (type casts, for example)
can be directly modeled by edge-annotations.  Given these
possibilities, I will now explore further how the framework of
\i{carriers} and \i{channels} fits into a code-graph context.
}
\spsubsectiontwoline{Channel Complexes, Code Graphs, and Carrier Transfers}
\p{For this discussion, assume that \fOne{} and \fTwo{} are implemented functions with
code graphs \gammaOne{} and \gammaTwo{}, respectively.  Assume
furthermore that some statement
or expression in \fOne{} involves a call to \fTwo{}.  There are several specific
cases that can obtain: the expression calling \fTwo{} may be nested in a
larger expression; \fTwo{} may be called for its side effects alone,
with no concern to its return value (if any); or the result of \fTwo{}
may be bound to a symbol in \fOne{}'s scope, as in \yeqfx{}.  I'll take
this third case as canonical; my discussion here extends to the other
cases in a relatively straightforward manner.
}
\p{A statement like \yeqfx{} has two parts: the expression \fx{} and the symbol
\ySym{} to which the results of \fFun{} are assigned.  Assume that this
statement occurs in the body of function \fOne{}; \xSym{} and
\ySym{} are then symbols in \fOne{}'s scope and the symbol \fFun{} designates
(or resolves to) a function which corresponds to what I refer to here as
\fTwo{}.  Assume \fTwo{} has a signature like \fIntI{}  As such, the
expression \fx{}, where \xSym{} is a carrier in the context of
\fOne{}, describes a \i{carrier transfer} according to which the value
of \xSym{} gets transferred to the carrier \iSym{} in \fTwo{}'s context.
}
\p{In the terms I used earlier, \fTwo{}'s signature represents a
channel \q{package} \mdash{} which, in the current example, has
a \lambdach{} channel of size one (with one carrier of
type \int{}) and a \returnch{} channel of size one (\fTwo{}
returns one \int{}).  Considered in the context of carrier-transfers
between code graphs, a Channel Package can be seen as a description
of how two distinct code-graphs may be connected via carrier
transfers.  When a function is called, there is a channel
complex which I'll call a \i{payload} that supplies values to
the channel \i{package}.  In the concrete example, the
statement \yeqfx{} is a \i{call site} describing a channel
\i{payload}, which becomes connected to a function
implementation whose signature represents a channel
\i{package}: a collection of transfers \carrOneOpTransferTwo{}
together describe an overall transfer between a
\i{payload} and a \i{package}.
}
\p{More precisely, the \fx{} example represents a carrier transfer whose
target is part of \fTwo{}'s \lambdach{} channel, which we can notate
\carrOneOpTransferTwolambda{}.  Furthermore, the full statement \yeqfx{}
shows a transfer in the opposite direction: the value in
\fTwo{}'s \i{\returnch{}} channel is transferred to the carrier \ySym{} in
the \i{payload}.  This relation, involving a return
channel, can be expressed with notation like \carrTwoOpTransferOnereturn{}.
The syntax of a programming language governs how code at a call site
supplies values for carrier transfers to and from a function body: in
the current example, binding a call-result to a symbol always
involves a transfer from a \i{\returnch{}} channel, whereas handling an
exception via code like \catchexce{} transfers a value from a called
function's \i{\exceptionch{}} channel.  The syntactic difference between
code which takes values from \returnch{} and \exceptionch{} channels,
respectively, helps reinforce the \i{semantic} difference between
exceptions and \q{ordinary} returns.  Similarly, method-call syntax
like \objfx{} visually separates the values that get transferred to a
\q{\sigmach{}} channel (\obj{} in this case)
from the \q{ordinary} (\lambdach{}) inputs, reinforcing Object semantics.
}
\p{To consolidate the terms I am using: we can interpret both function
\i{signatures} and \i{calls} in terms of channels.  Both involve
\q{carrier transfers} in which values are transferred \i{to} or \i{from}
the channels described by a function signature.  The distinction between
functions' \q{inputs} and \q{outputs} can be more rigorously stated,
with this further background, as the distinction between channels
in function signatures which receive values \i{from} carriers at
a call site (inputs), and those \i{from which} \mbox{values are obtained
as a procedure has completed (outputs)}.
}
%\vspace{-1em}
\p{A Channel Expression Language (\CXL{}) can describe channels both in
signatures and at call-sites.  The aggregation of channels generically
described by \CXL{} expressions I am calling a \i{Channel Complex}.
A Channel Complex representing a function \i{signature} I am calling
a \i{Channel Package}, whereas complexes representing a function
\i{call} I am calling a \i{Channel Payload}.  Input channels
are then those whose carrier transfers occur in the payload-to-package
direction, whereas output channels are the converse.
}
\vspace{1em}
\itcl{rz}
\vspace{-1em}
\p{The demo code for this chapter uses an Intermediate
Representation which builds complexes representing both
function signatures and dynamically evaluated
function-calls.  Figure \ref{lst:rz} shows
code in the special Interface Description Language
which includes a simple signature (\OneOverlay{}).
It also has several dynamic calls, including one to
create a new object of a \Cpp{} class called \FnDoc{}
(\TwoOverlayu{}) and one to retrieve a preconstructed
object of a class called \KCMEnv{} (\ThreeOverlay{})
(in this language the \dbleq{} operator is used
when the right-hand-side expression
is a value constructor, and there are assignment
operators, notated with a preceding back-slash as in
\Seq{} or \Sdbleq{}, specifically for assigning values
to previously-uninitialized symbols).
}
\vspace{1em}
\itcl{rzd}
\vspace{-1em}
\p{Corresponding to this sample, the Intermediate Representation
(actually Common \Lisp{}) in Figure \ref{lst:rzd},
to which the above code compiles, shows both channel
constructions for signatures (\OneOverlay{}) and
dynamic calls (\TwoOverlay{}, \ThreeOverlay{}, and
\FourOverlay{}).  These calls vary in the channel formations
used \mdash{} for instance, \FourOverlay{} uses a sigma
channel (and maps to a \Cpp{} method via the \Qt{} meta-object
system); whereas \ThreeOverlay{} uses a predefined
\Cpp{} function exposed to the scripting engine
(in this case one which retrieves a predefined \KCMEnv{}
object that is globally part of the scripting
environment; \envv{} stands for \q{environment value}),
and \TwoOverlay{} default-constructs an object of
the requested type (again, the \Qt{} code called from
this \Lisp{} form will employ \Qt{} meta-object code
to actually allocate the object, providing a type identifier based on
the type name).  These structural differences are reflected in how
the \Cpp{} callback methods are named (here visible indirectly by
the naming convention for some \Lisp{} symbols, which get
translated behind the scenes to
method names of a certain \Qt{}-based class) \mdash{}
\kprom{}, for example, maps to the particular \Cpp{} function
which the engine
uses to default-construct a value given a corresponding type name.
This code shows that channel complexes can be constructed
for several purposes, including but not limited to
dynamic evaluation (when evaluation is desired,
it can be triggered by \kcmde{}, as at \FiveOverlay{}).
}
\p{In addition to the payload/package distinction, we can also
understand Channel Complexes at two further levels.  On the one hand,
we can treat Channel Complexes as found \i{in source code},
where they describe the general pattern of payload/package
transfers.  On the other hand, we can represent Channel Complexes
\i{at runtime} in terms of the actual values and types held by
carriers as transfers are effectuated prior to, and then
after, execution of the called function.  Accordingly,
each Channel Complex may be classified as a \i{compile-time}
payload or package, or a \i{runtime} payload or package, respectively.
The code accompanying this chapter includes a \q{Channel Complex library}
\mdash{} for creating and analyzing Channel Complexes via a special
(\Lisp{}-based) Intermediate Representation \mdash{}
that represents complexes of each variety, so it can be used
both for static analysis and for enhanced runtimes and scripting.
}
\p{This formal channel/complex/package/payload/carrier vocabulary codifies
what are to some degree common-sense frameworks through which programmers
reason about computer code.  This descriptive
framework (I would argue) more effectively integrates the
\i{semantic} and \i{syntactic} dimensions of source code and program
execution (Figures \ref{lst:rzsyn} and \ref{lst:rzsem},
listed as an addendum the end of this section, sample the
syntactic \mdash{} or grammar/parsing \mdash{} and semantic/evaluative components
of the demo code compiler pipeline, respectively).
}
\p{Computer programs can be understood \i{semantically} in
terms of \mOldLambda{}-Calculi combined with
models of computation (call-by-value or
by-reference, eager and lazy evaluation, and so forth).  These
semantic analyses focus on how values change and are passed between
functions during the course of a running program.  From this
perspective, source code is analyzed in terms of the semantics
of the program it describes: what are the semantic patterns and
specifications that can be predicted of running programs on the
basis of source code in itself?  At the same time, source code
can also be approached \i{syntactically}, as well-formed
expressions of a formal language.  From this perspective,
correct source code can be matched against language grammars and,
against this background, individual code elements
(like tokens, code blocks, expressions, and statements) \mdash{} and
their inter-relationships \mdash{} may be identified.
}
\p{The theory of Channel Complexes straddles both the semantic and syntactic
dimensions of computer code.  Semantically, carrier-transfers capture the
fundamental building blocks of program semantics: the overall
evolving runtime state of a program can be modeled as a succession
of carrier-transfers, each involving specific typed values plus
code-graph node-pairs, marking
code-points bridged via a transfer.  Meanwhile,
syntactically, how carriers belong to channels \mdash{}
the carrier-to-channel map fixing carriers' semantics \mdash{} structures and
motivates languages' grammars and rules.  In particular,
carrier-transfers induce relationships between code-graph nodes.
As a result, language
grammars can be studied through code-graphs' profiles
insofar as they satisfy \RDF{}
and/or \DH{} Ontologies.
}
\p{In sum, a \DH{} and/or Semantic Web representation of computer code
can be a foundation for both semantic and syntactic analyses, and this
may be considered a benefit of Channel Complex representations
even if they only restate what are established semantic patterns
mainstream programming language \mdash{} for example, even if they
are restricted to a \sigmach{}-\lambdach{}-\returnch{}-\exceptionch{}
Channel Algebra modeled around, say, \Cpp{} semantics
prior to \CppEleven{} (more recent \Cpp{} standards also call for
a \q{\capturech{}} channel for inline \q{lambda} functions).
}
\p{At the same time, one of my claims in this chapter is that
more complex Channel Algebras can lead to new tactics for
introducing more expressive type-theoretic semantics in
mainstream programming environments.  As such, most of the
rest of this section will explore additional Channel Kinds
and associated Channel Complexes which extend, in addition
to merely codifying, mainstream languages' syntax and semantics.
}
%\p{%To flesh out Channels' \q{transfer semantics} further, it is necessary to
%expain the intersection between channels and types in more detail, which I will do by
%(rather casually) appproaching types through Category Theory.
%}
%\subsection{Enhancing Type Systems via Channel Kinds}
%\p{%}

\input{section3a.ngml}
\subsubsection{Channelized-Type Interpretations of
Larger-Scale Source Code Elements}
\p{By intent, Channel Algebras provide a machinery for modeling
function-call semantics more complex than \q{pure} functions
which have only one sort of input parameter (as in lambda
abstraction) \mdash{} note that this is unrelated to
parameters' \i{types} \mdash{} and one sort of (single-value) return.
%As just discussed,
Examples of a more complex paradigm
come from Object-Oriented code, where there are two varieties of
input parameters (\q{lambda} and \q{sigma}).
Consider method calls in Object-Oriented languages: these
typically have a special syntax with one
distinguished carrier.  This
carrier is in a sense privileged: the type of
the \q{\thisc{}} carrier establishes the class to which
function belongs, influencing when the function may be
called and how polymorphism is resolved.  Moreover,
\q{chaining} method calls means that the result of
one method becomes an object that may then receive 
another method (the following one in the chain).  Such 
chaining allows for an unambiguous function-composition operator:
since functions in general take multiple arguments, there is
no single operator to pass the result of one function to
another; but since most methods have one return value and
one \thisc{} object, it is straightforward to
define a method-chain operator.
}
\p{Another case-study
is offered by exceptions.  The option to \q{throw} an exception
can be considered an alternative kind of output channel.
A function throws an exception instead of returning a value.
As a result, \returnc{} and \errorc{} channels typically
evince a semantic requirement (which earlier \mdash{} see
the notations on page \hyperref[retexc]{\pageref{retexc}} \mdash{} I sketched
as an algebra stipulation): when functions
have both kinds of channels, only one
may have an initialized carrier after the function
returns.
Usually, thrown-exception values can only be bound to carriers
in \catchddd{} formations \mdash{} once held in a carrier
they can be used normally, but carriers in \exceptionch{}
channels themselves can only transfer values to
other carriers in narrow circumstances
(this in turn depends on things like code
blocks, which in turn will be reviewed below).
So \errorc{} channels are not a sugared
form of ordinary returns, any more than objects are sugar for
functions' first parameter; there are axiomatic criteria
defining possible formations of \errorrc{} and \returnrc{}
channels and carriers, criteria which are more transparently
rendered by recognizing \errorrc{} and \returnrc{} as
distinct channels of communication available
within function bodies.
}
\p{%This terminology serves as a basis for exploring the
%variation in how functions receive and share values.
In general, extensions to \lambdaCalculus{} are
meaningful because they model semantics other than
ordinary lambda abstraction.  For example, method-calls
(usually) have different syntax than non-method-calls,
but \sigmaCalculi{} aren't trivial extensions or
syntactic sugar for \lambdach{}s; the more significant difference is
that sigma-abstracted symbols and types have different
consequences for overload resolution and function
composition than \lambdach{}-abstractions.  Similarly,
exceptions interact with calling code differently
than return values.  These differences do not belong precisely
to \lambdaCalculus{}, because they are consequential more
in the calling context than in the called implementation ---
though not entirely, since (for example) throwing exceptions
aborts lexical finalization.  Nor do they belong
precisely to process calculi, because they are most complex
in the context of sequential procedures
(for example, a function cannot catch exceptions thrown from
functions it has spawned in new threads).  As intended here,
\q{Channel Algebra} suggests an intermediate formalization
combining features of both (generalized) lambda and
process calculi (or algebras).  Instead of scattered
\mOldLambda{}-extensions, Channel Algebra unifies multiple
expansions by endowing functions (their signatures,
in the abstract, and function-calls, in the concrete)
with multiple channels, each of which can be
independently modeled by some \mOldLambda{}-extension
(objects, exceptions, captures, and so forth).
}
\p{Specific examples of unorthodox \lambdas{} (objects, exceptions, captures)
suggest a general case: relations or operators between
functions can be modeled as relations between their respective
channels, subject to certain semantic restrictions.  A \i{method} can
be described as a function with several different channels:
say, a \q{\lambdach{}} channel with ordinary arguments (as in
\mOldLambda{}-calculus); a \q{\sigmach{}} channel with a distinguished
\thisc{} carrier (formally studied via \q{\sigmaCalculus{}}); and
a \returnc{} channel representing the return value.  Because the contrast between
these channels is first and foremost \i{semantic} \mdash{} they have different
meanings in the semantics of the programs where they appear \mdash{}
channels may therefore have restrictions governed by
programs' semantics.  For example, as I mentioned
in the context of \q{method chaining}, it may be stipulated that both
\sigmac{} and \returnc{} channels can have at most one carrier;
as a result, a special channel-to-channel operator can be defined which
is specific to passing values between the carriers of
\returnc{} and \sigmac{} channels.  This operator is available because
of the intended semantics of the channel system.
}
\p{In general, a Channel Algebra identifies
several \i{kinds} of channels which
each have their own semantic interpretation, and accompanying axioms
or restrictions.  On the basis of these semantic details,
channel-to-channel operators can be defined, derived from
underlying carrier-to-carrier operators.  A Channel Algebra in this sense
is not a single fixed system, but an outline for modeling
function-call semantics in the context of different programming
languages and environments.
}
\p{As the preceding paragraphs have presupposed, different
functions may have different kinds of channels, which may
or may not be reflected in functions'
types (recall the question, can
two functions have the same type, if
only one may throw an exception)?  This may vary between type
systems; but in any case the contrast between channel \q{structures}
is \i{available} as a criteria for modeling type descriptions.
On this basis, as I will now argue, we can provide type-system
interpretations to source code structures beyond just
values and symbols.
}
\subsubsection{Statements, Blocks, and Control Flow}  
\p{The previous paragraphs discussed expanded channel structures \mdash{} with,  
for example,  
objects and exceptions \mdash{} that model call semantics more complex
than the basic lambda-and-return (of classical \mOldLambda{}-Calculus).  
A variation on this theme, in the
opposite direction, is to \i{simplify} call structures: functions
which lack a return channel have to communicate solely through
side-effects, whose rigorous analysis demands a
\q{type-and-effect} system.  Even further, consider functions
with neither \lambdac{} nor \returnc{} (nor \sigmac{}
nor, maybe, \errorc{}).
As an alternate channel of communication, suppose function bodies
are nested in overarching bodies, and can
\q{capture} carriers scoped to the enclosing function.  \q{Capture semantics}
specifications in \Cpp{} are a useful example, because
\Cpp{} (unlike most languages that support anonymous or \q{intra-expression}
function-definitions) mandates that symbols are explicitly captured
(in a \q{capture clause}), rather than allowing functions to access
surrounding lexically-scoped with no further notation:
this helps visualize the idea that captured symbols are a kind of
\q{input channel} analogous to \lambdac{} and \sigmac{}.
}
\p{I contend this works just as well for code blocks.  Any language which has
blocks can treat them as unnamed function bodies, with
a \q{\capturech{}} channel (but not \lambdach{} or \returnch{}).
When (by language design) blocks \i{can} throw exceptions, it is reasonable to
give them \q{\exceptionch{}} channels (further work, that I put off for
now, is needed for loop-blocks, with \codebreak{} and \codecontinue{}).
Blocks can then be typed as function-like values, under
the convention that function-types can be expressed through
descriptions of their channels (or lack thereof).
}
\p{Consider ordinary source-code
expressions to represent a transfer of values between graph structures: let
\gammaOne{} and \gamaTwo{} be code-graphs compiled
from source at a call site and at the callee's implementation,
respectively.  The function call transfers values from carriers
marked by \gammaOne{} nodes to \gammaTwo{} carriers; with the further
detail of \q{Channel Packages} we can additionally say that the
recipient \gammaTwo{} carriers are situated in a graph structure which
translates to a channel description.  So the morphology of
\gammaOne{} has to be consistent with the channel structure of \gammaTwo{}.
For regular (\q{value}) expressions, we can introduce a new kind
of channel (which I'll call \q{lookup}) acknowledging that the function
called by an expression may itself be evaluated by its own expression, rather
than named as a single symbol (as in a pointer-to-function call
like \pfx{} in \CLng{}).  A segment of source code represents a
value-expression insofar as an equivalent graph representation
comprises a \Gamma{} semantically and morphologically
consistent with the provision of values to channels required
by a function call \mdash{} including the lookup channel on the
basis of which the proper implementation (for overloaded functions)
is selected.  How the graph-structure maps to the appropriate
channels varies by channel kind: for instance the \returnc{} channel
is not passed \i{to} the callee, but rather bound to a carrier as the
right-hand-side of an assignment (an \i{rvalue}) \mdash{} or else passed to
a different function (thus an example of channel-to-channel connection
without an intervening carrier).  A well-formed \Gamma{} represents
part of a function implementation's code graph, specifically that
describing how a Channel Package is concretely provisioned with values
(i.e., a payload).
}
\p{I will use the term \i{call-clause} to designate the
portion of a code graph, and the associated collection of
source code elements, describing a Channel Payload.
Term a call-clause \i{grounded} if its
resulting value is held in a carrier (as
in \yeqfx{}), and \i{transient} if this value
is instead passed on (immediately) to another function
(as in \hfx{}); moreover a call-clause
can be \i{standalone} if it has no result value or this
value is not used; and \i{multiply-grounded} if it has
several grounded result values \mdash{} i.e., a multi-carrier
\returnc{} channel, assuming the type system allows as much.
Grounded and standalone call-clauses can, in turn, model \i{statements};
specifically, \q{assignment} and \q{standalone} statements, respectively.
}
\p{This vocabulary can be useful for interpreting program flow.
Assignment statements with no other side effects can be delayed
until their grounding carrier is convoluted with some other carrier.
Of course, the default choice of \q{eager} or \q{lazy} evaluation is
programming-language-specific, but for abstract discussion
of source code graphs, we have no \i{a priori} idea of temporality;
of a program executing in time.  This is not a matter of
concurrency \mdash{} we have no \i{a priori} idea of procedures running
at the \i{same} time any more than of them running sequentially
(cf. \q{detached}
evaluation as on page \hyperref[detachedeval]{\pageref{detachedeval}}).
\mbox{Any temporal} direction through a graph is an interpretation
\i{of} the graph, and as such it is useful to assume that
graphs in and of themselves assert no temporal
ordering among their nodes or edges.  When modeling
eager-evaluation languages, particular edge-types
can be designated as forcing a temporal order or else
edges can be annotated with additional temporalizing
details.  Without this extra documentation, however,
execution order among graph elements can be
evaluated based on other criteria.
}
\p{In the case of statements, an assignment without side effects
has temporalizing relations only with other statements
using its grounding carrier.  In particular, the order
of statements' runtime need not replicate the order in which
they are written in source code.  For sake of argument, consider
this the default case: Channel Algebras, in principle, model
\q{lazy} evaluation languages, in the absence of any
temporalizing factors.  The actual runtime order among
sibling statements \mdash{} those in the same block \mdash{} then depends,
in the absence of further information, on how their grounding carriers are
used; this in turn works backward from a function's return channel
(in the absence of exceptions or effectual calls).  That is,
runtime order works backward from statements that initialize
carriers in the return channel, then carriers used in those
statements, etc.
}
\p{This order needs to broken, of course, for statements with side-effects.
A case in point is the expansion of \q{\doNotation{}} in Haskell: without
an \i{a priori} temporality, Haskell source code relies on
the asymmetric order of values passed into lambda abstractions to
enforce requirements that effectual expressions evaluate before other
expressions (Haskell does not have \q{statements} per se).  Haskell's
\doH{} \q{blocks} can be modeled (in the techniques used here) as
a series of assignment statements where the grounding carrier of each
statement becomes (i.e., transfers its value to) the sole occupant of a
\lambdach{} channel marking a new function body, which includes all the
following statements (and so on recursively).  There are two concepts
in play here: interpreting any sequence of statements (plus one
terminating expression, which becomes a statement
initializing a \returnch{} carrier) as a function body (not just those covering
the extent of a \q{block}); and interpreting assignment statements
as passing values into \q{hidden} lambda channels.  Of course,
Haskell backs this syntactic convention with monad semantics \mdash{}
the value passed is not the actual value of the monad-typed
carrier but its \q{contained} value.
For sake of discussion, let's call this a \i{monad-subblock} formation.
}
\p{The temporalizing elements in this formation are the \q{hidden lambdas}.
In a multi-channel paradigm, we can therefore consider
\q{monad-subblocks} with respect to other channels.  Consider how
individual statements can be typed: like blocks, statements can
select from symbols in scope and can potentially
result in thrown expressions, so their channel structure is something
like \capturech{} plus \errorch{}.  Even without hidden lambdas, observe
that the runtime order of statements can be fixed in situations
where an earlier statement can affect the value (via non-constant capture)
of a carrier whose value is then used by a later statement.  So
for languages with a more liberal treatment of side-effects than Haskell,
we can interpret chains of statements \i{in fixed order} as successively
capturing (and maybe modifying) symbols which occur in multiple
statements.  Having discussed convoluted \i{carriers},
extend this to channels: in particular, say two \capturech{} channels are
convoluted if there is a modifiable carrier in the first
which is convoluted with a carrier in the second (this is an ordered relation).
One statement must run before a second
if their \capturech{} channels are convoluted, in that order.
}
\p{This is approaching toward a \q{monad-subblock} formation using \i{\capturech{}}
in place of \lambdach{}.  To be sure,
Haskell monad-subblock does have the added gatekeeping
dimension that the symbol occurring after its appearance as grounding
an assignment statement is no longer the symbol with a monad
type, but a different (albeit visually identical) symbol
taken from the monad.  Between two statements, if the
prior is grounded by a monad, the implementation of
its \bind{} function is silently called, with the subsequent
(and all further) statements grouped into a block passed in
to \bindFn{}, which in turn (by design) both extracts its
wrapped value and (if appropriate) calls the passed function.
But this architecture can certainly be emulated on
non-\lambdach{} channels \mdash{} a transform that would belong to the
larger topic of treating blocks as function-values
passed to other functions, to which I now turn.
}
\subsubsection{Code Blocks as Typed Values}
\p{Insofar as blocks can be typed as functions (in a sense),
they may readily be passed around: so loops, \ifthenelse{},
and other control flow structures can plausibly be modeled
as ordinary function calls.  This requires some extra semantic
devices: consider the case of \ifthenelse{} (I'll use this
also to designate code sequences with potential \q{\elseif{}}s), which
has to become an associative array of expressions and functions
with \q{block} type (e.g., with only \capturech{} and \errorch{}
channels).  We need, however, a mechanism to suppress
expression evaluation.  Recall that expressions are
concretized channel-structures which include a
\lookupch{} channel providing the actual implementation to call.
Assume that \lookupch{} can be assigned a flag which suppresses
evaluation.  Assume also that carriers can be declared which
hold (or somehow point to) \i{expressions} that evaluate
to typed values,
in lieu of holding these values directly
(note that this is by intent orthogonal
to a type system: the point is not that carriers can hold values
whose type is designed to encapsulate potential
computations yielding another type, like \stdFuture{} in \Cpp{}).
Consider again the nested-expression variant of
\carrOneOpTransferTwo{}: when the result
of one function call becomes a parameter to another
function, the value in the former's
\returnch{} carrier (assume there is just one) gets
transferred to a carrier in the latter's \lambdach{} channel
(or \sigmach{}, say).  This handoff can be described before being
effectuated: a language runtime is free to vary the
order of expression-evaluation no less than of statements.
The semantics of a carrier-transfer between \Twof{}'s return
and \Onef{}'s lambda does not stipulate
that \Twof{} has to \i{run} before \Onef{}; language
engines can provide semantics for \carrOneOpTransferTwo{} allowing \carOne{} to hold a delayed capability to
evaluate the \Twof{} expression.  Insofar as this is an
option, functions can be given a signature \mdash{} this
would be included in the relevant \TXL{} \mdash{} where
some carriers are of this \q{delayed} kind.
Functions like \ifthenelse{} can then be declared in
terms of these carriers.
}
\p{To properly capture \ifthenelse{} semantics, it is also
appropriate to notate conditions on how blocks
passed as function values are used.  In the
case of \ifthenelse{}, the implementation
will receive multiple function values with the caller
expecting \i{exactly one} function to be called.
This constitutes a requirement on how carriers are
used, analogous to mandating that a \returnch{}
(or either \returnch{} or \errorch{}) carrier be initialized,
or that mutable references passed
to a function for initialization (as an alternative
to returning multiple values) are initialized.
While such requirements can potentially be
described as annotations on carriers/channels,
this is a general issue outside the scope of
blocks-as-function-values.
}
\p{A thorough treatment of blocks-as-functions also needs to
consider standard procedural affordances like \codebreak{}
and \codecontinue{} statements.  Since blocks
can be nested, some languages allow inner blocks to express the
codewriter's intention to \q{break out of} an outer block from
an inner block.  One way to model this via Channel Algebra is
to introduce a special kind of return channel for blocks
(called a \q{\breakch{}}, perhaps) which, when it has
an initialized carrier, uses this channel to hold a value that the
enclosing block interprets in turn: by examining the
inner \breakch{} the immediately outer block can decide whether it
itself needs to \q{break} and, if so, whether its own
\breakch{} channel needs to have an initialized carrier.  The
presence of such a \breakch{} can type-theoretically
distinguish loop blocks from blocks in (say)
\ifthenelse{} contexts.
}
\p{Further discussion of code models via Channel Complexes and
Channel Algebras is outside the scope of this chapter, but
is demonstrated in greater detail in the accompanying code-set.
Hopefully, the best way to present Channel Semantics outside
the basic \lambdach{}/\sigmach{}/\returnch{}/\exceptionch{}
quartet is via demonstrations in live code.
In that spirit, the demo code-set focuses more
on practical engineering and problem-solving where Channel models
can be useful, and I'll briefly review its structure and its
organizing rationales in the Conclusion.
}

\section{Conclusion}
\p{The upshot of my argument lat section was 
that the realm of mereologuical theory should 
be divided into \i{cognitive} and \i{multiscientific}.  
Criteria of mereological non-trivialoty are 
derived in the first case from \i{cognitive frames}, 
and in the second case from an \q{$R$-patchwork}: 
a patchwork of scientific theories delimiting the 
extent of different scientific registers and the 
dynamics and autopoesis sustaining the 
integrtity of (in various domains) 
non-trivial wholes.  Moreover, there is a structural 
resonance and interpenetration between 
cognitive frames and the \q{$R$-patchwork}.
}
\p{Whether from the perspective of cognitive frames or the 
\q{extramental} perspective of the $R$-patchwork, 
mreological theories can be rasonably elaborated 
in parallel to theories of the cognitive/ontological 
status of wholes participating in \xhppy{} style 
parthood assertions.  Any formalization of 
parthood relations can assume that 
}
\p{}

\subsection{Types, Sets, and Concepts}
\p{In formal/computational contexts, types can be defined as sets of both values and
\q{expectations} \cite{MathieuBouchard} (meaning assumptions which may be made about
all values covered by the type); alternatively, we can (perhaps better) consider types as
\i{spaces} of values.  Types' extensions have internal structure; there
can be \q{null} or \q{invalid} values, default-constructed values, and
so forth, which are \q{regions} of the conceptual space spanned or 
encompassing types.\footnote{Conceptual Space theory is outside the scope of this paper, but I'll note that 
it suggests a promising link between natural linguistics and formal/computational 
type theory, as suggested by computational or scientific expositions of 
the original Conceptual Space account developed by Peter \Gardenfors{}: 
cf. \cite{Zenker}, \cite{RaubalAdams} \cite{RaubalAdamsCSML}, 
\cite{Strle}.  Meanwhile, projects to develop formal models for 
Cognitive Grammar have also adopted Conceptual Space theory as an 
underlying semantics: \cite{KennethHolmqvist}, \cite{HolmqvistDiss}, 
\cite{MattSelway}, \cite{InteractingConceptualSpaces}.    
}
There is definitional interdependence
between types and functions: a function is defined in terms of the types it accepts as parameters and
returns \mdash{} rather than its entire set of possible inputs and outputs, which can
vary across computing environments.
\footnote{Moreover, expectations in a particular case
may be more precise than what is implied by the type itself \mdash{} it is erroneous
to assume that a proper type system will allow a correct \q{set of values} to
be stipulated for each point in a computation (the kind of contract enforced via
by documentation and unit testing).  So state-space in a given context may include many
\q{unreasonable} values, implying that within the overall space there is a \q{reasonable}
subspace, except that this subspace may not be crisply defined.
}  These are some reasons why in theoretical
Computer Science types are not \q{reduced} to underlying sets; instead, extensions
are sometimes complex spaces that model states of, or internal organization of comparisons
among, type instances.
}
\p{An obvious paradigm is organizing type-extensions around prototype/borderline
cases \mdash{} there are instances which are clear examples of types and
ones whose classification is dubious.  I contend, however, 
that common resemblance is not always a good marker
for types being well-conceived \mdash{} many useful concepts are common
precisely because they cover many cases, which makes defining
\q{prototypes} or \q{common properties} misleading.  
Also, sometimes the clearest
\q{representative} example of a type or concept is actually not a
\i{typical} example: a sample latter or model home is actually not (in
many cases) a real letter or home.  So resemblance-to-prototype is at
best one kind of \q{inner organization} of concepts' and types' spaces
of extension.  
}
\p{Sets, concepts, and types represent three different primordial thought-vehicles for
grounding notions of logic and meaning.  To organize systems around \i{sets} is
to forefront notions of inclusion, exclusion, extension, and intersection,
which are also formally essential to mathematical logic and undergird the
classical interdependence of sets, logic, and mathematics.
To organize systems around \i{concepts} is to forefront practical engagement
and how we mold conceptual profiles, as collections of ideas and pragmas,
to empirical situations.  To organize systems around \i{types} is to forefront
\q{functions} or transformations which operate on typed values, the interrelationships
between different types (like subtypes and inclusion \mdash{} a type can itself
encompass multiple values of other types), and the conceptual abstraction
of types themselves from the actual sets of values they may exhibit
in different environments.  Sets and types are
formal, abstract phenomena; whereas concepts are characterized by
gradations of applicability, and play flexible roles in thought and language.
The cognitive role of concepts can be discussed with some rigor, but there is a
complex interplay of cognitive schema and practical engagements which
would have to be meticulously sketched in many real-world scenarios, if
our goal were to translate conceptual reasoning to formal structures
on a case-by-case basis.  We can, however, consider in general
terms how type-theoretic semantics can capture conceptual structures
as part of the overall transitioning of thoughts to langauge.
}
\p{A concept does not merely package up a definition, like \q{restaurant} as
\q{a place to order food}; instead concepts link up with other concepts
as tools for describing and participating in situations.  Concepts are
associated with \q{scripts} of discourse and action, and find their
range of application through a variegated pragmatic scope.
We should be careful not to overlook these pragmatics, and
assume that conceptual structures can be simplistically
translated to formal models.
Cognitive Linguistics critiques
Set-Theoretic or Modal Logic reductionism (where a concept is just a set
of instances, or an extension across different possible worlds) \mdash{} George Lakoff and Mark Johnson,
prominently, argue for concepts' organization around
prototypes (\cite[p. 18]{LakoffJohnson}; \cite[p. 171, or p. \textit{xi}]{Johnson})
and embodied/enactive patterns of interaction (\cite[p. 90]{LakoffJohnson};
\cite[p. 208]{Johnson}).
Types, by contrast, at least in linguistic applications of type theory, are abstractions
defined in large part by quasi-functional notions of phrase struture.
Nevertheless, the \i{patterns} of how types may inter-relate
(mass-noun or count-noun, sentient or non-sentient, and so forth)
provide an infrastructure for conceptual understandings to be
encoded in language \mdash{} specifically, to be signaled by which typed
articulations conversants choose to use.  A concept like
\i{restaurant} enters language with a collection of understood
qualities (social phenomena, with some notion of spatial location and
being a \q{place}, etc.) that in turn can be marshaled by sets of
allowed or disallowed phrasal combinations, whose parameters
can be given type-like descriptions.  Types, in this sense,
are not direct expressions of concepts but vehicles for
introducing concepts into language.
}
\p{Concepts (and types also) are not cognitively the same as their
extension \mdash{} the concept \i{restaurant}, I believe, is distinct from
concepts like \i{all restaurants} or \i{the set of all restaurants}.
This is for several reasons.  First, concepts can be pairwise different
not only through their instances, but because they highlight different
sets of attributes or indicators.  The concepts \q{American President} and \q{Commander in Chief}
refer to the same person, but the latter foregrounds a military role.
Formal Concept Analysis considers \i{extensions} and \q{properties}
\mdash{} suggestive indicators that inhere in each
instance \mdash{} as jointly (and co-dependently) determinate: concepts
are formally a synthesis of instance-sets and property-sets \cite{YiyuYao},
\cite{Belohlavek}, \cite{Wille}.  Second,
in language, clear evidence for the contrast between \i{intension} and
\i{extension} comes from phrase structure: certain constructions specifically
refer to concept-extension, triggering a mental shift from thinking of the
concept as a schema or prototype to thinking of its extension (maybe in some context).
Compare:
\begin{sentenceList}\sentenceItem{} \label{itm:rhinor} Rhinos in that park are threatened by poachers.
\sentenceItem{} Young rhinos are threatened by poachers.
\end{sentenceList}
Both sentences focus a conceptual lens in greater detail than \i{rhino} in general, but
the second does so more intensionally, by adding an extra indicative criterion; while
the former does so extensionally, using a phrase-structure designed to operate on
and narrow our mental construal of \q{the set of all rhinos}, in the sense of
\i{existing} rhinos, their physical place and habitat, as opposed to
the \q{abstract} (or \q{universal}) type.  So there is a familiar semantic
pattern which mentally transitions from a lexical type to its extension and
then extension-narrowing \mdash{} an interpretation that, if accepted, clearly
shows a different mental role for concepts of concepts' \i{extension} than the
concepts themselves.\footnote{There is a type-theoretic correspondence between intension and
extension \mdash{} for a type \Tnoindex{} there is a corresponding \q{higher-order} type
of \i{sets} whose members are \Tnoindex{} 
(related constructions are the type of \i{ordered sequences} of \Tnoindex{};
unordered collections of \Tnoindex{} allowing repetition; and stacks, queues, and
deques \mdash{} double-ended queues \mdash{} as \Tnoindex{}-lists that can grow or shrink
at their beginning and/or end).  If we take this (higher-order)
type gloss seriously, the extension of a concept is not its \i{meaning}, but a
different, albeit interrelated concept.  Extension is not definition.
\i{Rhino} does not mean \i{all rhinos} (or \i{all possible rhinos}) \mdash{} though arguably
there are concepts \i{all rhinos} and \i{all restaurants} (etc.) along with the concepts
\i{rhino} and \i{restaurant}.
}  
}
\p{Concepts, in short, do not mentally signify sets, or
extensions, or sets-of-shared-properties.  Concepts, rather, are cognitive/dialogic tools.
Each concept-choice, as presentation device,
invites its own follow-up.  \i{Restaurant} or \i{house} have meaning not via
idealized mental pictures, or proto-schema, but via kinds of things
we do (eat, live), of conversations we have, of qualities we deem relevant.  Concepts do not
have to paint a complete picture, because we use them as part of ongoing situations
\mdash{} in language, ongoing conversations.  Narrow concepts \mdash{} which may best exemplify
\q{logical} models of concepts as resemblance-spaces or as rigid designators to
natural kinds \mdash{} have, in practice, fewer use-cases \i{because} there
are fewer chances for elaboration.  Very broad concepts, on the other hand, can have,
in context, too \i{little} built-in \i{a priori} detail.
(We say \q{restaurant} more often than \i{eatery}, and
more often than \i{diner}, \i{steakhouse}, or \i{taqueria}).  Concepts dynamically play
against each other, making \q{spaces} where different niches of meaning, including
levels of precision, converge as site for one or another.  Speakers need freedom to choose
finer or coarser grain, so concepts are profligate, but the most oft-used trend toward middle
ground, neither too narrow nor too broad.  \i{Restaurant} or \i{house} are useful because they are noncommittal, inviting more detail.
These dynamics govern the flow of inter-concept relations (disjointness, subtypes, partonymy, etc.).
}
\p{Concepts are not rigid formulae (like instance-sets or even attributes fixing when
they apply); they are mental gadgets to initiate and
guide dialog.  Importantly, this
contradicts the idea that concepts are unified around instances' similarity (to each other or
to some hypothetical prototype): concepts have avenues for contrasting
different examples, invoking a \q{script} for further elaboration, or for building temporary filters.  
In, say, 
\begin{sentenceList}\sentenceItem{} Let's find a restaurant that's family-friendly.
\end{sentenceList}
allowing such one-off narrowing is a feature
of the concept's flexibility.
}
\p{In essence: no less important, than acknowledged similarities across all instances, are well-rehearsed ways
\visavis{} each concept to narrow scope by marshaling lines of \i{contrast}, of \i{dissimilarity}.
A \i{house} is obviously different from a \i{skyscraper}
or a \i{tent}, and better resembles other houses; but there are also more nontrivial \i{comparisons}
between houses, than between a house and a skyscraper
or a tent.  Concepts are not only spaces of similarity, but of \i{meaningful kinds of differences}.
}
\p{To this account of conceptual breadth we can add the conceptual matrix spanned by
various (maybe overlapping) word-senses: to \i{fly}, for example, names
not a single concept, but a family of concepts all related to airborn
travel.  Variations highlight different features: the path of flight (\i{fly to Korea}, \i{fly over the mountain});
the means (\i{fly Korean air}, \i{that model flew during World War II});
the cause (\i{sent flying (by an explosion)}, \i{the bird flew away (after a loud noise)},
\i{leaves flying in the wind}).  Words allow different use-contexts
to the degree that their various \i{senses} offer an inventory of aspects for
highlighting by \i{morphosyntactic} convention.  Someone who says \i{I hate to fly} is not
heard to dislike hand-gliding or jumping off mountains.\footnote{People, unlike birds, do not fly \mdash{} so the verb, used intransitively
(not flying \i{to} somewhere in particular or \i{in} something in particular),
is understood to refer less to the physical motion and more to the socially
sanctioned phenomenon of buying a seat on a scheduled flight on an airplane. The construction
highlights the procedural and commercial dimension, not the physical mechanism and
spatial path.  But it does so \i{because} we know human flight is
unnatural: we can poetically describe how the sky is filled with flying leaves or birds,
but not \q{flying people}, even if we are nearby an airport.
}  Accordant variations
of cognitive construal (attending more to mode of action, or path, or motives, etc.),
which are elsewhere signaled by grammatic choices, are also spanned by a conceptual
space innate to a given word: senses are finer-grained meanings availing themselves to one construal or another.
}
\p{So situational construals can be signaled by word- and/or
syntactic form choice (locative, benefactive, direct and indirect
object constructions, and so forth).  Whereas conceptual organization
often functions by establishing classifications, and/or invoking
\q{scripts} of dialogic elaboration, cognitive structure tends to apply more
to our attention focusing on particular objects, sets of objects, events, or
aspects of events or situations.  
So the contrast between singular, mass-multiples, and count-multiples,
among nouns, depends on cognitive
construal of the behavior of the referent in question (if singular, its
propensity to act or be conceived as an integral whole; if multiple, its
disposition to either be divisible into discrete units, or not).
Or, events can be construed in terms of their causes
(their conditions at the outset), or their goals (their conditions at
the conclusion), or their means (their conditions in the interim).
Compare \i{attaching} something to a wall (means-focused) to
\i{hanging} something on a wall (ends-focused); \i{baking} a cake
(cause-focus: putting a cake in the oven with deliberate intent to cook it)
to \i{burning} a cake (accidentally overcooking it).\footnote{We can express
an intent to bake someone a cake, but not (well, maybe comedically) to
\i{burn} someone a cake (\q{burn}, at least in this context, implies
something not intended); however, we \i{can} say
\q{I burnt your cake}, while it is a little jarring to say
\q{I baked your cake} \mdash{} the possessive implies that some
specific cake is being talked about, and there is less apparent reason
to focus on one particular stage of its preparation (the baking) once
it is done.  I \i{will} bake a cake, in the future, uses
\q{bake} to mean also other steps in preparation (like \q{make}), while,
in the present, \q{the cake \i{is} baking} emphasizes more its
actual time in the oven.  I \i{baked your cake} seems to focus
(rather unexpectedly) on this specific stage even after it is completed,
whereas \i{I baked you a cake}, which is worded as if the recipient
did not know about the cake ahead of time, apparently uses \q{bake} in
the broader sense of \q{made}, not just \q{cooked in an oven}.
Words' senses mutate in relation to the kinds of situations where they are used
\mdash{} why else would \i{bake} mean \q{make}/\q{prepare} in the past or future tense but
\q{cook}/\q{heat} in the present?
}
These variations are not random assortments of polysemous words' senses:
they are, instead, rather predictably distributed according
to speakers' context-specific knowledge and motives.
}
\p{I claim therefore that \i{concepts} enter language complexly, influenced by
conceptual \i{spaces} and multi-dimensional semantic and syntactic selection-spaces.
Concepts are not simplistically \q{encoded} by types, as if for
each concept there is a linguistic or lexical type that just
disquotationally references it \mdash{} that the type \q{rhino} means the concept
\i{rhino} (\q{type} in the sense that type-theoretic semantics would model lexical
data according to type-theoretic rules, such as \i{rhino} as subtype of \i{animal} or
\i{living thing}).
Cognitive schema, at least in the terms I just laid out, select particularly
important gestalt principles (force dynamics, spatial frames, action-intention)
and isolate these from a conceptual matrix.  On this basis, we can argue that
these schemata form a precondition for concept-to-type association; or,
in the opposite logical direction, that language users' choices to employ
particular type articulations follow forth from their prelinguistic
cognizing of practical scenarios as this emerges out of collections
of concepts used to form a basic understanding of and self-positioning within them.
}
\p{In this sense I called types \q{vehicles} for concepts: not that types \i{denote}
concepts but that they (metaphorically) \q{carry} concepts into language.  
\q{Carrying} is enabled by types' semi-formal rule-bound
interactions with other types, which are positioned to capture concepts' variations and
relations with other concepts.
}
\p{To express a noun in the benefactive case, for example, which can be seen as attributing to
it a linguistic type consistent with being the target of a benefactive,
is to capture the concept in a type-theoretic gloss.
It tells us, I'm thinking about this thing in such a way that it
\i{can} take a benefactive (the type formalism attempting to capture
that \q{such a way}).
A concept-to-type \q{map}, as I just
suggested, is mediated (in experience and practical reasoning) by
cognitive organizations; when (social, embodied) enactions take
linguistic form, these organizing principles can be encoded in how
speakers apply morphosyntactic rules.  
}
\p{So the linguistic structures,
which I propose can be formally modeled by a kind of type theory, work
communicatively as carriers and thereby signifiers of cognitive
attitudes. The type is a vehicle for the concept because it takes part in constructions
which express conceptual details \mdash{} the details
don't emerge merely by virtue of the type itself.
I am not arguing for a neat concept-to-type correspondence; instead, a type system provides a
\q{formal substrate} that models (with some abstraction and simplification) how
properties of individual concepts translate
(via cognitive-schematic intermediaries) to their
manifestation in both semantics and syntax.
}
\p{Continuing with declention as a case study,
consider how an \q{ontology} of word senses 
can interrelate with the benefactive.
A noun as a benefactive target most often is a person or some other
sentient/animate being; an inanimate benefactive is most likely
something artificial and constructed (cf., \i{I got the car new tires}).
How readily hearers accept a sentence \mdash{} and the path they
take to construing its meaning so as to make it grammatically acceptable
\mdash{} involves interlocking morphological and type-related considerations;
in the current example, the mixture of benefactive case and which noun
\q{type} (assuming a basic division of nouns into e.g.
animate/constructed/natural) forces a broader or narrower
interpretation.  A benefactive with an \q{artifact} noun, for example,
almost forces the thing to be heard as somehow disrepaired:
\begin{sentenceList}\sentenceItem{} I got glue for your daughter.
\sentenceItem{} I got glue for your coffee mug.
\end{sentenceList}
We gather (in the second case) that the mug is broken \mdash{} but this is never spelled out
by any lexical choice; it is implied indirectly by using benefactive case.  
It is easy to design similar examples with other cases:
a locative construction rarely targets \q{sentient} nouns, so in
\begin{sentenceList}\sentenceItem{} We're going to Grandma!
\sentenceItem{} Let's go to him right now.
\sentenceItem{} Let's go to the lawyers.
\sentenceItem{} Let's go to the press.
\end{sentenceList}
we mentally substitute the person with the place where they live or work.
}
\p{Morphosyntactic
considerations are also at play: \i{to the lawyers} makes \q{go} sound more like \q{consult with},
partly because of the definite article (\i{the} lawyers implies conversants have some prior involvement
with specific lawyers or else are using the phrase metonymically, as in \q{go to court} or
\q{to the courts}, for
legal institutions generally; either reading draws attention away from literal spatial implications of
\q{go}). \i{Go to him} implies that \q{he} needs
some kind of help, because if the speaker just meant going to wherever he's at, she probably would
have said that instead.
}
\p{Similarly, the locative in \i{to the press} forces the mind to
reconfigure the landmark/trajector structure, where \i{going} is thought not as a literal
spatial path and \i{press} not a literal destination \mdash{} in other words, the phrase must be
read as a metaphor.  But the \q{metaphor} here is not \q{idiomatic} or removed from linguistic rules
(based on mental resemblance, not language structure); here it
clearly works off of formal language patterns: the landmark/trajector
relation is read abstracted from literal spatial movement because the locative is applied
to an expression (\i{the press}) which does not (simplistically) meet
the expected interpretation as \q{designation of place}.
}
\p{In short, there are two different levels of \i{granularity} where we 
can look for agreement requirements: a more fine-grained level where e.g. 
\i{locative} draws in a type-specification of a \i{place} or \i{location}; 
and a coarser level oriented toward Parts of Speech, and typologies 
of phrasal units.  On the finer scale, what linguistics can draw from type 
theory gravitates toward type-coercions, \q{dependent types}, and 
topics based on programming language type systems, like \q{monads} 
\mdash{} \cite{ZhaohuiLuo}, \cite{LuoSoloviev},
\cite{ChatzikyriakidisLuo},  
\cite{AsherPustejovsky}, 
\cite{BarkerShan}, \cite{ShanMonads}, \cite{ShanThesis}.  
On the coarser scale, on the other hand, 
analyses can focus on the interconnections between types and functions: 
most top-level linguistic types fit the formal model of 
(in computational contexts) \q{functional} types, associated with 
types of \q{inputs} and \q{outputs}.  For instance, 
assuming we have a primordial \i{noun} type and a \i{proposition} type 
(the type assigned to complete sentences), a \i{verb} is then at some 
abstract reading a \q{function} from nouns to propositions \mdash{} insofar 
as verbs produce propositions when combined with nouns.  Similarly 
an adjective maps nouns to other nouns (in a conceptual sense; noun-phrases, 
literally speaking); adverbs map verbs to other verbs, and so forth.              
}
\p{The proposition type (say, \q{\Prop{}}) provides a
type attribution for sentences, but also for sentence parts: \i{he is at school},
for example, presents a complete idea, either as its own sentence or part of a
larger one.  In the latter case, a \Prop{} phrase would typically be preceded with a
word like \i{that}; \i{syntactically}, \q{that} is essentially a 
connector, helping sentence-parts link with each other:
\begin{sentenceList}\sentenceItem{} I think he is at school.
\sentenceItem{} \label{itm:bt} I believe that he is at school.
\end{sentenceList}
Type-theoretically, however, we may want to assign types to every word, even those
\mdash{} like \i{that} in (\ref{itm:bt}) 
which seem auxiliary and lacking much or any semantic content of their own.
Arguably, \i{that} serves to \q{package} an assertion, encapsulating
a proposition as a presumed fact designated as
one idea, for the sake of making further comments, as if \q{making a noun} out
of it: \PropToN{}.  Perhaps our intuitions are more as if \i{that he is at school}
is also a proposition, maybe a subtly different kind, by analogy to how
questions and commands are also potentially \Prop{} variants.  Since \thatPhrases{} are \q{arguments} for verbs,
the choice then becomes whether it is useful to expand our type picture of verbs
so that they may act on propositions as well as nouns,
or rather type \q{encapsulated} propositions as just nouns
(maybe special kinds of nouns).
}
\p{In either case, \i{I know that ...} clearly involves a verb with subject and direct
object: so either \VisNNtoProp{} or \VisNProptoProp{}.  
Using \i{link grammar}, 
Consider the role of a \TS{}-link here:
specifically, \TS{} connects the verb to the assertorial direct object (most
directly, to \i{that}).\footnote{Link Grammar models parses via interword relations classified according to 
some four dozen recognized syntactic and semantic connectors; 
a complete parse yields a labeled \q{graph} of each sentence, similar to 
graphs derived from Depenency Grammars, with link-kinds notated via 
labels like \q{\TS{}} (the exact 
vocabulary of link kinds establishes specific grammars for specific 
natural languages) \mdash{} cf. \cite{SleatorTamperley};
\cite{Debusmann} 
\cite{DebusmannThesis}.  
\cite{DebusmannDuchierRossberg}, 
\cite{Nivre}.
Link Grammar is distinguished from Depenency Grammar in that the former 
does not directly model head/dependent relations within word-pairs, but 
it \i{does} develop a theory of individual words carrying incompleteness 
or expectations which define \q{compatibility} between words, allowing 
links to be formed; the link/dependency comparison is investigated in 
\cite{HongEisner}, ; and other writings compare both grammatic schools to 
more hierarchical phrase-structure grammars: 
\cite{KongRushSmith}, \cite{Schneider}, \cite{XiaPalmer}. 
}  The purely formal consideration is ensuring that
types are consistent: either the \TS{} target is \Prop{}, 
with the verb type modified accordingly; or the \TS{} target is a noun,
though here it is fair to narrow scope.  For this particular kind of
link, the target must express a proposition: either typed directly as
such or typed as, say, a noun \q{packaging} a proposition, which would then
be a higher-order type relation (just as \q{redness} is a noun \q{packaging}
an adjective, or \q{running} is an adjective packaging a verb).  In other words,
it is difficult to state the type restrictions on the link-pair without employing
more complex or higher-order type formations.
}
\p{On the other hand, this is
another example of the fuzzy boundary between syntax and semantics: given a sentence
which seems to link a verb calling for a belief or assertion (like \q{know},
\q{think}, \q{suggest}, \q{to be glad}) to something that is not proposition-like, is such a
configuration ungrammatical, or just hard to understand?  Clearly, the
\i{semantic} norms around verbs like \q{know} is that their \i{subject}
has some quality of sentience (or can be meaningfully attributed
belief-states, even if speakers know not to take it literally: \q{The function
doesn't know that this number will never be zero}); and their \i{object} should
be somehow propositional.  But applying type theory (or type theory in conjunction
with Dependency Grammar) leaves open various analytic preferences: these
requirements can be presented as rigid grammatic rules or as \q{post-parsing}
semantic regulations.  How to model the qualities of sentience (or at least of having
propositional attitudes broadly conceived), for the noun, and of propositionality,
for the direct object, are again at the discretion
of the analysis (subtypes, quality-associations, or etc.) \mdash{} Figure ~\ref{fig:Iknow} shows one potential,
rather simplified unpacking of the sentence; from this structure details
can be added perhaps as extra syntax constraints or perhaps more as cues
to interpretation.\input{figure.tex}  If these
requirements are seen as more syntactic, so qualities are incorporated into
data like Part of Speech (say, a noun designating something with propositional attitudes
being a subtype of of a generic \N{} type), then we are more likely to analyze
violations as simply incorrect (recall \q{The tree wants to run away from the dog}
\mdash{} ungrammatical or just somehow \q{exotic}?).
Some examples suggest less incorrectness as clever or poetic usage \mdash{} so a
richer analysis may recognize expressions as type- and link-wise acceptable, but
showing incongruities (which is not the same as impropriety)
at a more fine-grained type level.  That \i{to want} takes a subject
\i{associated} with sentience does not force type annotations to inscribe this
in grammatic or lexical laws; instead, these associations can be
introduced as potential \q{side effects}, \i{triggering} re-associations
such as forcing hearers to ascribe sentience to something (like a tree) where
such ascription is not instinctive.  The type effect in this case lies more
at the conceptual level, the language-user sifting conceptual
backgrounds to find a configuration proper to the type requirements (in what
sense can a tree \q{want} something?).  In this \q{tree} case we probably
appeal to concepts of \q{as if}: if the tree \i{were} sentient, it
would be nervous of the dog sniffing around \mdash{} a humorous way of calling
attention to the dog's actions (obliquely maybe alluding to people's background
knowledge that dogs sometimes do things, like pee, in inconvenient
places, from humans' perspectives).
}
\p{In brief, it is certainly possible \mdash{} though by no means mandatory \mdash{} to model
type requirements with greater flexibility at a provisional grammatical layer,
and then narrow in on subtypes or extra accumulations of qualifications on
type-instances in a transition from grammar to semantics.  Perhaps cognitive
schema occupy an intermediary role: progressing from basic recognition of
grammaticality, through cognitive schema, to conceptual framing, with type
machinery capturing some of the thought-processes at each \q{step}
(not that such \q{steps} are necessarily in a temporal sequence).  The basic
verb-subject-direct object articulation sets up an underlying cognitive
attitude (represented by a basic type-framing of verb, noun, and proposition,
like the \VisNNtoProp{} signature).  Cognitive ascriptions fill this out
by adding detail to the broader-hewed typing, associating sentience with the
subject and propositionality with the object (sub- or higher-order typing
modeling this stage).  And how the actual lexical choices fit these cognitive
expectations \mdash{} I call them cognitive because they are intrinsically tied
to structurational schema in the type, morphology, and word-order givens
in the encountered language \mdash{} compels conversants to dip into background
beliefs, finding concepts for the signified meanings that hew to the
intermediary cognitive manipulations (finding ways to conceptualize the
subject as sentient, for example).  This also has a potential type model,
perhaps as forcing a type conversion from a lexical element which does
not ordinarily fit the required framing (such as giving inanimate things
some fashion of sentience).  Type theory
can give a window onto unfolding intellection at these multiple stages,
although we need not conclude that the mind subconsciously doing this thinking
mimics a computer that churns through type transformations mechanically and exactly.
}
\p{I envision the unfolding that I have just sketched out as something Phenomenological
\mdash{} it arises from a unified and subjective consciousness, one marked by
embodied personal identity and social situation.  If there are structural stases
that can be found in this temporality of experience, these are not constitutive
of conscious reality but a mesh of rationality that supports it, like the veins in
a leaf.  Stuctural configurations can be lifted from language insofar as it is a
conscious, formally governed activity, and lifted from the ambient situations which
lend language context and meaning intents.  So any analytic emphasis on
structural fixpoints threaded through the lived temporality of consciousness is an
abstraction, but one that is deliberate and necessary if we want to make scientific
or in any other manner disputatable claims about how language and congition works.
In that spirit, then, I will try to condense the three \q{layers} of unfolding
understanding, which as I have sketched them are posited in the metaphysical order
of temporal experience \mdash{} \q{unfolding} in likely overlapping, blending ways
\mdash{} I will \q{read into} them a more static and logically stacked meta-structure.
Where I have sketched three layers or stages of unfolding language understanding,
I will transition to proposing three \q{tiers} of language organization, in particular
three levels where type-theoretic models can be applied.
}

\section{Phenomenology, AI, and Holism}
\p{This situation is not irrelevant to either linguistics or to Cognitive
Phenomenology.  In Computational Linguistics,
for example, linguistic IR models seem
to be valued based on their utility in AI-driven
Natural Language Processing.  This presents a disciplinary
bifurcation, where potential computational models
are \i{either} connoted rather informally as part of
a thoretical investigation among linguists
(or philosophers of language, etc.), \i{or}
concretely implemented, in some kind of software package,
but then measured as components in a Natural Language Processing
system: assessed on the basis of how the system overall
approximates human language understanding via artifical means.
Another genre of formal models,
such as the type- or monadic theories I have alluded to here,
may also have potential software incarnations but
tend to be developed instead in a mathematical
style, effectively \q{programming} in the abstract
space of theorems and syllogisms rather than actual
compuers.  Each of these methodologies skirt
around the potential intermediary tie:
concrete computational systems that are
designed as exemplifications of semantic, grammmatical,
or pragmatic theories, presented as hands-on
software to anchor theoretical discussion but also intended
as tools to advance the human study of language, rather
than as steps toward synthetic avatars.
}
\p{At the same time, there is another side to the story: software
implementations offer a focal center for research, something
tangible that scholars can experiment and collaborate on.
The AI story provides a target goal; it helps developers understand
the local, technical code they are working on by
connecting it to a larger system.
Whatever philosophical objection one may have to AI
initiatives, we should recognize the value of
expanding academic and institutional practice beyond
just writing and reading research papers.  Insofar as
part of one's scholarly modus operandi can include
writing computer code, and studying code
repositories developed by others, we can benefit from a
hands-on, even trial-and-error kind of experimentation.
}
\p{In effect: software which can be given concrete tasks \mdash{} if it
does \i{this} properly (whatever \i{this} is),
then there is some larger thoretical point that
is demonstrated \mdash{} and then, incrementally, evolves
to realize those tasks, provides a distinct form
of intellectual engagement.  We get \i{that} to work, then
\i{that}, then fix \i{that}.  This kind of
\q{code-and-fix} cycle is quicker than conventional
research, especially in the humanities, where the
routines of authorship and publication and conferences can
feel like they are unfolding in slow motion.
}
\p{Perhaps for this reason, some of the most interesting
cognitive models hacve come from
computational and academic environemnts informed by
ambitious \q{Artificial General Intelligence}
programmes, like Carnegie Melon University's
OpenCog, and the \q{lmnTal} project at Waseda University
in Tokyo.  These projects both employ
formal-semantically expressive, hypergraph-oriented
data systems that embody both the structural
and procedural dimensions of computer systems \mdash{}
manifesting theories of both the execution
of computational processes and the representation of
formalized information.  These are important
models even outside the Artificial General Intelligence
ideology.
}
\p{In fact, these are models which in linguistics
and phenomenology may deserve more attention
than Artificial General Intelligence \i{qua}
philosophy.  But we should not discredit the
role that Artificial General Intelligence may
provide as a kind of intellectual compass helping
scholars and engineers reason through the
intrestitial machinery which may in fact be
more real than the philosophical vision, but
also less effective as theoretical \i{vie ferrate}.
Metaphors can triangulate research
whereas analogies guide transfers of theories
or methods between fields \mdash{} i.e., analogies are
more trustworthy landmarks than metaphors
for surveying the envisioned future of a science
\mdash{} but metaphors can still be intuitive guides;
maybe AI and Artificial General Intelligence can
stabilize into our overall science and
philosophy as a modest but suggestive metaphor.
}
\p{}
\p{}

\spsubsectiontwoline{Phenomenology and the Limits of Formalization}
\p{Considering how Analytic and Continental Philosophy 
has diverged, it is easy to forget that the early 
mature works of Husserl \mdash{} arguably the root 
stock of 20th-century \q{Continental} thought \mdash{} 
emerged from the same circle of logico-mathematical 
questions and evolutions as the ideas 
of Frege, Russerl, Wittgenstein, et. al. that 
would evolve into \q{Analytic Philosophy}.  
Nor is it necessarily true that analytic philosophers 
have remained closer in spirit to the unfolding 
horizons of mathematical knowledge in
the 20th century.  
}
\p{The perpetuation of 20th-century philosophical tropes
\mdash{} Analytic vs. Continental Philosophy; humanities vs.
logicao-mathematics \mdash{} perhaps diagnoss something incomplete in
20th century thought that the 21st century needs to transcend, but
it also bears witness to intellectual structures that directly
abut phenomenology itself.  The idea, for example, that
logico-mathematical analysis is distinct from cultural-humanities
theory \mdash{} that works of art, for example, do not have logical
or mathematical forms; that the structures evinced by cultural
and artistic productions (and any system of signs and communication
necessarily has \i{some} rigorous structure) does not itself
have a logical and mathematical dimension, or a structural
form that if specified in isolation from its signifying context
would reveal logical and/or mathematical properties.  A parallel
assumption is that logical and mathematical structures need to be
expressed \mdash{}and indeed \i{located} \mdash{} via a few canonical
theories, such as first-order predicate logic, first-order
quantification, and (perhaps modally inflected) set theory.
So 20th-century Analytic Philosophy unfolded in a
kind of semi-explicit dialog with disciplines such as
linguistics and computer science which promoted, from the
full spectrum of mathematical domains, a handful of
specific paradigms \mdash{} like Zermelo-Frankel Set Theory,
Tarskian quantificational-predicate logic, and
Hindley-Milner type theory \mdash{} that in turn became equally
paradigmatic in fields like formal language analysis and
\q{knowledge enginering}.  This is not to imply that
linguists and computer scientists took their marching orders
from philsophers \mdash{} it's probably equally true that Analytic
Philosphers gravitated to these kinds of mathematical frameworks
because they were proving empircally useful in other fields.
But we can say that a group of disciplines aiming to
technically examine subjects at the intersection between
human thought and formal systems \mdash{} language, knowledge,
valid inference, Artificial Intelligence \mdash{} collectively
progressed toward a vision of what logicomathematical
structure \i{is}, a paradigm (or network of paradigms)
which undergirds any reception of phenomenology insofar
as a \q{Naturalizing Phenomenology} project seeks to
initiate dialog and comparisons between how phenomenology
and Analytic Philosophy, respectively, treat logic,
mathematics, or formal-systematic explanations in general.
}
\p{Interestingly, the logico-mathematical formulations that
arguably have come to dominate methodology on the scholarly
\i{frontier} between science and humanities \mdash{} lingustics,
cognitive science, computer science \mdash{} are different from
those that seem to dominate mathematics and physical sciences
themselves.  Of course, some logical form is intrinsic to
any formal method since formal systems cannot be blatantly
self-contradictory.  But the vehicles by which criteria
of logical consistency are expressed in, say,
linguistics and philosophy \mdash{} presenting sentence-meaning via
first-order logical glosses of their predicate
(locutionary, assertorial) semantics, for instance, or
equating meanings to the conditions in the world
(or possible worlds) that would make linguistic expressions true
\mdash{} these first-order and set-theoretic intuitions do not
necessarily coincide with working mathematicians' visualization of
the spaces where their own axioms and formulae are in force.
Category theory, for example, recognizes that set and predicate logic
are in a sense emergent properties of the different formal domains
where mathematics is enacted \mdash{} sheaves, topoi, manifolds, sites, and
so on \mdash{} which are not isomorphic, propagating to variation
in how logic itself is understood.  Insofar as there is no single
formalization to specify what a \i{predicate} is, or a \i{set},
we are on shakier foundation if we want to apply these notions to
conceptual or linguistic investigations \mdash{} to consolidate the concept
\q{cat}, say, as the set of all \i{felix catus}, or
the meaning of \i{Hugo is sleeping on the sofa} as a particular state
of affairs pertaining wherein a certain tabby cat sleeps on a
certain piece of furniture.  The problem with this
sort of logico-mathematical gloss is that it either employs sets and
predicates in an informal manner \mdash{} which brings us back to
the circularity problems I analyzed earlier, since
an honest assertion that Hugo is sleeping on the sofa obvious
intends to report, modulo the imperfections of informal logic,
that as a point of fact Hugo is sleeping on the sofa \mdash{} or
it tries to elevate the logic from informal to formal,
which legitimately then adds explanatory value to our technical
grasp of the sentence, but then we get tangled in questions like
what sets are.
}
\p{To put it differently, it certainly seems as if the
sentence \q{Hugo is sleeping on the sofa}, as an intended and
sensed artifact in a human world, in some fashion encapsulates and
\i{designates} a proposition, and/or the state of affairs, that this
tabby cat is in repose on that piece of furniture.  The sentence
(deliberately produced sound patterns, perhaps) provides us
with mental access or handle to that proposition, making it an
unitary and isolated object of thought, as well as a thought that
can be shared between people via language.  There is a predicate
there that seems to be \i{intended} in the phenomenological sense
by the sentence as a linguistic gestalt, just as Hugo is intended
in our visual consciousness when we explicitly see him on the couch,
or at least analogous to that regime of intentionality
insofst as Husserl finds both parallels and contrasts between
\i{apophantic} and preceptual-integrative analysis.
That is, the predicate or fact of Hugo's sleeping at that spot
is the noema to the sentence's phenomena, just as Hugo is the
noema to our seeing-Hugo.  In this sense it is non-circular to
proceed with \q{Hugo is sleeping on the sofa} to mean the
fact of Hugo's sleeping on the sofa, since we are exploiting the
seemingly reundant repetition of the sentence to designate
indirectly a predicate which we have no way of expressing except
through the sentence.  So there is \i{something} there that seems to
have a logical form we need (at least if our goal is theoretical
completeness) to cast light on.
}
\p{But when we then pass from informal to formal logic
we encounter the gap separating mathematics circa 2019 from
mathematics circa 1919.  Hugo, for example, is a cat,
which (in the precincts of informal logic, i.e. a logic
solid enough to shape communication and coordinations of
thought but backgrounded as only one of many structural
parameters delimiting how thoughts and the world avolves)
means that Hugo is a member of the set of all cats.  But
what is a set?  Garfield (a fictional, comic-strip cat)
is also a cat but he does not \q{exist}, so the set of all cats
cannot be enumerated just by (even as a mental exercise) pointing
to them one-by-one.  Moreover, what is referenced by the word
or idea \q{Hugo} in the thought that Hugo belongs to the set of
cats?  We can look at and point to Hugo to get a rough idea; but
the fur Hugo just shed is not a member of the set of cats, even
though it is (or was) a part of Hugo.  It is by no means
obvious that we can, even as another mental exercise, define
\q{Hugo} by an enumeration (metaphysically possible even
if logistically infeasible) of hairs, or cells or molecules, that
via mereological summation grounds the name in a
logically well-established referential framework.
}
\p{On course, it is possible to logically restate a sentence like
\q{Hugo, who is a cat, sleeps on the couch} in first-order logic,
which seems to be arriving at a formal exposition of the predicate
informally intended by the sentence \mdash{} the noema to the
sentence's phenomena.  But formal machinery
only papers over the logical incompleteness of concepts like
\q{set} and \q{predicate}: a logical denotation of the sentence
\i{seems} to refer to the state of affairs that the sentence
tokens (and not just the sentence itself) via quantifier
symbols and other logic artifacts.  I'd argue, however, that these
elements are meta-discursive cues rather than theoretical posits:
they invite us to entertain the thought of the predicate lurking
beyond the sentence via a departure from discursive norms,
the same way that simply repeating the sentence does.  In short,
it may be non-circular to gloss \q{Hugo, who is a cat,
sleeps on the couch} via formations like \ExistsX{} and the
logical form gets us closer to apprehending the abstract predicate,
but this is not because of logical symbols like \q{\Exists{}}
in themselves, but because of their rhetorical effects as
interrupting the flow of conventional philosophical writing.
Symbols of quantification, set membership, or part/whole
inclusion only carry the imprimateur of logical rigor
insofar as we can synchronize what flavor of mathematics
we are using with which to theorize quantification or parthood.
}
\p{We can, for example, argue that Hugo is an extended
three-dimensional thing (or maybe a four-dimensional
perdurant, since he's still here day after day), and
any reasonably compact and continuous \ndim{} manifold can
potentially be designated in thought and language.  This is
reasonable, but it positions our infrastructure for the semantics
of proper names (and a large class of referring expressions)
against a topological backdrop in which we can talk about sets
of points, or perhaps regions spanned by extremum points
(Hugo's fur as the convex hull of Hugo's solid mass) or just
reasons as first-class posits in a point-free algebraic
topology.  We can then present a case for the semantics of
\q{Hugo} as a conventionalized sound adopted by a small group
of humans to designate a particular four-dimensional manifold
\mdash{} at least if we have a sufficient topological or Categorical
foundation to give logical exactitude to concepts like
\q{four-dimensional manifold}.  This still neglects
merological problems, like Hugo losing some fur
but still being himself; either the manifold Hugo
is scattered wherever his fur (and, well, peepee and poopie)
are or the Hugo manifold has fragmented into lots of disjoint
other manifolds, only one of which is \q{Hugo}.  This latter
option suggests that we need to identify the \q{essential parts}
of Hugo, not the inessential parts like those in the litter box;
then the Hugo manifold is the one that contains those parts
\mdash{} i.e. we need a mereological formalization to embed in
the topological one.
}
\p{It is certainly possible that by patching philosophical incompletions
and importing mathematical vocabularies we can piece together a
reasonably compelling theory of mereological/topological
complexes, maybe with graded sets (some parts are more essential than
others), to logically model the semantics of singular referents.
By taking sets of such complexes, maybe
with a modal machinery (for, say, fictional cats) or
graded membership (some things are only partially
subsumed by their concept: is a pet malawolf a dog?) we can
similarly flesh out a set-theoretic semantics, and combining
the singular and set analyses perhaps achieve a logically
inflected semantics applicable to a broad class of linguistic events.
But have we therefore truly surpassed informal logic,
coiled within dialogic conventions, toward a rigorous
theory of linguistic meanings via formal logic?  Do we need
topology, mereology, graded sets, and modal logic because these
are intrinsic to the epistemic domain where signification
happens, or are we just grabbing an assortment of mathematical
tools to plug the leaks in a philosophical attitude wherein
there is a relatively straightforward formalization of
quantifier-predicate logic that can polish the informal
logic of language, but which perhaps failed to
anticipate the mathematical complexities of formalizing what
exactly are the \q{thises} that quantifiers can quantify over?
}
\p{These questions are directly relevant to phenomenology \mdash{}
and to Analytic Philosophy's reception of and dialog with
phenomenology \mdash{} because a phenomenological account of
issues like manifold synthesis and mereological essentialism
can easily bring forward formulations that overlap
with a \q{logistic} paradigm in some details but not others,
and profitably pursuing the dialog depends on identifying both
similarities and differences.  The story of Hugo as a
four-dimensional manifold might not be told in a phenomenological
context, but we might study our visual episodes involving Hugo
in terms of perceptual synthesis: seeing Hugo from the front
we anticipate that we \i{could} see Hugo's tail if we were
standing behind him and it was not as much tucked beneath him;
this anticipation is also a formula for how we integrate
successive impressions as our gaze travels from his ears to
his tail that we see partially.  There is a certain continuity and
topography of our accumulation of successive visuals that
retraces a sort of manifold, even if now we're talking
not about Hugo (or his body) himself but rather the
interplay between the space where he's laying and my own
perspctive and vantage.  So, even
on a phenomenological basis we encounter fragments of analysis
that point toward something like a mathematical treatment \mdash{}
but that doesn't mean we are necessarily committed to a
quantitative-predicate logical theory of preception, or
equally important, of the communicative acts that
perception yields.
}
\p{In pursuit of logical rigor, 20th-century Analytic Philosophy
(plus, say, linguistics, cognitive science, and fields related
to knowledge engineering and information technology, like
AI-driven medical diagnosis) have highlighted
specific logico-mathematical representation-regimes,
building up a paradigm of logicality (as a philosophical
quality, something one senses in analytic methods)
and propositionality (as a corresponding analyzand).
The various layers of this construction address different
aspects of logic, or \q{logicality} as such, such as
predication (that which can be qualified as true or
false in the first place), quantification (establishing
a scope for empirical truthness, so logic does not not decay
into tautology \mdash{} that \i{one} thing or \i{some}
members of set have a property where others don't, modeling
contingency and a posteriority), modality
(capturing the further contingency that most facts
could have been otherwise), mereology and individuation,
quantification scope, supervenience, and other logical
themes that seem intrinsic to any intellectual territory
where predication and true-falseness pertains.
Codification of these facets of logic not only
fine-tunes philosophical analysis; it also provides a
semantics for natural language insofar as we accept that
meaning is centered on locutionary truth-conditions (from
which phenomena like illocutionary force and polarity
\mdash{} our feelings and pragmatic intentions \visavis{}
states of affairs at the predicate substrate of language \mdash{}
can be reached in supplemental analysis) and a philosophical
basis for the formal/computational representation of facts
and knowledge (such as, in database design).  So
insofar as a philosophical-computational-linguistic consensus
emerged on the broad outlines of what logic entails,
this paradigm spread into multiple disciplines,
scientific as well as humanistic, and applied as well
as purely speculative.  Insofar as Naturalizing Phenomenology
involves dialog with Analytic Philosophy, it
intersections with these different faces of logicality and
the mathematical systems which prior generations believed best
captured them.
}
\p{But this does not mean phenomenology needs to endorse the full
stack of paradigms layered in Analytic Philosophy's
excavation of logic; that we need to approach logical-mathematical
eidetics via quantifier-predicate logic, for example, or to
prefer the mathematical frameworks through which
quantification and predication are usually modeled.
Mathematics itself pulls in different directions: quantification,
precication, and set-membership have different
theories in different contexts, e.g. mereotopology compared to
algebraic topology compared to Category theory.  These
lines of development are largely tangential to Analytic
Philosophy partly because the philosophical (and, say,
cognitive-scientific) push toward formalized logic is not really
a push toward \i{mathematics}, but rather a vision of
bringing rigor to what is still essentially \i{informal} logic.
That is, pursuing for instance a formal treatment of
quantification to undergird natural language semantics
does not mean building up to a rigorous, maybe
Categorial theory of quantification in terms
of (say) dependent coproducts, because what \q{logicality}
does in language does not depend on that degree of technical
detail.  But that means in effect that operationalizing
the \q{logicalness} of language is not so much a problem
in \i{formal} (mathematical) logic but a formalized
\i{informal} logic, logico-linguistic or logico-mental
(as in \q{folk theory of other minds}) more than
logico-mathematical.  Fair enough; but this then places other
analytic conveniences (like referents as 4-dimensional
manifolds) on shakier ground, because we're no longer
in an all-encompassing logico-\i{mathematical} paradigm.
Topology and Category Theory, in effect, reveals the difference
between formal-mathematical and formalized-informal logic.
Trying to negotiate with mathematics as source of methodological
rigor, mathematics proves to be an elusive partner.
}
\spsubsectiontwoline{Phenomenology and the Limits of Logic-Compositionality}
\p{The trajectory through Algebraic Topology and Category Theory
has been one vector of post-1950 mathematical progress.
Another has been computer proofs, proof-assistants, and
computer-inspired approaches to mathematical foundations, like
Homotopy Type Theory.  Regarding \i{types} rather than
\i{sets} as the primordial elements of logic and mathematics
positions mathematics and computer science as two different
worlds evolving from a common type-theoretic universe.
In this environment, mathematics as abstract thought gets
mixed with mathematics as a technological discipline;
specialists for example debate the merits of mathematical
systems formalized with a priority to automated proof-checking
vs. presentations more aligned with mathematical precedents.
The proper place for technology and automation becomes
debated.  Computer theorem-provers after all are not oracles;
they need their own quality-checking, their own design theory.
Do we trust the theorem-prover (itself a human artifact) more
than the conventional deliberations validating a proof
in the eyes of mathematicians (and if so, why)?  Are design
principles of proof-assistant software and languages topics
for mathematics or computer science?  Does a test suite to
demonstrate a proof assistant's trustworthiness bear witness
to mathematics becoming suddenly an empirical rather than
eidetic science?  Does the empirical contingency of the
correctness of software used to provde the 4-color theorem
make the theorem something other
than \q{synthtic \i{a priori}}?
}
\p{There are analogs to these questions in the realm of
Analytic-Philosophical logic.  For example, the idea that cognitive
or linguistic meanings can be reduced to first-order predicate
formulae can be investigated by considering computer technology
which does in turn model digital content, in effect, as predicate
assertions together with rules of inference (that can
in principle be modeled via first-order logic): the Prolog
programming language, for instance; certain database
query systems; the Semantic Web and \q{formal Ontologies}.
These technologies operationalize logistic paradigms, but they
also put an empirical face on a logico-mathematical background
that for Russell, Frege, and Wittgenstein would have been
abstract speculation, akin to how proof assistants
put an emprical tie on abstract mathematics.  Similarly,
formal linguistic approaches (like Categorial Combinatory
Grammar) base themselves on abstract logico-mathematical
paradigms (like Hindley-Milner type sysetms), but by 2019 we
have a panoply of linguistic schools whose models
are realized in computational linguistic technologies,
digital platforms that represent and/or automatically
identify semantic and grammatic patterns according to a
specific thory of language.  In this context, type theory
informs linguistic research at a theoretical level as part
of its logico-mathematical underpinnings, but it also belongs
to the technological infrastructure wherein linguistic
technology is implemented.  These two roles can be complimentary:
insofar as Classical Hindley-Milner type theory has been
supplanted by competing paradigms in software which proves
technologically advantageous for modeling natural language,
we can explore whether more modern type theory is
superior for linguistic analysis at a purely theoretical level too.
}
\p{Insofar as Analytic Philosophy converged in the last century upon, it
seems, a fairly consistent nucleus of logical paradigms
\mdash{} covering the various intrinsic facets of quantifier-predicate
logic \mdash{} the heritage of that philosophy is similarly
rebased by 21st-century technology.  The persistence of first-order
logic as a reified example and medium of logic as such was driven by
the abstract reality that almost any logical system can be translated
first- or secod-order logic in principle (at least if we can
tolerate whatever metaphysical posits may be needed
to ground the resulting logic expressions, like Possible Worlds).
As long as logic belongs just to pure thought, the formal/theortical
isomorphism between superficially different theoretical
architectures can be treated as a challenge to overcoming:
converging on first-order predicate-quantifier logic
as a common ground beneath various shallower logics may seem
(and perhaps did seem to early Analytic Philosophers)
to be explanatory progress, a theoretical repositioning
which really does evince a situation where philosophical
disputation really can progress toward greater clarity.  However,
now that different logics are indeed manifest in a concrete sense
in different digital technologis, their contrasts no longer seem
so shallow.  It could be argued that the illusion of a \q{deep}
logic underlying superficial variation (like the illusion
of a \q{deep grammar} stabilizing the morphosyntactic surface,
perhaps) is itself the superficial presumption, failing
to duly appreciate surface contrasts which are only manifest
in the register of applied technology, rather
than abstract thought.
}
\p{For example, one can make a strong
semi-theoretical and semi-implementational case that Graph Theory
or Process Calculii are a better foundation for projects like
programming-language compilers and database query systems than
conventional predicate logic \mdash{} notwithstanding that
propositions in graph or process algebras can potentially be
reformulated in first-order logic (and vice-versa).  The
issue with such in-principle reductions is that they are of
dubious explanatory merit; the pro-forma restatement of
graph structures in predicate logic, for instance, has no
particular importance unless we are prior committed to
predicate logic being somehow fundamntal, in which case
the reducibiloty of graph theory (since this reducibility
can be inverted \mdash{} we can provide graph models of
attribute predication, boolean operators, etc.) can't
justify the commitment on pain of circularity.  What
\i{can} be observed, however, is that graph-database
queries cannot be optimized in their first-order
translations (which is why graph databases are a different
technology than relational databases), which in turn suggests that
the structural differences btween graphs and propositions are
philosophically non-trivial, and consequential.
}
\p{Similarly, concrete technologies can point
us to an intuition of
sentence-comprehension as a calculus of interlocking
cognitive processes, not the mechanical unzipping
of a predicate compact, seeing that there are substantial
implementational differences between process calculii and
predicate logic as paradigms for programming languages.
In short, technology reveals how formal systems may be
substantially divergent from predicate logic even if the
logical structures governing their operation have
isomorphs in predicate logic: technology reveals that such
\q{in principle} isomorphs are less important, less
indicative of philosophical depth, than we may have believed a
century ago.  But this is revealed not within philosophy
alone, but in an interdisciplinary spirit where insights are
mined from fields like compiler design and database engineering;
and from practical experience as well as speculation.  We can
appreciate the substantiveness of
\q{surface-level} logico-structural variation when we actually
\i{write} compilers or database query engines.
}
\p{The phenomenological de-centering of metaphysics does
not mean a complete abdication of metaphysical priorities,
of the need for a ledger to assess the reach and limits
of thought, or science or intellectual explanation;
to place the norms of philosophical exposition and
conversation on a disputable, theorizable conceptual bedrock.
Bracketing metaphysics does not mean suspending what
metaphysics is supposed to do to philosophical performance.
But Phenomenology seeks to bracket metaphsyics as something
received, habitual, or institutioanlized; something that
influences scholarship more than it is studied.
Phenomenology brackets metaphysics in the hope that it can
be recreated and return, but more transparent and self-aware
than before, like a professional athlete who returns
from a league suspension chastened and technically sounder.
Metaphysics is suspended but not expelled.  There is an
analogous dialectic in open-source progrommaing;
institutional norms for software quality are both suspended
and resinstated in a more transparent, technical
incarnation.  The de-centering of \i{institutional} controls
does not foreclose a re-centering of scientific
convergence in a shared technical understanding of the
project.  Similarly, the decentering of \i{institutional}
metaphysics \mdash{} not just the ideas, but the academic workflows
and publishing conventions that support them \mdash{} does
not foreclose a recovery of metaphysics in the
philosopher's relation to her own consciousness, in
writings as an existential trace more
than a social production.
}
\p{All of this then produces the question: what should metaphysics
be if it is not the paradigms we have received?  Where
does a phenomenological analysis take us that we were not before
it started?  As with any writing, we have to start with
the question: now that the writing has finished,
what has changed?  Aniticpating the writings of the future:
what do we want to have changed?
}
\p{We have to communicate this to Analytic Philosophers
(and practitioners of many other branches of scholarship,
like cognitive science and linguistics): what changes
when a phenomenological study is complete?  What
changed when there was a \i{Thing and Space} or a
\i{Phenomenology of Perception} compared to before
there was?  Failure to understand what a writing
wants to change can easily yield to failure
to understand the writing.
}
\p{I have often felt that Analytic and phenomenological
practitioners were talking past each other precisely in
this sense: a specialist in the philosophy of mind or
of science might approach a reading or conversation
thinking the phenomenologists is working to one kind
of telos, while the other is going somewhere else.
In apparently the most common iteration, the
Philosopher of Mind thinks the phenomenologist is aiming
toward a kind of skeptically inflected science of
perception, a theory of the deep structures behind
perceptual experience that potentially takes us further toward
an unbiased and even scientific theory by bracketing
presupposions about how consciousness and cognition
are supposed to work.  The Philosopher of Mind might
then think that phenomenology should, let's say,
uncover the synthesizing and interpretive acts that
constitute my raw sensate awareness as the grasping
and judging that Hugo is sleeping on the sofa, or that
it is half-past four and the sun is setting outside, or that
the car double-parked outside is red.  By attending
in a hopefully presupposionless way to the immediacy
of consciousness, we can see the underlying stitching
and unifying that normally passes beneath the threshold
of awareness, as if we were observing our
own consciousness like a shaman on ayahuasca.
}
\p{I have never in a philosophy classroom, nonetheless, felt
that my role was to issue oracular pronouncements as if from
a drug-induced trance.  Consciousness is not usually
extravagant and cosmic; my awareness of Hugo sleeping
on the couch is, really, not much more than
just Hugo sleeping on the couch.  Notice the structural
analogue to logistic philosophies of language: just as
adding logical symbols to \q{Hugo is sleeping on the sofa} does
not guarantee we are uncovering some deeper structure in the
language, some precious scientific insight, nor does
superimposing talk of a Passive Synthesis of Perception or
Hyletic Immediacy on the mundane episode wherein I see
Hugo sleeping on the sofa make that perception suddenly
scientifically tractable, as if we can solve the
mind-body problem just by concentrating hard enough.
It is what it is.  The things themselves are the things
themsevles; if we're going to them we can't assume we're
en route to some revelation or science, as if the
laws of consciousness are present in the room alongside
Hugo, the sofa, and the lamp.
}
\p{But a reframing of the sentence \q{Hugo is sleeping on the sofa}
\i{can} take us at least modestly further toward a good theory
of language, something beyond the disquotational
verity that a sentence means what it means, just as
\q{it is what it is}; \q{Snow is white} means that snow is white.
Similarly, I believe reframing perceptual episodes can
shed light on the formal \mdash{} and perhaps even scientific
\mdash{} structures of cognition, but we have to clarify ahead
of time what magnitude of clarity we are hoping for;
what precision of understanding landmarks a successful
analysis.  If a scientist reads phenomenology
hoping for an \expose{} of consciosuness that is both
first-person accurate and naturalistically revelatory
\mdash{} something experientially faithful like a
physchological novel but clinically analyzble
like a Brain Simulation \mdash{} then what is actually on
offer will fall short.  Both attending to consciousness
as it is experienced and honoring professional and
scientific standards of intellectual rigor are guiding
principles of phenomenological research, but these
are principles, not minimal standards like
a one-drink minimum at a night club: a very good
phenomenological writing will \i{not} truly be
either stream-of-consciousness realistic like a
novel nor scientifically complete like a Neural Network
platform presented at a computer science conference.
The point is via phenomenology to move a little closer
to both experiential honesty (to accept
consciousness as it is experienced, not as something
we hypothetically transform to fit cognitive and
physical explanations) and scientific rigor
(both the science of neurophysical explanations and
the science of cognitive or intelligence simulations),
but while doing so acknowledging the incompleteness of
both projects and the tensions between them.
}
\p{In the case of perceptual synthesis, we have good reseon
to believe that biophysiology and experimental psychology
gives us a well-informed acccount of perceptual phenoema,
in their episodic unity, at the neurological and
sensori-motor levels; presumably I have some brain
function that unifies white-sensations and brown-sensations
and distributes the color regions to different concepts
(sofa, Hugo).  Since I'm not aware of these
levels they are not phenomenological subjects per se
(which doesn't mean we should reject them entirely,
since I have reasonably firm beliefs about the scientific
basis of perception, beliefs which I don't deliberately
question any more than I deliberately apply them to
mundane perceptions).  Perhaps our scientific training
casuses our actual perceptual content, at last sometimes, to be
subtly different in our experience of it.  I am reminded
of the rainbow color of a thin slick of spilled oil:
knowing to some approximation how we think that phenokenon
occurs makes me explicitly associate the rainbow with the
oil, and therefore to almost unconsciously direct my
gaze to the border lines where the colors are brownish and
less pronounced.  Somone who did not impose a mental
explanation on what they're seeing, even if unintentionally,
would perhaps attend more intently to the color and
experience it more vigorously.  So scientific beliefs
are relevant in the limited sense that their trace can
sometimes be detected in actual experience.
With this qualification, however, any scientific
story about how percption \i{works} lies outside
phenomenology proper.
}
\p{Over-analysis, \i{excess} clarity, \i{too much}
resolution distorts experience even if it aims
to explain it.  So phenomenology does not neatly
fit within a scientific paradigm where magnification
often leads to elucidation and discovery: we cannot take
a microscope to consciousness and find new candidates
for scientific treatment, the way we can advance science
by making more powerful telescopes.  For a scientist
or someone with scientific intiutions it can seem
hard how to go forward: you don't study cancer by staring
at tumors; so how can you study consciousness without some
analysis, some distortion and defamiliarization?  Otherwise
we seem fated to endless circularity: \q{it is what it is}.
The things themselves are what they are.
}
\p{This point, relative to perceptual synthesis, has
similar forms in other aspects of cognitivity activity:
judging that the animal on the sofa is a cat, indeed Hugo,
that he's sleeping, that it's daytime becauuse it is
light outside, that the men in uniform skating on
an ice rink on television are hockey players, and so forth.
Phenomenology cannot ignore that many of the perceptual
and conceptual judgments that we are aware of take the
form in consciousness of settled fact; we can't do first-person
analysis of chains of reasoning that we don't actually
experience ourselves performing.  So while
presumably there is a detailed story to be told about
how optical and inferential thought-processes lead
us to the conclusions that it's daytime and that Hugo
is sleeping, the perceptual record shows that
none of this labor was special enough to warrant conscious
attention.  It's not like I first speculated that some
unfamiliar creature was on the sofa, and subsequently
concluded that it's just Hugo; or that I had to stand
and ponder who the cat is until I could assemble
enough visual detail to reach that conclusion.
I honestly don't know what I was unconsciously
thinking to presume passively that it's Hugo \mdash{} honestly,
I usually don't know what I'm unconsciously thinking.
I cannot, accordingly, use unconscious or
too-mundane-to-experience judgments as part of
phenomenology at least during argumentational stretches
where I'm sticking to phenomenological reports.
}
\p{Insofar as this is characeristic of phenomenology in general,
we should not expect that phenomenology will advance
us further to some sort of rigorous philosophy
or experientially faithful science by revealing things
about experience that we have not yet noticed due to
lack of attention: we cannot magnify consciousness
the way a microscope enlarges different things and a
telescope reveals distant ones; metacognition
and first-person reflection don't enlarge or
polish the things in consciousness so that we can
describe them more thoroughly or scientifically.
Phenomenology does take us to some new place by
making us aware of pieces of experiential synthesis
that we don't experience directly otherwise,
as we don't experience the molecules in
water; or minutiae of conceptual
reasoning that we would otherwise take for
granted, like how we figure out when an
animal is a cat.
}
\p{Of course, however, often we \i{are} aware of
perceptual synthesis and conceptual judgment.  If
after seeing Hugo I walk toward him,
with my perspecive gradually changing, I
consciously register a change in how Hugo appears to me
which produces an explicated synthsis,
a self-conscious accumulation of visual evidence telling
me about Hugo's current state, at least if I am
looking at him with express interest him as opposed to looking in
his directly and passively noticing him while attending to
something else, like the television.  Likewise I may passively
observe that he is sleeping if I previously thought he
was sleeping and just happened to glance at him again,
but there's an alternative scenario where I don't
specifically assume he is sleeping, but upon looking
at him rather attentively, with his lack
of movement and his relaxed posture, I reach the
conclusion self-consciously.  Walking outside I may
see a familiar stray cat and put no thought into recognizing
her, but I may also spot a movement in the distance
and look purposefully before identifying the animal as a cat.
In short, perceptual synthesis and conceptual
judgments are neither \i{intrinsically} conscious
or pre-conscious, active or passive; many
cognitive episodes that we are explicitly aware of in their
unfolding could, in slightly different circumstances,
remain below the threshold of conscious awareness.
Conversely, many judgements we make without
much conscious attention could alternately be
experienced with more deliberation and attention.
Indeed, sometimes we revise our passive judgments
when something goes awry in our passive estimation of
what's going on around us \mdash{} we assume we hear
the dog barking in the back yard, but then
we see her eating pizza crusts in the kitchen, so now
we have to account for the mysterious barking outside.
We all have the faculty to quickly adopt the
mode of conscious investigation, trying to gather specific
information, when the passive judgments we have just
assumed to be true turn out to have some incongruity, some
gap in their logical order.
}
\p{Now we are moving toward some insight, since although we
can't give a phenomenological account of preconscious
judgment, we can at least confirm that \i{some}
judgments are \q{active} and part of
explicit awareness, while others are passive and
not (in any inner details) noticed; we can also say that
at least some passive jugments seem similar in form
and outcome to active ones.  This raises several
avenues for further, maybe scientific continuation:
why do we have both active and passive judgment?  What
benefits accrue to us biologically as animals that depend
on precise judgments to survive that we make many
instinctive judgments and some additional deliberated ones?
Why are active and passive judgments often similar?
What triggers a decision to actively attend to specific
components of experience so that associated judgments in
that one area will, at least in the immediate aftermath,
be more active than passive?  What is actually
neurophysically different in the two cases so that,
in a pure subjective and first-person sense, I know when
a judgment is active and passive?
}
\p{Furthermore, since many judgments \i{do} rise to a level
of consciousness \mdash{} and not just the minimal awareness
of believing such-and-such but some explicit experience of
how the judgment is formed, of their inner pieces and
intemediate stages \mdash{} these conceptual and perceptual episodes
\i{are} candidates for phenomenological analysis.   We are
certainly aware of perceptual episodes in their
unfolding; how perception scans over a room, or
any other surrounding, like the pan of a camera;
while, also, different episodes piece together,
with shifts in attention (and, often, kinaesthtic
actions) \mdash{} like looking up from a book \mdash{}
juncturing disparate perceptual phases.  It seems self-evident
that perception unfolds according to multiple time scales,
where episodes measured in seconds join to become the perceptual
background as we engage in activities over the course of minutes,
within errands that stretch toward hours and engagements
over the course of a day (we cook, go to the store, ride
a subway, watch a show), and of course we plan for a future
and remember a past over monrhs and years.  Cognitive acts and
perceptual interludes can be slotted into different time
scales, affecting how we experience them and should analyze
(or even name) them: a perceptual episode, brief enough to appear
in consciousness as a single event, needs a different
theoretical registration than experiences implying lengthier
segments.  Practical actions, also, bridging
perception, conceptualization, and planning, sort into different
time-spans: opening a bottle vs. preparing a meal;
petting the dog vs. walking the dog.
}
\p{We have, as such, I would argue, a collection of observations
and distinctions we can make about consciousness and
conscious experience, that emananate from consciousness
itself and not from prior scientific or
metaphysical commitments.  These observations
seem presuppositionless, or at least common-sensical.
We should not take that too far: the \q{suspension of
belief} should not commmit us to a virgin brilliance
exorcising millenia of thought, posturing us as priests
of intellectual purity that can be defrocked by the
smallest trace of naturalistic thinking \mdash{} gotcha, a
presupposition!  Philosophy as volumes forging profound
insight from pure thought alone seems silly in our time,
more art created from ideas than rational exposition.
But we can make a good-faith effort to present
patterns in consciousness based on rather casual and
common-sense reflections not explicitly aligned with
scientific thories, or with schools of philsophy,
and I claim at least some of the intellectual structure
we can get to via that route takes the form of observations
I have highlighted: passive vs. active judgment; the
various scales of perceptual and enactive temporality; the
individuality, aggregation, and junctioning between
perceptual episodes.
}
\p{Insofar as these observations
and distinctions form the rudimentary constructs in
a theory of consciousness, the explanatory value
of that theory will depend on how it can be
consistently extended.  Remaining within phenomenology,
we cannnot analyze deeper into consciousness \mdash{}
capturing more fine-grained layers of perception, uncovering
more details of passive synthesis \mdash{} because we cannot
impose scientific-empirical data or philosophical speculation
on subjective experience.  So if we can only articulate a
handful of organized but commonsensical comments about
experience, we need to find away to make the analysis meaningful:
for it to make a difference.
}
\p{One way to do this, I believe, is to make the common-sensical
observations more than just observations; not in the
sense of importing interpretive claims that don't
unfold from them organically, but in the sense of trying to
express their as a formal structure, something which can
be then set in relation to other structures and other
accounts of mental processes.  There are many formal systems
that have been proposed to cover various features of
cognition and intelligence; but only a few
of these, I suspect, are formal views on experientially vivid
acts of mind; on judgments or percpetual syntheses that
we can experience in their unfolding and reflect upon
first-personally.
}
\p{There is a considerable body of cognitive-scientific
literature that aspires to formalize (and sometime simulate)
mental activity via structures that have some blend
of logical, mathematical, and computational
specification.  There is also a considerable body of
literature that remains true only to the first-person
perspective, not explicitly trying to align with
scientific explanations and world-view.  But where
these bodies overlap I believe lays the potential for
overlapping theories, like a Naturalized Phenomenology.
}
\p{}

\section{Link Grammar and Type Theoretic Semantics}
\p{From one perspective, grammar is just a
most top-level semantics, the primordial Ontological division of language into designations of
things or substances (nouns), events or processes (verbs), qualities and attributes (adjectives),
and so forth.  Further distinctions like count, mass, and plural nouns add
semantic precision but arguably remain in the orbit of grammar (singular/plural
agreement rules, for example); the question is whether semantic detail gets
increasingly fine-grained and somewhere therein lies a \q{boundary} between syntax and
semantics.  The mass/count distinction is perhaps a topic in grammar more so than
semantics, because its primary manifestation in language is via agreement
(\i{some} wine in a glass; \i{a} wine that won a prize; \i{many} wines
from Bordeaux).  But are the distinctions between natural and constructed objects,
or animate and inanimate kinds, or social institutions and natural
systems, matters more of grammar or of lexicon?  Certainly they engender
agreements and propriety which appear similar to
grammatic rules.  \q{The tree wants to run away from the dog} sounds wrong \mdash{} because
the verb \q{want}, suggestive of propositional attitudes, seems incompatible with
the nonsentient \q{tree}.  Structurally, the problem with this sentence seems analogous
to the flawed \q{The trees wants to run away}: the latter has incorrect singular/plural linkage,
the former has incorrect sentient/nonsentient linkage, so to speak.  But does this
structural resemblance imply that singular/plural is as much part of semantics as grammar, or
sentient/nonsentient as much part of grammar as semantics?  It is true that there are no
morphological markers for \q{sentience} or its absence, at least in English \mdash{} except
perhaps for \q{it} vs. \q{him/her} \mdash{} but is this an accident of English or revealing
something deeper?
}
\p{To explore these questions it is first necessary to consider how a grammar theory can be
extended to and/or connected with a formal or, to some measure, informal semantics.
Here I will present one approach to make this extension \visavis{} Link Grammar.
}
\p{Insofar as grammatic categories do provide a very basic \q{Ontological} viewpoint,
it is reasonable to build semantic formalization on top of grammar theories.
Link Grammar, for example, explicitly derives \q{link types} \mdash{} species of word-to-word
relations \mdash{} by appeal to \q{Categorial} grammars which define parts of
speech in terms of their manner of composition with other, more \q{fundamental} parts
of speech \cite{Kiselyov}, \cite{Rossi}, \cite{MartinPollard}.  The most primordial
grammatic categories are generally seen to be nouns and
\q{propositions} (self-contained sentences or sentence-parts which assert individual states of
affairs), and categories like verbs and adjectives are derived on their basis.  For example, a
verb \q{combines} with a noun to produce a proposition.  \i{Students} is an abstract
concept; \q{Students complained}, tieing the noun to a verb, tethers the concept to an
assertorial flesh, yielding something that expresses a belief or observation.
Meanwhile, Categorial Grammar models not only the semantic transition from abstract to concrete, but
surface-level composition: in English and other \SVO{} language for example the verb should
immediately follow the noun; in German and all \SVO{} languages the verb tends the come
last in a sentence, and can be well apart from its subject.  The semantic pattern in the
link is how the verb/noun pair yields a new semantic category (propositional) whereas the
grammatic component lies in how the link is established relative to other words
(to the left and not the right, for example, and whether or not the words are adjacent).
}
\p{Assuming that surface-level details can be treated as grammar rules and abstracted from
the semantics, we can set aside Categorial Grammar notions like connecting
\q{left} vs. \q{right} or \q{adjacent} (near) vs. \q{nonadjacent} (far).
With this abstracting, Categorial Grammar becomes similar
to a Type-Theoretic Semantics which recognizes, in Natural Language, operational
patterns that are formally studied in mathematics and computer science
\cite{ZhaohuiLuo}, \cite{MaillardClarkGrefenstette},
\cite{ZhaohuiLuo}, \cite{Pollard}, \cite{MeryMootRetore}.
A verb, for example,
\i{transforms} a noun into a sentence or proposition (at least an intransitive verb;
other kinds of verbs may require two, or even three nouns).  In some schematic sense a
verb is analogous to a mathematical \q{function}, which \q{takes} one or more nouns
and \q{yields} propositions, much like the \q{square} function takes a real number and
yields a non-negative real number.  To make this analogy useful, however,
it is necessary to clarify how \q{types} in a mathematical or computational
context may serve as appropriate metaphors for syntactic and/or semantic groupings in language.
}
\subsection{Types, Sets, and Concepts}
\p{Most Computer Science rests on types rather than (for example) sets,
because abstract reasoning about data types
requires some abstraction from practical limitations about how particular values may be
digitally encoded.
Types can be defined as sets of both values and
\q{expectations} \cite{MathieuBouchard} (meaning assumptions which may be made about
all values covered by the type); alternatively, we can (perhaps better) consider types as
\i{spaces} of values.  Types' extensions have internal structure; there
can be \q{null} or \q{invalid} values, default-constructed values, and
so forth, which are \q{regions} of type-space and can
be the basis of topological or Category-Theoretic rather than
set-based analyses of type-extension.  Also, expectations intrinsically include
functions which may be \q{called on} types.  There is definitional interdependence
between types and functions: a function is defined in terms of the types it accepts as parameters and
returns \mdash{} rather than its entire set of possible inputs and outputs, which can
vary across computing environments.  These are some reasons why in theoretical
Computer Science types are not \q{reduced} to underlying sets; instead, extensions
are sometimes complex spaces that model states of, or internal organization of comparisons
among, type instances.
}
\p{An obvious paradigm is organizing type-extensions around prototype/borderline
cases \mdash{} there are instances which are clear examples of types and
ones whose classification is dubious.  I will
briefly argue later, however, that common resemblance is not always a good marker
for types being well-conceived \mdash{} many useful concepts are common
precisely because they cover many cases, which makes defining
\q{prototypes} or \q{common properties} misleading; this reasoning
arguably carries over to types as well.  Also, sometimes the clearest
\q{representative} example of a type or concept is actually not a
\i{typical} example: a sample latter or model home is actually not (in
many cases) a real letter or home.  So resemblance-to-prototype is at
best one kind of \q{inner organization} of concepts' and types' spaces
of extension.  Computer Science develops other
pictures of types' \q{state space}, reflecting the trajectory of symbols
or channels which hold type instances, which at different moments in time
become initialized \mdash{} acquiring a value obtained from a \i{constructor}
function (one \q{type space region} is then demarcated by which values
can be direct results of constructors) \mdash{} then possibly subject to change in the value they hold, and finally
(often) transitioning to a state where the held value is no longer \q{valid}.\footnote{Managing the \q{lifetime} of values from many types, especially \q{pointer} types
(that hold a numeric value representing the current memory address of some other value),
has been a notorious source of programming errors, especially in older computer
languages.  Of late, also, data types often need to be designed to minimize the
risk of data corruption, theft, and malicious code.  For these reasons,
Cybersecurity takes particularly interest in studying types' extensions and transitions
between different values (morphisms within a type space) to formally describe
states or state-transitions which are security vulnerable.
}  Type \i{spaces} have potentially complex patterns of regions and equivalence classes of
inter-value mappings (in the sense of behavioral equivalence relative to
code analysis, testing, or security)
\mdash{} the \i{conceptual} properties of types are expressed in the
\i{internal structuration} of their associated state-space.  Putting this in mathematical
language, an in-depth treatment of types cannot work \q{in the Category} of
sets, even for basic type-extension, but rather (for instance) the
Category of Topological Spaces.
}
\p{Moreover, expectations in a particular case
may be more precise than what is implied by the type itself \mdash{} it is erroneous
to assume that a proper type system will allow a correct \q{set of values} to
be stipulated for each point in a computation (the kind of contract enforced via
by documentation and unit testing).  So state-space in a given context may include many
\q{unreasonable} values, implying that within the overall space there is a \q{reasonable}
subspace, except that this subspace may not be crisply defined.
A value representing someone's age may be assigned
a type for which a legal value is, say, 1000 years, which is obviously unreasonable
\mdash{} the conceptual role served by the \i{particular} use of a type in some context can
be distinct from the entire space of values exhibited by the type.  It is possible to
construct types which are narrowed down to more precise ranges, but in many cases this is
unnecessary or poorly motivated: while 1000 years is clearly too large for an age,
it would be arbitrary to specify a \q{maximum allowed} age (recall that assuming a
\q{maximum allowed} year of 1999 \mdash{} so that the year in decimal only required two digits
\mdash{} led to costly reprogramming of archaic legacy code during \YtwoK{}).
In this kind of situation programmers
usually assign types based on properties of binary representation \mdash{} what number of
binary digits is optimal for memory and/or speed, even if this allows \q{absurd} values
like 1000 years old.  Run-time checks, rather than type restrictions, may be used to
flag nonsensical data and prevent data corruption.  In these scenarios, types represent a
compromise between \i{concepts}, which can be fuzzy and open-ended, and
\i{sets}, which conceptually are nothing more than the totality of their extension.\footnote{Nevertheless, there is interesting (and potentially practically useful) research in
how formal type-constructions model conceptual organization: for example,
\Gardenfors{} Conceptual Space Theory has seen formal implementations
\cite{RaubalAdams}, and it is very interesting to juztapose
scientific and mathematical treatments of Conceptual Spaces (as in \cite{Strle} or
\cite{Zenker}) with mathematical (e.g., topological) thoeries of
data types \cite{Escardo}, \cite{StellConvexivity}.
}
}
\p{Sets, concepts, and types represent three different primordial thought-vehicles for
grounding notions of logic and meaning.  To organize systems around \i{sets} is
to forefront notions of inclusion, exclusion, extension, and intersection,
which are also formally essential to mathematical logic and undergird the
classical interdependence of sets, logic, and mathematics.\footnote{Recent work in mathematics, however (partly
under the influence of computational proof engines and foundations research
like Homotopy Type Theory) shows that type and/or Category theory
may replace sets as a groundlevel for logico-mathematical reasoning (if
not notation) in the future \cite{HOTT} (It is worth pointing out that
despite their similar ordinary meanings, mathematically \i{type} is much different
from \i{Category} even though these respective theories can be usefully integrated).
}
To organize systems around \i{concepts} is to forefront practical engagement
and how we mold conceptual profiles, as collections of ideas and pragmas,
to empirical situations.  To organize systems around \i{types} is to forefront
\q{functions} or transformations which operate on typed values, the interrelationships
between different types (like subtypes and inclusion \mdash{} a type can itself
encompass multiple values of other types), and the conceptual abstraction
of types themselves from the actual sets of values they may exhibit
in different environments.  Sets and types are
formal, abstract phenomena; whereas concepts are characterized by
gradations of applicability, and play flexible roles in thought and language.
The cognitive role of concepts can be discussed with some rigor, but there is a
complex interplay of cognitive schema and practical engagements which
would have to be meticulously sketched in many real-world scenarios, if
our goal were to translate conceptual reasoning to formal structures
on a case-by-case basis.  We can, however, consider in general
terms how type-theoretic semantics can capture conceptual structures
as part of the overall transitioning of thoughts to langauge.
}
\p{A concept does not merely package up a definition, like \q{restaurant} as
\q{a place to order food}; instead concepts link up with other concepts
as tools for describing and participating in situations.  Concepts are
associated with \q{scripts} of discourse and action, and find their
range of application through a variegated pragmatic scope.
We should be careful not to overlook these pragmatics, and
assume that conceptual structures can be simplistically
translated to formal models.
Cognitive Linguistics critiques
Set-Theoretic or Modal Logic reductionism (where a concept is just a set
of instances, or an extension across different possible worlds) \mdash{} George Lakoff and Mark Johnson,
prominently, argue for concepts' organization around
prototypes (\cite[p. 18]{LakoffJohnson}; \cite[p. 171, or p. \textit{xi}]{Johnson})
and embodied/enactive patterns of interaction (\cite[p. 90]{LakoffJohnson};
\cite[p. 208]{Johnson}).
Types, by contrast, at least in linguistic applications of type theory, are abstractions
defined in large part by quasi-functional notions of phrase struture.
Nevertheless, the \i{patterns} of how types may inter-relate
(mass-noun or count-noun, sentient or non-sentient, and so forth)
provide an infrastructure for conceptual understandings to be
encoded in language \mdash{} specifically, to be signaled by which typed
articulations conversants choose to use.  A concept like
\i{restaurant} enters language with a collection of understood
qualities (social phenomena, with some notion of spatial location and
being a \q{place}, etc.) that in turn can be marshaled by sets of
allowed or disallowed phrasal combinations, whose parameters
can be given type-like descriptions.  Types, in this sense,
are not direct expressions of concepts but vehicles for
introducing concepts into language.
}
\p{Concepts (and types also) are not cognitively the same as their
extension \mdash{} the concept \i{restaurant}, I believe, is distinct from
concepts like \i{all restaurants} or \i{the set of all restaurants}.
This is for several reasons.  First, Concepts can be pairwise different
not only through their instances, but because they highlight different
sets of attributes or indicators.  The concepts \q{American President} and \q{Commander in Chief}
refer to the same person, but the latter foregrounds a military role.
Formal Concept Analysis considers \i{extensions} and \q{properties}
\mdash{} suggestive indicators that inhere in each
instance \mdash{} as jointly (and co-dependently) determinate: concepts
are formally a synthesis of instance-sets and property-sets \cite{YiyuYao},
\cite{Belohlavek}, \cite{Wille}.  Second,
in language, clear evidence for the contrast between \i{intension} and
\i{extension} comes from phrase structure: certain constructions specifically
refer to concept-extension, triggering a mental shift from thinking of the
concept as a schema or prototype to thinking of its extension (maybe in some context).
Compare these sentences:
\sentenceexamples{\sentenceexample{Tigers in that park are threatened by poachers.}
\sentenceexample{Young tigers are threatened by poachers.}
}
Both sentences focus a conceptual lens in greater detail than \i{tiger} in general, but
the second does so more intensionally, by adding an extra indicative criterion; while
the former does so extensionally, using a phrase-structure designed to operate on
and narrow our mental construal of \q{the set of all tigers}, in the sense of
\i{existing} tigers, their physical place and habitat, as opposed to
the \q{abstract} (or \q{universal}) type.  So there is a familiar semantic
pattern which mentally transitions from a lexical type to its extension and
then extension-narrowing \mdash{} an interpretation that, if accepted, clearly
shows a different mental role for concepts of concepts' \i{extension} than the
concepts themselves.
}
\p{There is a type-theoretic correspondence between intension and
extension \mdash{} for a type \Tnoindex{} there is a corresponding \q{higher-order} type
of \i{sets} whose members are \Tnoindex{}.\footnote{Related constructions are the type of \i{ordered sequences} of \Tnoindex{};
unordered collections of \Tnoindex{} allowing repetition; and stacks, queus, and
deques (double-ended queues) as \Tnoindex{}-lists that can grow or shrink
at their beginning and/or end.
}  If we take this (higher-order)
type gloss seriously, the extension of a concept is not its \i{meaning}, but a
different, albeit interrelated concept.  Extension is not definition.
\q{Tiger} does not mean \i{all tigers} (or \i{all possible tigers}) \mdash{} though arguably
there are concepts \i{all tigers} and \i{all restaurants} (etc.) along with the concepts
\i{tiger} and \i{restaurant}.  Concepts, in short, do not mentally signify sets, or
extensions, or sets-of-shared-properties.  Concepts, rather, are cognitive/dialogic tools.
Each concept-choice, as presentation device,
invites its own follow-up.  \i{Restaurant} or \i{house} have meaning not via
idealized mental pictures, or proto-schema, but via kinds of things
we do (eat, live), of conversations we have, of qualities we deem relevant.  Concepts do not
have to paint a complete picture, because we use them as part of ongoing situations
\mdash{} in language, ongoing conversations.  Narrow concepts \mdash{} which may best exemplify
\q{logical} models of concepts as resemblance-spaces or as rigid designators to
natural kinds \mdash{} have, in practice, fewer use-cases \i{because} there
are fewer chances for elaboration.  Very broad concepts, on the other hand, can have,
in context, too \i{little} built-in \i{a priori} detail.
(We say \q{restaurant} more often than \i{eatery}, and
more often than \i{diner}, \i{steakhouse}, or \i{taqueria}).  Concepts dynamically play
against each other, making \q{spaces} where different niches of meaning, including
levels of precision, converge as site for one or another.  Speakers need freedom to choose
finer or coarser grain, so concepts are profligate, but the most oft-used trend toward middle
ground, neither too narrow nor too broad.  \i{Restaurant} or \i{house} are useful because they are noncommittal, inviting more detail.
These dynamics govern the flow of inter-concept relations (disjointness, subtypes, partonymy, etc.).
}
\p{Concepts are not rigid formulae (like instance-sets or even attributes fixing when
they apply); they are mental gadgets to initiate and
guide dialog.  Importantly, this
contradicts the idea that concepts are unified around instances' similarity (to each other or
to some hypothetical prototype): concepts have avenues for contrasting
different examples, invoking a \q{script} for further elaboration, or for building temporary filters
(\q{Let's find a restaurant that's family-friendly}; allowing such one-off narrowing is a feature
of the concept's flexibility).
No less important, than acknowledged similarities across all instances, are well-rehearsed ways
\visavis{} each concept to narrow scope by marshaling lines of \i{contrast}, of \i{dissimilarity}.
A \i{house} is obviously different from a \i{skyscraper}
or a \i{tent}, and better resembles other houses; but there are also more nontrivial \i{comparisons}
between houses, than between a house and a skyscraper
or a tent.  Concepts are not only spaces of similarity, but of \i{meaningful kinds of differences}.
}
\p{To this account of conceptual spaces we can add the conceptual matrix spanned by
various (maybe overlapping) word-senses: to \i{fly}, for example, names
not a single concept, but a family of concepts all related to airborn
travel.  Variations highlight different features: the path of flight (\i{fly to Korea}, \i{fly over the mountain});
the means (\i{fly Korean air}, \i{that model flew during World War II});
the cause (\i{sent flying (by an explosion)}, \i{the bird flew away (after a loud noise)},
\i{leaves flying in the wind}).  Words allow different use-contexts
to the degree that their various \i{senses} offer an inventory of aspects for
highlighting by \i{morphosyntactic} convention.  Someone who says \i{I hate to fly} is not
heard to dislike hand-gliding or jumping off mountains.\footnote{People, unlike birds, do not fly \mdash{} so the verb, used intransitively
(not flying \i{to} somewhere in particular or \i{in} something in particular),
is understood to refer less to the physical motion and more to the socially
sanctioned phenomenon of buying a seat on a scheduled flight on an airplane. The construction
highlights the procedural and commercial dimension, not the physical mechanism and
spatial path.  But it does so \i{because} we know human flight is
unnatural: we can poetically describe how the sky is filled with flying leaves or birds,
but not \q{flying people}, even if we are nearby an airport.
Were \q{flying people} used jokingly, it would be in bad taste, like
\q{cat all over all over the driveway} from Pinker
\cite{Pinker} on page 119 and Langacker's \q{Nouns and Verbs} \cite{Langacker87} on page 67.
}  Accordant variations
of cognitive construal (attending more to mode of action, or path, or motives, etc.),
which are elsewhere signaled by grammatic choices, are also spanned by a conceptual
space innate to a given word: senses are finer-grained meanings availing themselves to one construal or another.
}
\p{So situational construals can be signaled by word- and/or
syntactic form choice (locative, benefactive, direct and indirect
object constructions, and so forth).  Whereas conceptual organization
often functions by establishing classifications, and/or invoking
\q{scripts} of dialogic elaboration, cognitive structure tends to apply more
to our attention focusing on particular objects, sets of objects, events, or
aspects of events or situations.  \i{Conceptual} is more abstract and belief-oriented;
\i{Cognitive} is more concrete and phenomenological.  Concepts organize our
\q{background knowledge} \cite{SmithMcIntyre}; cognitions allow it to be latent
against the disclosures of material consciousness
\cite{DavidWoodruffSmith}, \cite{DavidWoodruffSmithChapter},
\cite{JordanZlatev}, \cite{SeanDorranceKelley}.
So the contrast between singular, mass-multiples, and count-multiples,
among nouns, depends on cognitive
construal of the behavior of the referent in question (if singular, its
propensity to act or be conceived as an integral whole; if multiple, its
disposition to either be divisible into discrete units, or not).
Or, events can be construed in terms of their causes
(their conditions at the outset), or their goals (their conditions at
the conclusion), or their means (their conditions in the interim).
Compare \i{attaching} something to a wall (means-focused) to
\i{hanging} something on a wall (ends-focused); \i{baking} a cake
(cause-focus: putting a cake in the oven with deliberate intent to cook it)
to \i{burning} a cake (accidentally overcooking it).\footnote{We can express
an intent to bake someone a cake, but not (well, maybe comedically) to
\i{burn} someone a cake (\q{burn}, at least in this context, implies
something not intended); however, we \i{can} say
\q{I burnt your cake}, while it is a little jarring to say
\q{I baked your cake} \mdash{} the possessive implies that some
specific cake is being talked about, and there is less apparent reason
to focus on one particular stage of its preparation (the baking) once
it is done.  I \i{will} bake a cake, in the future, uses
\q{bake} to mean also other steps in preparation (like \q{make}), while,
in the present, \q{the cake \i{is} baking} emphasizes more its
actual time in the oven.  I \i{baked your cake} seems to focus
(rather unexpectedly) on this specific stage even after it is completed,
whereas \i{I baked you a cake}, which is worded as if the recipient
did not know about the cake ahead of time, apparently uses \q{bake} in
the broader sense of \q{made}, not just \q{cooked in an oven}.
Words' senses mutate in relation to the kinds of situations where they are used
\mdash{} why else would \i{bake} mean \q{make}/\q{prepare} in the past or future tense but
\q{cook}/\q{heat} in the present?
}
These variations are not random assortments of polysemous words' senses:
they are, instead, rather predictably distributed according
to speakers' context-specific knowledge and motives.
}
\p{I claim therefore that \i{concepts} enter language complexly, influenced by
conceptual \i{spaces} and multi-dimensional semantic and syntactic selection-spaces.
Concepts are not simplistically \q{encoded} by types, as if for
each concept there is a linguistic or lexical type that just
disquotationally references it \mdash{} that the type \q{tiger} means the concept
\i{tiger} (\q{type} in the sense that type-theoretic semantics would model lexical
data according to type-theoretic rules, such as \i{tiger} as subtype of \i{animal} or
\i{living thing}).
Cognitive schema, at least in the terms I just laid out, select particularly
important gestalt principles (force dynamics, spatial frames, action-intention)
and isolate these from a conceptual matrix.  On this basis, we can argue that
these schema form a precondition for concept-to-type association; or,
in the opposite logical direction, that language users' choices to employ
particular type articulations follow forth from their prelinguistic
cognizing of practical scenarios as this emerges out of collections
of concepts used to form a basic understanding of and self-positioning within them.
}
\p{In this sense I called types \q{vehicles} for concepts: not that types \i{denote}
concepts but that they (metaphorically) \q{carry} concepts into language, as a bus
carries people into a city.  \q{Carrying} is enabled by types' semi-formal rule-bound
interactions with other types, which are positioned to capture concepts' variations and
relations with other concepts.
To express a noun in the benefactive case, for example, which can be seen as attributing to
it a linguistic type consistent with being the target of a benefactive,
is to capture the concept in a type-theoretic gloss.
It tells us, I'm thinking about this thing in such a way that it
\i{can} take a benefactive (the type formalism attempting to capture
that \q{such a way}).
A concept-to-type \q{map}, as I just
suggested, is mediated (in experience and practical reasoning) by
cognitive organizations; when (social, embodied) enactions take
linguistic form, these organizing principles can be encoded in how
speakers apply morphosyntactic rules.  So the linguistic structures,
which I propose can be formally modeled by a kind of type theory, work
communicatively as carriers and thereby signifiers of cognitive
attitudes. The type is a vehicle for the concept because it takes part in constructions
which express conceptual details \mdash{} the details
don't emerge merely by virtue of the type itself.
I am not arguing for a neat concept-to-type correspondence; instead, a type system provides a
\q{formal substrate} that models (with some abstraction and simplification) how
properties of individual concepts translate
(via cognitive-schematic intermediaries) to their
manifestation in both semantics and syntax.
}
\p{Continuing with benefactive case as a case study (no pun intended),
consider how an ontology of word senses (which could plausibly be
expressed by types and subtypes) can interrelate with the benefactive.
A noun as a benefactive target most often is a person or some other
sentient/animate being; an inanimate benefactive is most likely
something artificial and constructed (cf., \i{I got the car new tires}).
How readily hearers accept a sentence -- and the path they
take to construing its meaning so as to make it grammatically acceptable
-- involves interlocking morphological and type-related considerations;
in the current example, the mixture of benefactive case and which noun
\q{type} (assuming a basic division of nouns into e.g.
animate/constructed/natural) forces a broader or narrower
interpretation.  A benefactive with an \q{artifact} noun, for example,
almost forces the thing to be heard as somehow disrepaired:
\begin{sentenceList}\sentenceItem{} I got glue for your daughter.
\sentenceItem{} I got glue for your coffee mug.
\end{sentenceList}
We gather (in the second case) that the mug is broken \mdash{} but this is never spelled out
by any lexical choice.  It is implied indirectly by benefactive case along with
notions of classification, on the grammar/semantic border, that have a potential
type-theoretic treatment.  It is easy to design similar examples with other cases:
a locative construction rarely targets \q{sentient} nouns, so in
\begin{sentenceList}\sentenceItem{} We're going to Grandma!
\sentenceItem{} Let's go to him right now.
\sentenceItem{} Let's go to the lawyers.
\sentenceItem{} Let's go to the press.
\end{sentenceList}
we mentally substitute the person with the place where they live or work.  Morphosyntactic
considerations are also at play: \i{to the lawyers} makes \q{go} sound more like \q{consult with},
partly because of the definite article (\i{the} lawyers implies conversants have some prior involvement
with specific lawyers or else are using the phrase metonymically, as in \q{go to court} or
\q{to the courts}, for
legal institutions generally; either reading draws attention away from literal spatial implications of
\q{go}). \q{Go to him} implies that \q{he} needs
some kind of help, because if the speaker just meant going to wherever he's at, she probably would
have said that instead.  Similarly, the locative in \i{to the press} forces the mind to
reconfigure the landmark/trajector structure, where \q{going} is thought not as a literal
spatial path and \q{press} not a literal destination \mdash{} in other words, the phrase must be
read as a metaphor.  But the \q{metaphor} here is not \q{idiomatic} or removed from linguistic rules
(based on mental resemblance, not language structure); here it
clearly works off of formal language patterns: the landmark/trajector
relation is read abstracted from literal spatial movement because the locative is applied
to an expression (\i{the press}) which does not (simplistically) meet
the expected interpretation as \q{designation of place}.
We need to analyze syntactic details like noun case and
forms of articles, but also finer-grained (though not purely lexicosemantic) classifications
like sentient/nonsentient or spatial/institutional.
}
\p{One way to engage in classification in this kind of example is just to consider subtyping:
divide nouns into sentient and non-sentient, the former into human and animal and the latter
into artifacts and natural things, and so forth.  But other options are less blunt.
For example, notions like sentient/nonsentient can be construed as \q{higher-order types},
meaning that for broadly-hewed types like nouns or verbs, there are sentient (and non-sentient) variants,
just as for a type \TypeCat{} there are mass-plural and count-plural collections of \TypeCat{},
ordered and unordered \TypeCat{} collections, and so on.  Subtyping, higher-order types, inter-type
associations and various other formal combinations are options for encoding grammatic and
semantic classification in something like a formal type theory.  The key properties of
type systems are not only meanings atttached to individual types but notions of functionality
(according to the central notion that a type system includes \q{function}
types which are mappings between other types; in Category Theory, any
formal type system is \q{Cartesian Closed}, meaning that if \T1{} and \T2{} are
types, there is necessarily a type \TSupT{} of functions between them).
So if adjectives, say, are most basically \NtoN{} (they modify nouns and yield noun-role phrases),
we can then consider how adjectives should be modeled when their modified nouns
are associated with or attributed sentience, mass-plural, or any other variation (whether
via subtyping or some other association).  How these \q{variations} are modeled in accord
with one single type is less important than how they \q{propagate} via applicative
structures, where \q{function-like} types apply transformations and produce phrases.
}
\p{To build up a linguistic type theory, I assume, then, a framework of
types and type associations with a few underlying properties, such as
these:
\renewcommand{\labelitemi}{$\bullet$}
\begin{itemize}\item{} Types have a spectrum of granularity, from the very broad (Parts
of Speech) to the much narrower,
including (at the fine end of the scale) where
they incorporate lexical data
(types can potentially include \i{tiger}, \i{house}, and so on).
In between are constructions related to \q{Ontology},
like sentient/nonsentient, pointwise/extended, artifact/institution, among many others.
\item{} Types are neither strictly grammatic nor strictly semantic, but
their gradations of precision cross between grammar and semantics.
\item{}  Returning to \q{Ontology}: types have associated qualities like sentient/nonsentient;
spatially (and/or temporally) extended, pointwise, or non-spatial
(/non-temporal); caused, self-causing, self-determining, affected by
other things, affecting other things; objects, events, processes, or
institutions; abstracta or spatetime present things; observables
or subjectives like emotions or sensations, which are temporally present
for someone but not (directly) encountered by others.  These are qualities
pertaining to the manner of referents' appearing, causing, and extending
in the world and in consciousness, and to a \q{classification} of kinds
of entities (like a metaphysical Ontology, though the point is not
to reproduce Medieval philosophy but, more modestly, to catalog word senses).
I will refer to these qualities generically as \q{associations}.  They may be
introduced via subtyping or more complex type operators.
\item{}  Some types are \q{function like}: this means that they are
\i{applied} to senses which have their own types.  This introduces one
form of head/dependent relation, where a head word instances a
function-like type and is applied to one or more \q{dependents}.
\item{}  Type information \q{distributes over} Link Grammar pairs.  For
any pair of words which have a meaningful inter-word relation, we can
consider types which may be applicable to both words, and how these
types affect and are affected by the significance of the particular
kind of link.  Some kinds of links mandate particular type
interpretations of the links elements: \TS{} links,\footnote{http://www.link.cs.cmu.edu/link/dict/section-TS.html
} to cite a narrow
example, would only be formed between verb and \Prop{} types (at least
this is a plausible interpretation of the relevant Link Grammar rules.
Other type/link combinations are more
open-ended.
\item{}  Type information similarly \q{distributes} over clusters of
link-pairs, where the presence of one such link influences how a
connected link is understood (or whether it is allowed).  Type-related
qualifications can propagate from one link-pair to connected
link-pairs.\footnote{For example, we can say that the linkage structure in
\q{Three times students asked an interesting question} alters the
normal type-attribution of \q{students} as just a plural noun;
relative to the connected structure linking \q{three times} through
\q{students} to \q{a question}, we can say that \i{three times}
modifies \q{students} so that it may function, as subject of
\q{asked}, as if typed as singular, because \i{three times} acts as a
\q{space builder} and creates a mental frame wherein the students are
singular, even if the word is plural.  Because of this frame
phenomenon, the singular/plural status of students does not propagate
to \q{a question}; collectively they presumably did  not all ask just
one question.  Type annotation for \q{students} has to be defined,
in this case, relative to multiple \q{cognitive frames}.
}
\item{}  Type information also \q{distributes over} applicative
structures.  Given a function-like type we can consider how
associations for the head and dependent elements propagate to
associations on the resulting phrase \mdash{} again, via subtyping or some
other mechanism.
\end{itemize}
Such a \q{linguistic type theory} needs to model (at the least) these
aforementioned associations, the \q{distribution} of type details over
link and applicative structures, and the \q{propagation} of
associations and other type details.  While informal analyses in any
single case may be clear, integrating many case-studies into a unified
theory can be advanced by drawing ideas from rigorous, quasi-mathematical
type theories \mdash{} relevant research has adopted technical
formations like \q{dot-types}, higher-order types, dependent types,
Monoidal Categories, Tensors, Continuations, \q{Linguistic Side Effects},
Monads, Combinatory Logic, and (Mereo)Topology/Geometry.\footnote{Monoids: \cite{DelpeuchPreller};
Tensors: \cite{MaillardClarkGrefenstette};
Continuations: \cite{BarkerShan};
Combinators: \cite{Villadsen};
Side Effects: \cite{ShanThesis};
Monads: \cite{GiorgoloAsudeh}, \cite{ShanMonads}, \cite{Kiselyov};
Topology: \cite{Petitot}, \cite{CasatiVarzi}.
}  Such techniques can marshal type-theoretic ideas without
falling back on simplistic type notions that can end up collapsing a type-system into a
one-dimensional \q{Ontological} classification, rather than exploring more advanced formulations
like higher-order types and (what I am callling) \q{associations}.
}
\p{With respect to Type Theory related to Link Grammar, consider again the \TS{} links
(there are dozens of potential link-grammar pairs, of which \TS{} are among the
less common, but they provide a useful example).  First, note that \Prop{} provides a
type attribution for sentences, but also for sentence parts: \i{he is at school},
for example, presents a complete idea, either as its own sentence or part of a
larger one.  In the latter case, a \Prop{} phrase would typically be preceded with a
word like \i{that}; in the case of Link Grammar, we can define words relative
to their semantic and/or syntactic role, which often lies primarily in linking
with other parts of a sentence or helping those parts link with each other.
Type-theoretically, however, we may want to assign types to every word, even those
which seem auxiliary and lacking much or any semantic content of their own.
Arguably, \i{that} serves to \q{package} an assertion, encapsulating
a proposition as a presumed fact designated as
one idea, for the sake of making further comments, as if \q{making a noun} out
of it: \PropToN{}.  Perhaps our intuitions are more as if \i{that he is at school}
is also a proposition, maybe a subtly different kind, by analogy to how
questions and commands are also potentially \Prop{} variants.  Since \thatPhrases{} are \q{arguments} for verbs,
the choice then becomes whether it is useful to expand our type picture of verbs
so that they may act on propositions as well as nouns,
or rather type \q{encapsulated} propositions as just nouns
(maybe special kinds of nouns).
}
\p{In either case, \i{I know that ...} clearly involves a verb with subject and direct
object: so either \VisNNtoProp{} or \VisNProptoProp{}.  Consider the role of a \TS{}-link here:
specifically, \TS{} connects the verb to the assertorial direct object (most
directly, to \i{that}).  The purely formal consideration is ensuring that
types are consistent: either the \TS{} target is \Prop{}, as I suggested
above, with the verb type modified accordingly; or the \TS{} target is a noun,
though here it is fair to narrow scope.  For this particular kind of
link, the target must express a proposition: either typed directly as
such or typed as, say, a noun \q{packaging} a proposition, which would then
be a higher-order type relation (just as \q{redness} is a noun \q{packaging}
an adjective, or \q{running} is an adjective packaging a verb).  In other words,
it is difficult to state the type restrictions on the link-pair without employing
more complex or higher-order type formations.
}
\p{On the other hand, this is
another example of the fuzzy boundary between syntax and semantics: given a sentence
which seems to link a verb calling for a belief or assertion (like \q{know},
\q{think}, \q{suggest}, \q{to be glad}) to something that is not proposition-like, is such a
configuration ungrammatical, or just hard to understand?  Clearly, the
\i{semantic} norms around verbs like \q{know} is that their \i{subject}
has some quality of sentience (or can be meaningfully attributed
belief-states, even if speakers know not to take it literally: \q{The function
doesn't know that this number will never be zero}); and their \i{object} should
be somehow propositional.  But applying type theory (or type theory in conjunction
with Dependency Grammar) leaves open various analytic preferences: these
requirements can be presented as rigid grammatic rules or as \q{post-parsing}
semantic regulations.  How to model the qualities of sentience (or at least of having
propositional attitudes broadly conceived), for the noun, and of propositionality,
for the direct object, are again at the discretion
of the analysis (subtypes, quality-associations, or etc.) \mdash{} Figure ~\ref{fig:Iknow} shows one potential,
rather simplified unpacking of the sentence; from this structure details
can be added perhaps as extra syntax constraints or perhaps more as cues
to interpretation.\input{figure.tex}  If these
requirements are seen as more syntactic, so qualities are incorporated into
data like Part of Speech (say, a noun designating something with propositional attitudes
being a subtype of of a generic \N{} type), then we are more likely to analyze
violations as simply incorrect (recall \q{The tree wants to run away from the dog}
\mdash{} ungrammatical or just somehow \q{exotic}?).
Some examples suggest less incorrectness as clever or poetic usage \mdash{} so a
richer analysis may recognize expressions as type- and link-wise acceptable, but
showing incongruities (which is not the same as impropriety)
at a more fine-grained type level.  That \i{to want} takes a subject
\i{associated} with sentience does not force type annotations to inscribe this
in grammatic or lexical laws; instead, these associations can be
introduced as potential \q{side effects}, \i{triggering} re-associations
such as forcing hearers to ascribe sentience to something (like a tree) where
such ascription is not instinctive.  The type effect in this case lies more
at the conceptual level, the language-user sifting conceptual
backgrounds to find a configuration proper to the type requirements (in what
sense can a tree \q{want} something?).  In this \q{tree} case we probably
appeal to concepts of \q{as if}: if the tree \i{were} sentient, it
would be nervous of the dog sniffing around \mdash{} a humorous way of calling
attention to the dog's actions (obliquely maybe alluding to people's background
knowledge that dogs sometimes do things, like pee, in inconvenient
places, from humans' perspectives).
}
\p{In brief, it is certainly possible \mdash{} though by no means mandatory \mdash{} to model
type requirements with greater flexibility at a provisional grammatical layer,
and then narrow in on subtypes or extra accumulations of qualifications on
type-instances in a transition from grammar to semantics.  Perhaps cognitive
schema occupy an intermediary role: progressing from basic recognition of
grammaticality, through cognitive schema, to conceptual framing, with type
machinery capturing some of the thought-processes at each \q{step}
(not that such \q{steps} are necessarily in a temporal sequence).  The basic
verb-subject-direct object articulation sets up an underlying cognitive
attitude (represented by a basic type-framing of verb, noun, and proposition,
like the \VisNNtoProp{} signature).  Cognitive ascriptions fill this out
by adding detail to the broader-hewed typing, associating sentience with the
subject and propositionality with the object (sub- or higher-order typing
modeling this stage).  And how the actual lexical choices fit these cognitive
expectations \mdash{} I call them cognitive because they are intrinsically tied
to structurational schema in the type, morphology, and word-order givens
in the encountered language \mdash{} compels conversants to dip into background
beliefs, finding concepts for the signified meanings that hew to the
intermediary cognitive manipulations (finding ways to conceptualize the
subject as sentient, for example).  This also has a potential type model,
perhaps as forcing a type conversion from a lexical element which does
not ordinarily fit the required framing (such as giving inanimate things
some fashion of sentience).  Type theory
can give a window onto unfolding intellection at these multiple stages,
although we need not conclude that the mind subconsciously doing this thinking
mimics a computer that churns through type transformations mechanically and exactly.
}
\p{I envision the unfolding that I have just sketched out as something Phenomenological
\mdash{} it arises from a unified and subjective consciousness, one marked by
embodied personal identity and social situation.  If there are structural stases
that can be found in this temporality of experience, these are not constitutive
of conscious reality but a mesh of rationality that supports it, like the veins in
a leaf.  Stuctural configurations can be lifted from language insofar as it is a
conscious, formally governed activity, and lifted from the ambient situations which
lend language context and meaning intents.  So any analytic emphasis on
structural fixpoints threaded through the lived temporality of consciousness is an
abstraction, but one that is deliberate and necessary if we want to make scientific
or in any other manner disputatable claims about how language and congition works.
In that spirit, then, I will try to condense the three \q{layers} of unfolding
understanding, which as I have sketched them are posited in the metaphysical order
of temporal experience \mdash{} \q{unfolding} in likely overlapping, blending ways
\mdash{} I will \q{read into} them a more static and logically stacked meta-structure.
Where I have sketched three layers or stages of unfolding language understanding,
I will transition to proposing three \q{tiers} of language organization, in particular
three levels where type-theoretic models can be applied.
}

\spsubsectiontwoline{Phenomenology and the Limits of Logic-Compositionality}
\p{The perpetuation of 20th-century philosophical tropes
\mdash{} Analytic vs. Continental Philosophy; humanities vs.
logicao-mathematics \mdash{} perhaps diagnoss something incomplete in
20th century thought that the 21st century needs to transcend, but
it also bears witness to intellectual structures that directly
abut phenomenology itself.  The idea, for example, that
logico-mathematical analysis is distinct from cultural-humanities
theory \mdash{} that works of art, for example, do not have logical
or mathematical forms; that the structures evinced by cultural
and artistic productions (and any system of signs and communication
necessarily has \i{some} rigorous structure) does not itself
have a logical and mathematical dimension, or a structural
form that if specified in isolation from its signifying context
would reveal logical and/or mathematical properties.  A parallel
assumption is that logical and mathematical structures need to be
expressed \mdash{}and indeed \i{located} \mdash{} via a few canonical
theories, such as first-order predicate logic, first-order
quantification, and (perhaps modally inflected) set theory.
So 20th-century Analytic Philosophy unfolded in a
kind of semi-explicit dialog with disciplines such as
linguistics and computer science which promoted, from the
full spectrum of mathematical domains, a handful of
specific paradigms \mdash{} like Zermelo-Frankel Set Theory,
Tarskian quantificational-predicate logic, and
Hindley-Milner type theory \mdash{} that in turn became equally
paradigmatic in fields like formal language analysis and
\q{knowledge enginering}.  This is not to imply that
linguists and computer scientists took their marching orders
from philsophers \mdash{} it's probably equally true that Analytic
Philosphers gravitated to these kinds of mathematical frameworks
because they were proving empircally useful in other fields.
But we can say that a group of disciplines aiming to
technically examine subjects at the intersection between
human thought and formal systems \mdash{} language, knowledge,
valid inference, Artificial Intelligence \mdash{} collectively
progressed toward a vision of what logicomathematical
structure \i{is}, a paradigm (or network of paradigms)
which undergirds any reception of phenomenology insofar
as a \q{Naturalizing Phenomenology} project seeks to
initiate dialog and comparisons between how phenomenology
and Analytic Philosophy, respectively, treat logic,
mathematics, or formal-systematic explanations in general.
}
\p{Interestingly, the logico-mathematical formulations that
arguably have come to dominate methodology on the scholarly
\i{frontier} between science and humanities \mdash{} lingustics,
cognitive science, computer science \mdash{} are different from
those that seem to dominate mathematics and physical sciences
themselves.  Of course, some logical form is intrinsic to
any formal method since formal systems cannot be blatantly
self-contradictory.  But the vehicles by which criteria
of logical consistency are expressed in, say,
linguistics and philosophy \mdash{} presenting sentence-meaning via
first-order logical glosses of their predicate
(locutionary, assertorial) semantics, for instance, or
equating meanings to the conditions in the world
(or possible worlds) that would make linguistic expressions true
\mdash{} these first-order and set-theoretic intuitions do not
necessarily coincide with working mathematicians' visualization of
the spaces where their own axioms and formulae are in force.
Category theory, for example, recognizes that set and predicate logic
are in a sense emergent properties of the different formal domains
where mathematics is enacted \mdash{} sheaves, topoi, manifolds, sites, and
so on \mdash{} which are not isomorphic, propagating to variation
in how logic itself is understood.  Insofar as there is no single
formalization to specify what a \i{predicate} is, or a \i{set},
we are on shakier foundation if we want to apply these notions to
conceptual or linguistic investigations \mdash{} to consolidate the concept
\q{cat}, say, as the set of all \i{felix catus}, or
the meaning of \i{Hugo is sleeping on the sofa} as a particular state
of affairs pertaining wherein a certain tabby cat sleeps on a
certain piece of furniture.  The problem with this
sort of logico-mathematical gloss is that it either employs sets and
predicates in an informal manner \mdash{} which brings us back to
the circularity problems I analyzed earlier, since
an honest assertion that Hugo is sleeping on the sofa obvious
intends to report, modulo the imperfections of informal logic,
that as a point of fact Hugo is sleeping on the sofa \mdash{} or
it tries to elevate the logic from informal to formal,
which legitimately then adds explanatory value to our technical
grasp of the sentence, but then we get tangled in questions like
what sets are.
}
\p{To put it differently, it certainly seems as if the
sentence \q{Hugo is sleeping on the sofa}, as an intended and
sensed artifact in a human world, in some fashion encapsulates and
\i{designates} a proposition, and/or the state of affairs, that this
tabby cat is in repose on that piece of furniture.  The sentence
(deliberately produced sound patterns, perhaps) provides us
with mental access or handle to that proposition, making it an
unitary and isolated object of thought, as well as a thought that
can be shared between people via language.  There is a predicate
there that seems to be \i{intended} in the phenomenological sense
by the sentence as a linguistic gestalt, just as Hugo is intended
in our visual consciousness when we eplicitly see him on the couch,
or at least analogous to that regime of intentionality
insofst as Husserl finds both parallels and contrasts between
\i{apophantic} and preceptual-integrative analysis.
That is, the predicate or fact of Hugo's sleeping at that spot
is the noema to the sentence's phenomena, just as Hugo is the
noema to our seeing-Hugo.  In this sense it is non-circular to
proceed with \q{Hugo is sleeping on the sofa} to mean the
fact of Hugo's sleeping on the sofa, since we are exploiting the
seemingly reundant repetition of the sentence to designate
indirectly a predicate which we have no way of expressing except
through the sentence.  So there is \i{something} there that seems to
have a logical form we need (at least if our goal is theoretical
completeness) to cast light on.
}
\p{But when we then pass from informal to formal logic
we encounter the gap separating mathematics circa 2019 from
mathematics circa 1919.  Hugo, for example, is a cat,
which (in the precincts of informal logic, i.e. a logic
solid enough to shape communication and coordinations of
thought but backgrounded as only one of many structural
parameters delimiting how thoughts and the world avolves)
means that Hugo is a member of the set of all cats.  But
what is a set?  Garfield (a fictional, comic-strip cat)
is also a cat but he does not \q{exist}, so the set of all cats
cannot be enumerated just by (even as a mental exercise) pointing
to them one-by-one.  Moreover, what is referenced by the word
or idea \q{Hugo} in the thought that Hugo belongs to the set of
cats?  We can look at and point to Hugo to get a rough idea; but
the fur Hugo just shed is not a member of the set of cats, even
though it is (or was) a part of Hugo.  It is by no means
obvious that we can, even as another mental exercise, define
\q{Hugo} by an enumeration (metaphysically possible even
if logistically infeasible) of hairs, or cells or molecules, that
via mereological summation grounds the name in a
logically well-established referential framework.
}
\p{On course, it is possible to logically restate a sentence like
\q{Hugo, who is a cat, sleeps on the couch} in first-order logic,
which seems to be arriving at a formal exposition of the predicate
informally intended by the sentence \mdash{} the noema to the
sentence's phenomena.  But formal machinery
only papers over the logical incompleteness of concepts like
\q{set} and \q{predicate}: a logical denotation of the sentence
\i{seems} to refer to the state of affairs that the sentence
tokens (and not just the sentence itself) via quantifier
symbols and other logic artifacts.  I'd argue, however, that these
elements are meta-discursive cues rather than theoretical posits:
they invite us to entertain the thought of the predicate lurking
beyond the sentence via a departure from discursive norms,
the same way that simply repeating the sentence does.  In short,
it may be non-circular to gloss \q{Hugo, who is a cat,
sleeps on the couch} via formations like \ExistsX{} and the
logical form gets us closer to apprehending the abstract predicate,
but this is not because of logical symbols like \q{\Exists{}}
in themselves, but because of their rhetorical effects as
interrupting the flow of conventional philosophical writing.
Symbols of quantification, set membership, or part/whole
inclusion only carry the imprimateur of logical rigor
insofar as we can synchronize what flavor of mathematics
we are using with which to theorize quantification or parthood.
}
\p{We can, for example, argue that Hugo is an extended
three-dimensional thing (or maybe a four-dimensional
perdurant, since he's still here day after day), and
any reasonably compact and continuous \ndim{} manifold can
potentially be designated in thought and language.  This is
reasonable, but it positions our infrastructure for the semantics
of proper names (and a large class of referring expressions)
against a topological backdrop in which we can talk about sets
of points, or perhaps regions spanned by extremum points
(Hugo's fur as the convex hull of Hugo's solid mass) or just
reasons as first-class posits in a point-free algebraic
topology.  We can then present a case for the semantics of
\q{Hugo} as a conventionalized sound adopted by a small group
of humans to designate a particular four-dimensional manifold
\mdash{} at least if we have a sufficient topological or Categorical
foundation to give logical exactitude to concepts like
\q{four-dimensional manifold}.  This still neglects
merological problems, like Hugo losing some fur
but still being himself; either the manifold Hugo
is scattered wherever his fur (and, well, peepee and poopie)
are or the Hugo manifold has fragmented into lots of disjoint
other manifolds, only one of which is \q{Hugo}.  This latter
option suggests that we need to identify the \q{essential parts}
of Hugo, not the inessential parts like those in the litter box;
then the Hugo manifold is the one that contains those parts
\mdash{} i.e. we need a mereological formalization to embed in
the topological one.
}
\p{It is certainly possible that by patching philosophical incompletions
and importing mathematical vocabularies we can piece together a
reasonably compelling theory of mereological/topological
complexes, maybe with graded sets (some parts are more essential than
others), to logically model the semantics of singular referents.
By taking sets of such complexes, maybe
with a modal machinery (for, say, fictional cats) or
graded membership (some things are only partially
subsumed by their concept: is a pet malawolf a dog?) we can
similarly flesh out a set-theoretic semantics, and combining
the singular and set analyses perhaps achieve a logically
inflected semantics applicable to a broad class of linguistic events.
But have we therefore truly surpassed informal logic,
coiled within dialogic conventions, toward a rigorous
theory of linguistic meanings via formal logic?  Do we need
topology, mereology, graded sets, and modal logic because these
are intrinsic to the epistemic domain where signification
happens, or are we just grabbing an assortment of mathematical
tools to plug the leaks in a philosophical attitude wherein
there is a relatively straightforward formalization of
quantifier-predicate logic that can polish the informal
logic of language, but which perhaps failed to
anticipate the mathematical complexities of formalizing what
exactly are the \q{thises} that quantifiers can quantify over?
}
\p{These questions are directly relevant to phenomenology \mdash{}
and to Analytic Philosophy's reception of and dialog with
phenomenology \mdash{} because a phenomenological account of
issues like manifold synthesis and mereological essentialism
can easily bring forward formulations that overlap
with a \q{logistic} paradigm in some details but not others,
and profitably pursuing the dialog depends on identifying both
similarities and differences.  The story of Hugo as a
four-dimensional manifold might not be told in a phenomenological
context, but we might study our visual episodes involving Hugo
in terms of perceptual synthesis: seeing Hugo from the front
we anticipate that we \i{could} see Hugo's tail if we were
standing behind him and it was not as much tucked beneath him;
this anticipation is also a formula for how we integrate
successive impressions as our gaze travels from his ears to
his tail that we see partially.  There is a certain continuity and
topography of our accumulation of successive visuals that
retraces a sort of manifold, even if now we're talking
not about Hugo (or his body) himself but rather the
interplay between the space where he's laying and my own
perspctive and vantage.  So, even
on a phenomenological basis we encounter fragments of analysis
that point toward something like a mathematical treatment \mdash{}
but that doesn't mean we are necessarily committed to a
quantitative-predicate logical theory of preception, or
equally important, of the communicative acts that
perception yields.
}
\p{In pursuit of logical rigor, 20th-century Analytic Philosophy
(plus, say, linguistics, cognitive science, and fields related
to knowledge engineering and information technology, like
AI-driven medical diagnosis) have highlighted
specific logico-mathematical representation-regimes,
building up a paradigm of logicality (as a philosophical
quality, something one senses in analytic methods)
and propositionality (as a corresponding analyzand).
The various layers of this construction address different
aspects of logic, or \q{logicality} as such, such as
predication (that which can be qualified as true or
false in the first place), quantification (establishing
a scope for empirical truthness, so logic does not not decay
into tautology \mdash{} that \i{one} thing or \i{some}
members of set have a property where others don't, modeling
contingency and a posteriority), modality
(capturing the further contingency that most facts
could have been otherwise), mereology and individuation,
quantification scope, supervenience, and other logical
themes that seem intrinsic to any intellectual territory
where predication and true-falseness pertains.
Codification of these facets of logic not only
fine-tunes philosophical analysis; it also provides a
semantics for natural language insofar as we accept that
meaning is centered on locutionary truth-conditions (from
which phenomena like illocutionary force and polarity
\mdash{} our feelings and pragmatic intentions \visavis{}
states of affairs at the predicate substrate of language \mdash{}
can be reached in supplemental analysis) and a philosophical
basis for the formal/computational representation of facts
and knowledge (such as, in database design).  So
insofar as a philosophical-computational-linguistic consensus
emerged on the broad outlines of what logic entails,
this paradigm spread into multiple disciplines,
scientific as well as humanistic, and applied as well
as purely speculative.  Insofar as Naturalizing Phenomenology
involves dialog with Analytic Philosophy, it
intersections with these different faces of logicality and
the mathematical systems which prior generations believed best
captured them.
}
\p{But this does not mean phenomenology needs to endorse the full
stack of paradigms layered in Analytic Philosophy's
excavation of logic; that we need to approach logical-mathematical
eidetics via quantifier-predicate logic, for example, or to
prefer the mathematical frameworks through which
quantification and predication are usually modeled.
Mathematics itself pulls in different directions: quantification,
precication, and set-membership have different
theories in different contexts, e.g. mereotopology compared to
algebraic topology compared to Category theory.  These
lines of development are largely tangential to Analytic
Philosophy partly because the philosophical (and, say,
cognitive-scientific) push toward formalized logic is not really
a push toward \i{mathematics}, but rather a vision of
bringing rigor to what is still essentially \i{informal} logic.
That is, pursuing for instance a formal treatment of
quantification to undergird natural language semantics
does not mean building up to a rigorous, maybe
Categorial theory of quantification in terms
of (say) dependent coproducts, because what \q{logicality}
does in language does not depend on that degree of technical
detail.  But that means in effect that operationalizing
the \q{logicalness} of language is not so much a problem
in \i{formal} (mathematical) logic but a formalized
\i{informal} logic, logico-linguistic or logico-mental
(as in \q{folk theory of other minds}) more than
logico-mathematical.  Fair enough; but this then places other
analytic conveniences (like referents as 4-dimensional
manifolds) on shakier ground, because we're no longer
in an all-encompassing logico-\i{mathematical} paradigm.
Topology and Category Theory, in effect, reveals the difference
between formal-mathematical and formalized-informal logic.
Trying to negotiate with mathematics as source of methodological
rigor, mathematics proves to be an elusive partner.
}
\p{}
\p{}
\p{}

\input{section7.ngml}
\input{section7a.ngml}
\subsection{Concreteness and Computation}
\p{I would be remiss to draw a contrast between 
\i{abstract} and \i{concrete} models without mntioning 
the emergent field of \q{Ontology Engineering} and the role 
of mereology therein.  Arguably, mereology 
in these contexts is neither abstract like logical-axiomatic 
models, nor concrete in the sense of \q{routing} through 
cognitive or scientific models.  Ontologies 
\mdash{} as this term applies in computer science and 
Information Systems \mdash{} are computational artifacts, effectively 
documents in a special-purpose language, enabling 
programmers and scientists to define and 
annotate data structures so as to describe logical 
patterns obtaining within data.  Ontologies allow 
data to be correctly shared and reproduced, and 
also sometimes to be logically analyzed, identifying 
structures within a data set that can be 
extracted via logical rules.  Embedding logical axioms, 
including those of mereology, 
in data modeling frameworks therefoe becomes 
part of an overall data nanagement technology.  Logical 
analysis becomes a part of data analytics in general. 
}
\p{Ontologies are abstract prototypes; they play a determinative 
modeling role only in conjunction with actual data sets.   
Also, although computer programs can perform 
some \q{reasoning} on modeled data sets, the 
explicit correlation between Ontologies and 
empirical data depends on people consciously 
aligning data structures to Ontologies' types and 
relations.  Data sets annotated via formal Ontologies 
therefore reprsent deliberate human attempts to 
relate conceptual and empirical 
tsructures \mdash{} they are both cognitively 
and scientifically concrete.
}
\p{Consider, for sake of demonstration, a derivation of 
the fact that an ingrown toenail should 
qualify as a malady of the foot.  An Ontology 
might represent ingrown toenails afflicting 
toes, and also that toes are part of the foot, 
without identifying an ingrown toenail as a 
malady of the foot per se.  Someone searching 
a database for foot afflictions, however, 
might well expect cases of ingrown toenails to be  
among the search results.  Thus mereological 
rules can extend the semantic reach of an information system.
}
\p{But that semantics depends first on conscious design 
of Ontologies to represent things like 
disease classifications and part/whole relations, and 
also on human c;assification activity.  Someone has 
to specifically notate that a patient's condition 
exemplifies an ingrown toenail.  The Information 
System is designed around human practice, according 
to different operational needs: entering a disagnostic code 
in a medical record, for example, and then 
querying a database to gather a set of results 
such as cases of diagnoses filtereed by some 
condition (say, conditions that might be 
treated by a podiatrist).  An Information 
System \mdash{} its technology as well as its 
actual data \mdash{} can evolve in response to 
operational preferences.  In short, 
deductive reasoning like \i{ingrown 
toenail is a kind of foot malady} is desired not 
out of some abstract concern for logical completeness, 
but it fits operational needs to 
pair a data entry about a patient with an 
ingrown toenail matching a potential database 
query about podiatric diagnoses.
}
\p{Mereology is, accordingly, part of an overall logical 
structure designed to facilitate the operational 
requirements of adding and filtering data 
from information spaces.  This adds a further dimension to 
mereological analysis, because the 
mereology (along with other logical 
relations) becomes part of an engineered system, 
where we have prior anticipation of \i{how} 
we want the system to perform.  We identify specific 
kinds of inferences which should be engineered 
into the system based on anticipated 
preferences of human users.
}
\p{To clarify, we can recognize in an Information 
System at least two different classes of 
\q{behaviors}: \i{input} and \i{output}.  
Input behaviors are the system's implementations, 
and the interactions it affords, for people 
entering new data (say a doctor entering a 
new ingrown-toenail diagnosis); output behaviors are 
implementations and afforances for getting information 
out of the system, such as by queries (perhaps, a 
list of patients with podiatric diagnoses).  
We can add \i{internal} behaviors involving 
internal processing to link inputs and outputs 
in (according to users' expectation) correct 
ways.  The system-outputs are products of 
the cumulative history of its inputs up to the 
point where output behavior is requisite.  
For example, the prior input of an ingrown-toenail 
diagnosis constitutes \i{relevant} input history 
for a query (demanding output 
behavior) for podiatric diagnoses.
}
\p{So mereology and other logical relations play a role when it 
comes to identifying \i{input history} that is \i{relevant} 
for a certain output rquest.  Logical reasoning 
expands the space of input history which will be 
considered relevant in a given output-context.  
This points to an interpretation of mereology in 
terms of \i{relevance}; properties or 
assertions involving a part (e.g., diagnoses) 
are relevant to corresponding whole.  
}
\p{Notionally, then, mereology serves as a kind of 
semantic enrichment of a system whose essential 
nature is a mapping from \i{input history} 
to \i{output behavior}.  This core 
operational motive \mdash{} \i{input history 
determines output behavior} \mdash{} should be reiterated 
in the context of the Semantic Web in particular, 
to help establish an underlying conceptual 
model to properly analyze (and implement) 
Semantic Web technology.
}
\p{Conceptually, the Semantic Web is one 
example of a data-sharing platform; 
and the key ingredient of data sharing is 
the replication of information spacs, 
in whole or in part, between different 
places.  Insofar as \q{input history 
determines output behavior} is the 
constituent feature of information spaces, 
replication means that the points 
where the data is shared will replicate 
analogous input-to-output patterns.  
Data sharing does not necessarily mean that 
all points process data identically; when data is 
exported between environments the new location 
may serve different operational purposes, 
so that analogous queries or interactions 
which receieve output data may be evaluated differently  
at different points in a data-sharing network.  
But assuming the \i{receiver} of shared data 
does differ in its output structures from 
the \i{senders}, these differences should be 
systematically accounted for \visavis{} how 
the respective systems are engineered. 
}
\p{Conventional Semantic Web data structures are 
organized around the principle of \i{labeled graphs}, 
where the building blocks of data models are 
so-called \q{triples}: \i{Subject/Predicate/Object} 
graph units where \i{Subject} and \i{Object} 
are graph-nodes and \i{Predicates} are labels 
applied to graph-edges, asserting that the \i{Object} 
bears to the \i{Subject} some 
identified relationship.  Ontologies then, 
first and foremost, provide lists 
of identifier for both nodes and edge-labels: 
this allows nodes to be associated with a kind or 
type (a \i{person}, say, or a \i{diagnostic code}) 
while edges are labelled according to a controlled 
vocabulary (so edges represent concrete 
relations, for example, that a person has a diagnosis).  
Ontologies then superimopose logical axioms 
pertaining to restrictions or guarantees among relations and/or 
nodes: that a relation is transitive, 
say, or that a given relation is only reasonable 
for certain \i{Subject} and \i{Object} types 
(e.g., a diagnostic relation is only 
coherent as a relation between a person or 
patient and a diagnosis or diagnostic code).
}
\p{Although a fair range of empirical semantics can be 
captured by individual \i{Subject/Predicate/Object} 
triples, in many cases a realistic semantics 
demands a model of how multiple relations aggregate 
to form meaningful units.  For example, instead of 
a diagnosis being modeled as a single relation between 
a patient and a diagnostic code, a systm may 
stipulate a more detailed representation where a single 
\q{diagnosis} has several associated piecs 
of information, such as a diagnostic code, a date, a 
doctor's name, and potential 
reference to laboratory or radiological findings that 
substantiate the diagnosis.  This relflects how, in 
most computational settings, the basic units of 
evaluationb are not single data points but internally 
structured data aggrgates.
}
\p{Thus far, no standard model has emerged to represent 
larger-scale data structures in the Semantic Web 
context.  Intrinsically, the \q{semantics} of the Semantic Web 
depends on computer environments where data is shard and 
received.  Since people generally do not consider data 
sets as a whole, any ,eaningfulness we can attribute 
to dignital information depends on selecting 
smaller amounts of data, which in turn depends on 
the digital systems \mdash{} user interface technology as 
well as data management software \mdash{} through which 
people interact with information 
senses.  So any discussion of 
the \q{semantics} of data, formally warehoused in 
information spaces, 
has to recognize the multiscalar organization 
intrinsic to information spaces. 
}
\subsubsection{Mereology and Hypergraphs}
\p{The interaction between humans and information spaces 
is itself a technological artifact which has 
to be engineered and implemented.  Such implementation 
in turns depends on implemented \i{procdures}, 
which classify, manipulate, and visually or 
textually present relevat data to human users.  
A major element in the semantics of information spaces 
is accordingly the semantics of procedures operating on this 
data \mdash{} what data structures are valid operands 
for procedures' operations; how procedures mody 
data structures or map inputs to outputs; how data 
manipulated procedurally must be internally structured, 
according to procedures' computational assumptions 
and preconditions.  The constituent elements of 
information spaces are internally structured data aggregates, 
more so than they are unstructured information atoms, 
because structured data is the kind of value which 
in the general case computational oprocedures operate upon. 
}
\p{For these reasons, part/whole aggregation and the rules for 
nesting data structures (directly or indirectly) 
should be essential topics for the Semantic Web (separate 
and apart from mereological axioms built into 
web Ontologies).  Thre are at last four different 
strategies that can be used to represent multiscale data in a 
graph or labeled graph context.  These are not mutually 
exclusive but have different emphases; insofar as 
labeled graphs form the notional foundation of 
the Semantic Web they can carry over into 
that context as well.  Graphs supporting multi-scale 
representations are generically called \q{hypergraphs}, 
although there are different hypergraph  
theories providing different spaces of models. For sake 
of illustration, I'll use the hypothetical case 
of a diagnosis encompassing a diagnostic code, 
date, and the diagnosing doctor's name, 
and refer to schematic diagrams in Fig1a-e.  
}
\p{One form of hypergraph is based on the 
idea of \i{hyperdges}, and further 
\i{directed} and \i{labeled hyperdges}, which 
generalize graph-edges to include more than two nodes.  
For instance, a record of a diagnosis even 
can involve a four-node edge where one 
node represents the patient, and the remaining 
three nodres represent a diagnositic code, date, 
and doctor's name.  Directed Hypergraphs 
group the nodes incident to a hyperedge 
into \i{source} and \i{target} nodes; in 
this example the patient's node is the 
lone source and the remaining nodes are targets, 
indicating several pieces of information associated 
with a diagnosis (see Fig1a).  The (potentially multi-node) 
sourc and target sts generalize the Semantic 
Web concpt of \i{Subject} and \i{Object} 
for context where a given predicate-relation 
needs to be expressd with multiple subjcts and/or 
objects.
}
\p{A second model of hypergraphs can be defined by allowing each 
node to be (or internally contain) its own graph.  
This leads to nestd levels of graphs, where 
a graph on one level can be incorporated into a graph on the 
next higher level by being nested into one of 
the latter graph's nodes.  Figure 1b shown how 
this could work in the running xample: it 
shows the fact of a doctor giving a diagnosis 
on a given date as a miniature graph whose edges 
connect the doctor's name to a \q{date} node and 
\q{diagnostic code} nod, respctively.  
The mini graph is then placed as a package 
via a node inside the higher-level graph, where 
that node is joined to a nod designating the 
diagnosed patient.  Models may or may not allow 
nodes on one level to be linked with nodes on another 
level, or allow two nodes nested inside two different higher-level 
nodes to be linked to each other.
}
\p{A third genre of hypergraph has similar aggregational features 
but does not distinguish between different levels.  Instead, each 
graph has one level, but a collection of nodes and/or 
edges can be separately identified and associated 
with a separate node, or presented as some aggregate 
structure with special significance.  Hypergraphs 
with this added structure function like conventional 
Semantic Web graphs upon which are superimposed 
added details.  Figure 1c shows the extra layer of 
information via a dotted line around three graph edges.  
The underlying data in this case is a regular 
(labeled, directed) graph (not a hypergraph).  
A separate node is inrroduced representing a 
\i{diagnostic event}, serving 
essentially as a \q{placeholder} to which 
estra nodes (carrying factual details) 
can be attached.  This is a common pattern 
sometimes called a \q{snowflake}: a \i{Subject} 
node (here the \i{patient}) is 
connected to an \i{Object} 
node (here a \i{diagnostic event} whose 
semantic role is to unify 
additonal nodes, that provide the relevant 
actual details; the net effect is nodes branching 
off from a center, like a snowflake.  The extra 
layer of information comes into play by unifying 
the three nodes and edges one step removed 
from the \i{patient}.  These edges would be grouped 
together to indicate that their respective objeect-nodes 
form a logical unit (three components of a 
diagnosis: medical code, date, and doctor's name).  
They could also indicate data requirements: an  
Ontology can stipulate that if a \i{diagnostic event} 
node is present in a graph, it must be accompanied 
by three additional edges providing 
pertinent info about that event.
}
\p{Consider also the possible interrelationships 
between the first and third cases (hyperedges 
and frames).  A consequence of hyperedges 
in \i{directed} hypergraphs is that each 
hyperedge can have distinct head- and tail-sets, 
which are called \q{hypernodes}.  So a directed 
hyperedge connects two hypernodes.  In a typical 
case, a hypernode packages together multiple 
interrelated values, like \i{name}, \i{adress}, 
and \i{telephone number} or \i{date}, 
\i{doctor's name}, and \i{diagnostic code}.  
Meanwhile, in a frame model, several \i{edges} 
may be grouped togther to provide an interrelated 
package of information.  Combining these two ideas, 
we can define a sort of \i{cross-reference graph} 
whose nodes are hypernodes and whose edges are 
defined by values within a hypernode that \q{point} to 
other hypernodes.  A hypothetic case is 
outlined in Fig1d, with the assumption that 
elements in a hypernode are proxies for 
additional complex data structures.  For 
example, diagnostic codes can refer to 
a collection of known medical conditions each 
linked to additional data, such as symptoms 
and treatments.  Similarly, a doctor's name 
can link to a profile of that doctor 
in a database, and a \i{date} represented as a single 
value can be treated as a stand-in for an 
aggregate of \i{day}, \i{month}, and 
\i{year}.  A set of hypernodes whose component 
values cross-reference to other hypernodeas can 
be seen as defining a hypergraph whose edges are 
cross-references themselves; a computer 
representation of such a data space may not even need to 
recognize a separate model of \q{edges} apart 
from cross-reference values. 
}
\p{A final option for representing hypergraphs 
comes into effect if we recognize a \i{type system} 
applied to a graph; a classification of nodes 
according to an interrelated collection 
of types.  One feature of most type
systems is that for any type \ty{} there are  
corresponding types representing \i{lists} of 
\ty{}-values with different characteristics.  Lists are 
often classified in terms of how they are 
modified.  \i{Stacks} are lists where the last value added 
(the one which has been on the list for the 
shortest time) is the is the first value removed 
(think of a stack of trays in a cafeteria).  
\i{Queues} are lists where the value in the 
list for the \i{longest} time comes off the 
list first (think of people 
lining up for tickets).  And \i{Deques} 
(double-ended queues) can have values added 
or removed from both ends.  (There are more 
complex aggregate data structures as well, like 
associative arrays where values of one type 
are used as indicies to retrieve values of 
a second type from a long list; consider how a phone 
book maps people's names to telephone numbers).  
Collectively, these multi-valued types 
are often called \q{container} types.
}
\p{One way to construct hypergraphs is to allow 
node-types to be container types as well as 
singled-value types.  Therefore, a node may 
encapsulate a many-valued data structure.  
This is similar to the case depictird 
in Fig1b, where nodes could contain nested 
graphs; the difference here is that the nested 
data structures encapsulated in nodes are not necessarily 
graphs; they can also be different sorts of lists or tables.  
In Fig1e, this is shown by 
imagining a table that represents a patient's diagnostic 
history.  Each new diagnosis is added to the list, 
so the node can expand over time.  The data structure 
\q{inside} the node then functions as a table, 
whose rows represent diagnoses made on a single 
occasion and whose columns are different facets 
of a diagnosis (date, code, and doctor's name).  
The entire table is treated as one node, linked 
to the \i{patient} node 
(indicating that the provided table is diagnostic 
history for that one patient).
}
\p{These five different kinds of hypergraphs are to some degree 
interchangable, in that algorithms can transform 
hypergraphs between these formats.   They are 
still more or less optimal for different contexts.  
The fifth solution is similar to how container types 
work in conventional programming 
languages \mdash{} and also to foreign 
key links in relational databases \mdash{} so hypergaphs 
using that system can be easier to 
integrate with conventional (not necessarily 
Semantic Web) applications.  On the 
other hand, the third solution is closer 
in spirit to the Semantic Web, adding an extra 
level of specification but preserving 
the underlying Semantic Web graph thory.  Meanwhile, 
the nested-graph model is similar to 
how popular enterprise graph databases are 
designed; and the directed-hyperedge approach 
is closest to mathematical graph theory.  
}
\p{The diversity of related by not identical approaches points 
to an incomplete dimension to the Semantic Web; 
multiscale representations have 
not been standardized.\footnote{There is an official 
\q{RDF Collections} model but this addresses only a small 
part of genuine multscale modeling.}  Not that a 
single form of multiscale structure should be preferred 
over others; there are competing priorities which come 
to the fore in different contexts.  
This actually reflects competing ideas about what the 
Semantic Web should try, first and foremost, to accomplish.  
The Semantic Web can be seen as a potentially 
global, open-ended data-sharing community; as a foundation 
for a future of AI-driven networking; as a kind 
of vast, decentralized graph database; or as an 
analytics ecosystem that allows logical analytic 
methods to be applied, alongside othr 
(e.g., statistical) forms of data analysis, 
to travserse and filter information from 
the internet.  Which priorities we emphasize 
dictate the kinds of tchnologies that have 
to be implemented, and the \q{semantics} of the 
Semantic Web can reflect the need to optimize 
representations and network interactions to 
benefit those priorotized technologies.  
}
\p{What this diversity of option reveals, however, is that 
the question of multiscale data represenatation 
is complex and multifacted.  Instead of just 
setting forth one system as a speculative enterprise, 
engineers have to identify different forms 
of multiscalar data \mdash{} for eample, 
I enumerated five versions of hypergaphs \mdash{} and study 
which flavor is best suited for different technological 
priorities and acceptable trade-offs.  
Implicitly, I believe that such a field of study represents 
a kind of computational investigation of 
mereology \mdash{} not mereology tied to 
predefined logical axioms, but an organizing 
pattern for computational phenomena 
(e.g., hypergraphs) whose properties 
have to be codified and optimized.
}
\p{This brings me to the final characterization 
of mereology I'll propose here: Hypergraph Mereology.  
I think that analyzing hypergraphs in the computational 
and software engineering contexts presents a philosophically 
interesting picture of formalized mereology 
no less relevant \mdash{} and perhaps more cognitively 
or scientifically appropriate \mdash{} then a \q{logical} or 
\q{axiomatic} mereology. 
}
\p{Note that Hypergraph Mereology is not necessarily 
set forth axiomatically; we do not specifically 
define hypergraphs by stating their axiomatic 
conditions.  We \i{could} proceeed in this 
way, of course; hypergraphs are a meaningful 
mathematical subject.  But often the formalization 
of hypergraph models is indirect, 
relying on software implementations rather than 
logical specification.  In other words, much of the 
intellctual substance behind research into 
hypergraphs is built from practical implementations 
of hypergraph code libraries, allowing us 
to study and compare how these libraries are
abstractly conceived and concretely implemented.
}
\p{This demonstrates a phenomeneon which I think Analytic Philosophy 
has not fully registered: the kind of 
formalizing efforts that in decades past may have been 
pursued via logical or logicomathematical axiomatics 
can increasingly be replaced by software and compurter 
code.  Testing philosophical intuitions in an 
intellectual environment oriented to concrete code 
devleopment rather than logicomathematical 
abstraction has several implications: there is a 
kind of feedback and trial-and error which 
can shape the intellectual process; the ecosystem
of open-source code and software development tools 
has diffferent institutional norms 
than communal deliberation over the 
validity of logical argumentation; and 
code development tangibly augments academic 
writing, in that philosophers (or writers from 
any other discipline) can pair their conventional 
publications with code repositories and data sets.  
This kind of research can lead 
philosophy in an empirical, experimental 
direction, but also simultaneously in a 
direction which accords to the individual philosopher, 
I think, a renewed agency and autonomy.  
In this spirit we are not only sharing 
with our peers rational claims, but 
exhibiting practical technology created 
under the inspiration of philosophical arguments.  
Well-designed, open-source code projects have a 
certain inner empowerment for their contributors; 
in a world of (often imperfctly 
apportioned) intellectual capital, they 
act as fiat currency, neither curated nor 
endorced by academicinstitutions 
in the conventional sense, and as such 
not diluted by the not-always-rational 
timelines of academic prestige.
}
\p{It should be uncontroversial to say that \q{Analytic} 
philosophy is the prestige dialect of the mdern philosophical 
community, especially in English.  Partly this is no 
doubt topical \mdash{} Analytic philosopers seem to entertain 
more \q{hard science} themes than the political 
or cultural concerns, or issues of subjective 
experience and personal identity, guiding 
\q{Continental} philosophy.  But partly it 
also reflects the presupposition of analytic 
rigor accompanying analyses that can be 
grounded in, or at least shown to be 
inspired by, some kind of logial 
or mathematical framework. 
}
\p{With that said, some of the divergence between 
a logically-inspired Analytic paradigm and a more sociocultural 
Continental thought is methodological, not just 
thematic.  Logical models which discount the 
nuance and complexity of language and experience 
risk becoming reductionistic:  
limited explanatory value.  
Potentially, then, the emergence of a computational discipline 
for addressing philosophical 
structures can change the familiar 
Analytic/Continental narrative.  Software models 
can be technically drawn without being either 
axiomatic nor reductionistic; they can 
engage a formalization in software implementation in 
lieu of specifications of logical structure;  
and they can present technological 
systems as suggestive approximations of 
human phenomena without committing 
to a reductive theory that cognitive 
processes \mdash{} in situational understanding, 
linguistic meaning, or conceptualization, 
for example \mdash{} are wholly characterized by the 
terms of the implemented software.\footnote{Type-theoretic analyses of natural-language semantics and syntax, 
for example, can represent linguistic artifacts 
according to formal/computational structures 
without implying that the human 
procss of actually converting surface-level 
language to these derived representation 
happens mechanically or  mathematically.  
Computational \i{models} of linguistic phenomena 
are different from Natural Language 
\i{Processing}, which makes the further 
metaphysical assumption that 
computers can \q{parse} language down to 
schematized representations analogously to people.
}
}
\p{It is possible, then, to treat computational 
models as methodologically useful adjuncts to 
philosophical analysis while rejecting 
reductive analysis of cognition or 
language as transparently itself computational 
or logically tractable.  In this 
way, computational models can bring a technical 
methodology to Continental thought 
analogous to the role of logical 
axiomatics on the Analytic side, but 
with different methodological 
commitments \mdash{} and also in a different 
metatheoretical perspctive.  Whereas Analytic 
philosphrs may look on logicomathematical 
techniqus as a kind of rational kernel, 
a core ground of secure entailments on which 
to rest other claims, the Continental 
engagement with computational models can be 
more figured within the general 
plurality of Theory; the same pull toward hybridization 
that would wed phenomenology to analyses of 
race, class, and gender, for example, to elucidate 
the experiential as well as cultural political 
dimensions of social relations.  
Humanities \q{theory} acts as a kind 
of intellectual gravity, pulling in phenomenology, 
Critical Theory, Structuralism, linguistic and discourse 
analysis, and so forth; there is no \i{a priori} 
obstacle that would provent the same force 
from reaching, say, formal and natural-language type 
theory and the analysis of multiscale information systems.
}
\p{To my mind the most well-developed body 
of working merging phenomenology with 
computational methods is that of Jean Petitot; 
although I would add some at least 
provisional formal elaborations 
of Cognitive Grammar (if I may take this 
as a kind of applied phenomenology in the linguistic domain) 
by, for insta,ce Kenneth Holmqvist and Matt Selway.  
I might add, implicitly, the larger tradition 
of type-thoretic semantics represented by the sometimes 
highly formal work of a Zhaohui Luo, 
James Pustehovsky, or Bruno Mery; which 
in turn establishd technical links to 
linguistic models inspired directly by 
computer programming languages and computing 
paradigms, as represented by linguists such as 
Ash Asudeh, Gianluca Giorgolo, Oleg Kiselyov, and 
Chung-Chieh Shan.  In short, we can develop an 
overarching model of linguistic representation 
which is anchored both in cognitive 
phenomenology and in computational 
procedures. 
}
\p{Alongside this linguistic line of development we 
can identify a parallel integration between 
phenomenology and theories of spatial comprehension, 
perceptual synthesis, and cognitive 
\q{visualistics}.  One line of analysis, for example, 
considers the phenomenon of experiential \q{presence}, 
motivating phenomenological treatment of 
computer imagery and Virtual Reality.  Along 
this subjective assessment we can consider the 
mathematical infrastructure for creating realistic 
depictions of three-dimensional scens and surfaces; 
and for applying lighting, textures, and space-filling 
effects (mimicking the impression of smoke or flowing 
water, say).  This technology does not 
necessarily recreate the neurological processes by which 
we cognize our spatial environment \mdash{} nor 
create perfectly lifelike graphics \mdash{} but 
it presents a suite of technological formalisms 
that bear a simulative and structuraly articulated 
relation to visual phenomenology.  As such it provides a 
ground for complimentary but non-reductive analyses 
that can be set alongside, for example, theoriesof 
spatial synthesis as found in, say, 
Husserl or Merleau-Ponty.
}
\p{These linguistic and spatial/visual examples have a 
common architecture, at least \visavis{} phenomenology, 
in that they present computarional structures that 
can be argued to encapsulate cognitive representations 
without eliminating the complexity, situationality, 
phenomenological immediacy, and interpersonal 
dimensions of consciousness.  Computational models 
are suggestive analogies rather than axiomatic 
reconstructions of cognitive \q{logic}.
}
\p{I think mereological theories can have an analogous 
intellectual architecture, albeit in a narrower 
topic area.  Formal systems developed in the 
context of computer code \mdash{} say, hypergraphs in 
one of the five models outlined above \mdash{} provide a 
tangible but technological medium for 
articulating mereological concepts in a 
rigorous domain, while also allowing these models 
to be adjunct prototypes for theories of mereology 
as a cognitive phenomenon.  Hypergraphs, in 
short, are to mereology what formal parse-representations and 
lexicons are to natural language, or computer-generated 
imagery is to visual experience.  
}
\p{}

\input{section9.ngml}
\input{section10.ngml}
\input{section10a.ngml}
\section{Conclusion}
\vspace{-.5em}
\p{There are several tactics to practically employ techniques I have 
reviewed here in concrete projects.  For sake of discussion, I will focus 
on the approach used by the demo code-set provided with this chapter.  
This code actually has multiple dimensions, because it has been designed 
to integrate with data sets accompanying other chapters in the 
present volume, not merely this chapter.  This has led to the code being 
planned as follows: code published with \i{this} chapter directly uses 
its Hypergraph and Channel Complex libraries to parse Interface Definition 
files associated with \i{other} chapters.  Those chapters' data sets in 
turn employ \Cpp{} code whose 
construction is influenced by Channel Algebra 
\mdash{} notably in the design of function groups which are overloaded 
via types that are reasonable for \q{hidden channels}, 
and in how functions are exposed to external scripting and reflection.  
The implementations are in normal \Cpp{}, 
but the associated \IDL{} annotations 
describe them via more sophisticated channel formations, which are 
mimicked (in terms of conventional input parameters) in 
the production code.
}
\p{The \q{hidden} values constructed while routing between logically 
related functions, as seen in the various code-sets, tend to 
fit into two classifications.  On the one hand, these may be values 
constructed as a test that function arguments conform to protocol: 
for instance, one way to check that a numeric value falls within some 
desired range is to cast this value to a range-checked type 
\mdash{} whether or not the second value is actually used.  
Alternatively, in some places supplemental values are 
enumerations of typestates \mdash{} allowing function bodies 
to achieve something resembling pattern-matching as a 
code-branching structure (by \switch{}ing on enum values).
}
\p{Manually creating residual values along these lines is 
admittedly cruder than pure type-level engineering, as exemplified 
by Idris's synthesis of Dependent Types, typestate, and 
effect-typing.  The solutions chosen for the present context are 
less automated; they require deliberate choice by developers.  
Nevertheless, the Interface Definition technology makes these choices 
explicit and documented, which can help guide developers toward 
embracing these more fine-grained coding conventions and 
confirming that they have done so consistently.
}
\p{It is a legitimate question whether establishing coding paradigms 
\mdash{} in a data set management context, for example \mdash{} in this 
convention-driven spirit is better or more maintainable than 
approaching data sets in a language like Haskell or Idris whose 
rigors are more formal than conventional.  But the 
\Cpp{}-based approach has the benefit of more immediate integration 
with \GUI{} code, reflected in the \Qt{} components published 
as wrappers and visualizers for data-set-specific datatypes.  
}
\p{In the best-case scenario, \GUI{} code is a natural extension 
of programming languages' and individual languages' 
type systems \mdash{} with rigorous mapping between value types and 
visual-object types that graphically display associated 
values.  So long as Functional Programming languages remain 
operationally separated from \GUI{} drivers \mdash{} needing a 
\q{smoke binding} or foreign-function 
interface to marshal data between a 
\CLang{}-oriented User Interface and the non-visual data 
(and database) management components of a project \mdash{} the strongly-typed 
rigor of Functional Programming environments will be 
incomplete.  
}
\p{Whether or not by technical necessity or just entrenched culture, 
\GUI{} frameworks are still predominantly built in procedural 
or \OO{} languages like \CLang{}, \Java{}, and \Cpp{}.  It seems likely 
that a truly integrated type system, covering User Interface as 
well as data management logic, will need a hybrid functional/procedural 
paradigm at some stratum \mdash{} either the underlying \GUI{} framework or 
at application-level code.  So long as \GUI{} frameworks remain 
committedly procedural, the most likely site for such hybrid 
paradigms to emerge is at the application level, and in the 
context of application-development \SLE{} tools.
}
\p{Implementing \GUI{} layers in a Functional environment is usually 
approached from the perspective of \i{functional reactive} 
programming, which emphasizes how the interface between visual 
components and controller logic can be structured in terms 
of \i{event-driven} programming.  In this paradigm, there is 
no \i{immediate} linkage between \GUI{} events and the 
functions called in response \mdash{} for example, no 
single function that automatically gets called when the user 
clicks a mouse button.  Instead, events are entered into a pool 
wherein each event may have a varying number of handlers 
(including being ignored entirely).  This style of 
programming accords well with paradigms that try to minimize 
the number of functions with side effects.  Event-handlers are free 
to post new signals (these are interpreted as events) which 
may in turn be handled by other functions \mdash{} so that 
signals may be routed between multiple functions entirely 
without side-effects.  That said, most events \i{should} cause 
side effects eventually \mdash{} for instance, after all, 
a user does not typically 
initiate an action (triggering an event) without intending 
to change something in the application data or display.  But 
events can be routed between pure functions until an eventful 
handler is called, so side-effects can be localized in a 
proportionately small group of functions.
}
\p{Moreover, certain qualities of \GUI{} design can be expressed as 
logical constraints rather than as application-level state enforced 
by procedural code.  For example, optimal design may stipulate that some 
graphical component must automatically be resized and repositioned to remain 
centered in its parent window, sustaining that geometry even when 
the window itself is resized.  
Functional-Reactive frameworks allow many of these constraints to be 
declared as logical axioms on the overall visual layout and properties 
of an application, constructing the procedures to maintain this 
state behind-the-scenes \mdash{} which minimizes the extent of 
procedural code needing to be explicitly maintained by application developers.
}
\p{But while Functional Reactive Programming is a strategy for 
providing \GUI{} layers on a Functional code base, it can equally 
be treated as an Event-Driven enhancement to Object-Oriented 
programming.  In an \OO{} context, events and signals constitute an 
alternative form of non-deterministic method-call, where signal-emitting 
objects send messages to receiver functions \mdash{} except indirectly, 
passing through event-pools.  Indeed, as documented 
by the demo code, events and signals 
have a natural expression in terms of Channel Algebra, where both signal 
emitters and receivers are represented via special \q{Sigma} channels.  
On this evidence, Functional Reactive Programming should be assessed not 
just as a \GUI{} strategy for Functional languages but as a 
hybrid methodology where Functional and Object-Oriented 
methodologies can be fused, and integrated.
}


%\vspace{.4em}

%\onecolumngrid
\onecolumn
\begin{thebibliography}{99}
	
\vspace{-.7em}

%
\bibitem{RaubalAdams}
Benjamin Adams and Martin Raubal, 
\bibtitle{A Metric Conceptual Space Algebra}.
\tinyurl{https://pdfs.semanticscholar.org/521a/cbab9658df27acd9f40bba2b9445f75d681c.pdf}

\bibitem{RaubalAdamsCSML}
Benjamin Adams and Martin Raubal, 
\bibtitle{Conceptual Space Markup Language (CSML): Towards the Cognitive Semantic Web}.
\tinyurl{http://idwebhost-202-147.ethz.ch/Publications/RefConferences/ICSC_2009_AdamsRaubal_Camera-FINAL.pdf}

\bibitem{InteractingConceptualSpaces}
Bob Coecke, \i{et. al.}, 
\bibtitle{Interacting Conceptual Spaces I:
Grammatical Composition of Concepts}.
\tinyurl{https://arxiv.org/pdf/1703.08314.pdf}

\bibitem{BrendanFong}
Brendan Fong, \bibtitle{Decorated Cospans}
\tinyurl{https://arxiv.org/abs/1502.00872}

\bibitem{BrendanFongThesis}
Brendan Fong, \bibtitle{The Algebra of Open and
Interconnected Systems}
\tinyurl{https://arxiv.org/pdf/1609.05382.pdf}


\bibitem{GardenforsZenker}
Peter \Gardenfors{} and Frank Zenker,  
\bibtitle{Theory Change as Dimensional Change: Conceptual Spaces 
	Applied to the Dynamics of Empirical Theories}.
\intitle{Synthese 190(6)}, pp. 1039-1058, 2013.  
\tinyurl{http://lup.lub.lu.se/record/1775234}


\bibitem{PetitotSyntaxe}
Jean Petitot, 
\bibtitle{Syntaxe Topologique et Grammaire Cognitive}
\tinyurl{https://www.persee.fr/doc/lgge_0458-726x_1991_num_25_103_1610}

\bibitem{SmithGibbons}
Michael Anthony Smith and Jeremy Gibbons,
\bibtitle{Unifying Theories of Objects}
\tinyurl{http://www.cs.ox.ac.uk/jeremy.gibbons/publications/uto.pdf}


\bibitem{SpivakKent}
David Spivak and Robert Kent, 
\bibtitle{Ologs: A Categorical Framework for Knowledge
Representation}
\tinyurl{https://journals.plos.org/plosone/article/file?id=10.1371/journal.pone.0024274&type=printable}

\bibitem{Strle}
Gregor Strle, \bibtitle{Semantics Within: The Representation of Meaning Through Conceptual Spaces}.  
Univ. of Novi Gorici, dissertation, 2012.

\bibitem{BaltasarTrancon}
Baltasar Tranc'on y Widemann, \i{et. al.}, 
\bibtitle{Automized Generation of Typed Syntax Trees via XML}
\tinyurl{http://pizza.cs.ucl.ac.uk/xse01/ready/10.pdf}


\bibitem{OrlinVakarelov}
Orlin Vakarelov, 
\bibtitle{Pre-cognitive Semantic Information}
\tinyurl{http://www.u.arizona.edu/~okv/Papers/Pre-CogSemInfo_KTB.pdf}

\bibitem{MartinWesthead}
Martin Westhead, \bibtitle{BinX -- The Binary XML Description Language}
\tinyurl{http://buphy.bu.edu/~brower/SciDAC/doc/BinXIntro2.pdf}

\end{thebibliography}

\end{document}
